\documentclass[output=paper,colorlinks,citecolor=brown]{langscibook}
\ChapterDOI{10.5281/zenodo.22078575}
\author{Hana Filip\orcid{}\affiliation{Heinrich Heine University Düsseldorf}}
% REPLACE THE ABOVE WITH YOU AND YOUR CO-AUTHORS, YOUR AFFILIATIONS, AND YOUR ORCID NUMBERS
% rules for affiliation: If there's an official English version, use that (find out on the official website of the university); if not, use the original
% orcid doesn't appear printed; it's metainformation used for later indexing

\SetupAffiliations{mark style=none}
% comment in (i.e. remove the % sign at the beginning of the line) the above line if you are a single author or all authors have the same affiliation


\title{Genericity and habituality in Czech}
% REPLACE THE ABOVE WITH YOUR PAPER TITLE
% provide a shorter version of your title in case it doesn't fit a single line in the running head
% in this form: \title[short title]{full title}

\abstract{In Slavic languages, the domain of genericity (and habituality) is traditionally\linebreak viewed as associated with contextually determined meanings of imperfective forms, and to a lesser extent of perfective forms.
What has so far remained largely unexplored, and misunderstood, is the verb marker \textit{-va-}, which is dedicated to deriving verbs  that only have a generic (and habitual) interpretation. 
%In West Slavic languages (Czech, Slovak and Polish) it is used productively in all registers, giving rise to a regular and transparent derivational pattern. In East and South Slavic languages it is also attested, but it is reported to be mostly unproductive, at least in their standard varieties. 
Based mainly on Czech data, 
it will be shown that the verb marker \textit{-va-}
has distributional and semantic properties that 
preempt its analysis as a marker of tense and/or aspect (\textit{pace} \citealt{dahl1995}, i.a.). Neither can its meaning be reduced to standard (modal) operators that are used in the analysis of progressive or imperfective aspect (\textit{pace} suggestions  in \citealt{cable2022}, i.a.). Its
quantificational and modal (intensional) components strongly suggest that it is a type of generic marker delimiting a subclass of generic sentence. The insights reached here support
the view that generic markers in natural languages not only exist, but are not as rare as is generally thought, and by the same token also 
bolster independent arguments made elsewhere for the independence of genericity (and habituality) from the categories of the TMA (tense, modality, aspect) system. 

\keywords{genericity, habituality, grammatical aspect, imperfectivity, perfectivity,  progressivity, Slavic languages, Czech}
}

\IfFileExists{../localcommands.tex}{
   \addbibresource{../localbibliography.bib}
   \usepackage{tabularx,multicol}
\usepackage{url}
\urlstyle{same}

 
\usepackage{langsci-optional}
\usepackage{langsci-lgr}
\usepackage{langsci-gb4e}

% ADDED BY VOLUME EDITORS

% \usepackage{langsci-osl} % macros specific to Open Slavic Linguistics; PLACED DIRECTLY IN localcommands.tex
\usepackage{siunitx}
\usepackage[linguistics]{forest}
% \usepackage{caption} %borik (?)
\usepackage{pifont} %geist

% 08_muellerreichau

% \usepackage{caption}
\usepackage{subcaption}
\usepackage{cleveref}

\usepackage{drs}
\usepackage{diagbox}

% 09_simik

\usepackage{multirow}
\usepackage{tabto}

% 10_stegovec

\usetikzlibrary{arrows,patterns}
\usepackage{listings}
% \usepackage{multirow}

% 01_gehrke-simik

% \usepackage{comment}

% 13_wagiel

% \usepackage{pifont} % for checkmarks (\ding{51}, \ding{55})


\usepackage{langsci-branding}

   % \newcommand*{\orcid}{}


\makeatletter
\let\thetitle\@title
\let\theauthor\@author
\makeatother

\newcommand{\togglepaper}[1][0]{
   \bibliography{../localbibliography}
   \papernote{\scriptsize\normalfont
     \theauthor.
     \titleTemp.
     To appear in:
     E. Di Tor \& Herr Rausgeberin (ed.).
     Booktitle in localcommands.tex.
     Berlin: Language Science Press. [preliminary page numbering]
   }
   \pagenumbering{roman}
   \setcounter{chapter}{#1}
   \addtocounter{chapter}{-1}
}

\newbool{bookcompile}
\booltrue{bookcompile}
\newcommand{\bookorchapter}[2]{\ifbool{bookcompile}{#1}{#2}}

\AtBeginDocument{\nogltOffset} %removes extra spacing before trans line in examples - SPECIFIC TO OPEN SLAVIC LINGUISTICS, PREVIOUSLY AGREED UPON


%%%%%%%%%%%%%%%%%%%%%%%%%%%%%%%%%%%%%%%%%%%%%%%%%%%%%%%%%%%%%%%%%%%%%
%%      File: langsci-osl.sty
%%    Author: Radek Simik and Language Science Press (http://langsci-press.org)
%%      Date: 2019-02-18
%%   Purpose: This file contains macros for the Open Slavic Linguistics at Language Science Press series and is included 
%%            into langscibook.cls
%%  Language: LaTeX
%%   Licence: 
%%%%%%%%%%%%%%%%%%%%%%%%%%%%%%%%%%%%%%%%%%%%%%%%%%%%%%%%%%%%%%%%%%%%%

% FIXING GRAY

\definecolor{lsDOIGray}{cmyk}{0,0,0,0.45}

% PACKAGES REQUIRED

\RequirePackage{relsize}
\RequirePackage{stmaryrd}
\RequirePackage{array}
\RequirePackage{tabularx}

% % KEYWORDS

% \newcommand{\keywords}[1]{\noindent\textbf{Keywords:} {#1}}

% SEMANTICS MACROS - GENERAL FOR ALL PAPERS

\newcommand{\sib}[1]{$\llbracket${#1}$\rrbracket$} % = semantic interpretation brackets; puts text into double square brackets
\newcommand{\stb}[1]{\langle{#1}\rangle} % = semantic type brackets; puts stuff into semantic type brackets; necessary to use in the form $\st{}$
\newcommand{\cnst}[1]{\textsf{\relsize{-1}\MakeUppercase{#1}}} %creates sans-serif small-caps, used for typesetting logical constants which have no other appropriate symbols (such as Greek letters)

% MINUS HSPACE MACRO

\newcommand{\minsp}[1]{#1\hspace{-2.5pt}} % use for non-glossed elements in glossed examples (typically stars, brackets, slashes, ...)

% CITATION MACROS

\newcommand{\citeposst}[1]{\citeauthor{#1}'s (\citeyear{#1})} % produces Chomsky's (1995)
\newcommand{\citeposstpg}[2]{\citeauthor{#1}'s (\citeyear[#2]{#1})} % produces Chomsky's (1995: page)
\newcommand{\citepossalt}[1]{\citeauthor{#1}'s \citeyear{#1}} % produces Chomsky's 1995
\newcommand{\citepossaltpg}[2]{\citeauthor{#1}'s \citeyear[#2]{#1}} % produces Chomsky's 1995: page

% BARPLOT

\usepackage{pgfplots}
\pgfplotsset{compat=1.18} 
\newcommand{\barplot}[5]{%
  \begin{tikzpicture}
    \begin{axis}[
	xlabel={#1},  
	ylabel={#2}, 
	axis lines*=left, 
        width  = {#3\textwidth},
	height = .3\textheight,
    	nodes near coords, 
	xtick=data,
	x tick label style={},  
	ymin=0,
	symbolic x coords={#4},
	]
	\addplot+[ybar,lsRichGreen!80!black,fill=lsRichGreen] plot coordinates {
	    #5
	}; 
    \end{axis} 
  \end{tikzpicture} 
}



%%%%%%%%%%%%%%%%%%%%%%%%%%%
%%% 04_filip %%%

% \newcommand{\mC}[1]{\multicolumn{1}{c}{#1}} % multicolumn with a single column; ALREADY DEFINED

\usetikzlibrary{arrows, arrows.meta}
\usetikzlibrary{positioning, decorations, 
                decorations.pathreplacing, calligraphy}

%%%%%%%%%%%%%%%%%%%%%%%%%
%%% 08_muellerreichau %%%

\captionsetup[subfigure]{subrefformat=simple,labelformat=simple}
\renewcommand\thesubfigure{(\alph{subfigure})}

%%%%%%%%%%%%%%%%%%%%%%%%%%%
%%% 09_simik

\newcommand{\citepossts}[1]{\citeauthor{#1}' (\citeyear{#1})} % produces Rawlins' (2008)
\newcommand{\citepossalts}[1]{\citeauthor{#1}' \citeyear{#1}} % produces Rawlins' 2008


%%%%%%%%%%%%%%%%%%%%%%%%%%%%%%%
%%% 10_stegovec

\newcommand{\poscite}[1]{\citeauthor{#1}'s (\citeyear{#1})}

%extra commands
\newcommand{\fst}{\textsc{1p}}
\newcommand{\snd}{\textsc{2p}}
\newcommand{\trd}{\textsc{3p}}
\newcommand{\refl}{\textsc{refl}}
\newcommand{\excl}{\textsc{excl}}
\newcommand{\incl}{\textsc{incl}}
\newcommand{\sg}{\textsc{sg}}
\newcommand{\pl}{\textsc{pl}}
\newcommand{\dl}{\textsc{dl}}
\newcommand{\fem}{\textsc{f}}
\newcommand{\masc}{\textsc{m}}
\newcommand{\neut}{\textsc{n}}
\newcommand{\ImpCl}{IMPERATIVE}


\newcommand{\tikzmark}[1]{\tikz[overlay,remember picture] \node (#1) {};}
\forestset{M/.style={
		for tree={s sep=5mm, inner sep=0.5pt, l=5mm,
			calign=fixed edge angles,
			calign primary angle=-65,
			calign secondary angle=65},
	},
	Mid/.style={
		for tree={s sep=0mm, inner sep=0pt, l=0mm,
			calign=fixed edge angles,
			calign primary angle=-62.5,
			calign secondary angle=62.5},
	},
	word/.style={before typesetting nodes={content=\emph{##1}},
		no edge,l=0,for parent={l sep=0}},
	feat/.style={before typesetting nodes={content={{[}##1{]}}},
		no edge,l=0,for parent={l sep=0}},
	feat2/.style={before typesetting nodes={content={##1}},
		no edge,l=0,for parent={l sep=0}},
	llap/.style={
		tikz+={%
			\edef\forest@temp{\noexpand\node[\option{node options},
				anchor=base east,at=(.base east)]}%
			\forest@temp{#1\phantom{\option{environment}}};
		}
	},
	rlap/.style={
		tikz+={%
			\edef\forest@temp{\noexpand\node[\option{node options},
				anchor=base west,at=(.base west)]}%
			\forest@temp{\phantom{\option{environment}}#1};
		}
	}
}

\makeatletter
\newcommand{\coloruline}[2]{%
	\newcommand\temp@reduline{\bgroup\markoverwith
		{\textcolor{#1}{\rule[-0.5ex]{2pt}{0.4pt}}}\ULon}%
	\temp@reduline{#2}%
}

\newcommand{\coloruuline}[2]{%
	\UL@protected\def\temp@uuline{\leavevmode \bgroup
		\UL@setULdepth
		\ifx\UL@on\UL@onin \advance\ULdepth2.8\p@\fi
		\markoverwith{\textcolor{#1}{\lower\ULdepth\hbox
				{\kern-.03em\vbox{\hrule width.2em\kern1\p@\hrule}\kern-.03em}}}%
		\ULon}%
	\temp@uuline{#2}%
}

\newcommand{\coloruwave}[2]{%
	\UL@protected\def\temp@uwave{\leavevmode \bgroup 
		\ifdim \ULdepth=\maxdimen \ULdepth 3.5\p@
		\else \advance\ULdepth2\p@ 
		\fi \markoverwith{\textcolor{#1}{\lower\ULdepth\hbox{\sixly \char58}}}\ULon}
	\font\sixly=lasy6 % does not re-load if already loaded, so no memory drain.
	\temp@uwave{#2}%
}
\makeatother

%%%%%%%%%%%%%%%%%
%%% 13_wagiel %%%

\newcommand\irregularcircle[2]{% radius, irregularity
  \pgfextra {\pgfmathsetmacro\len{(#1)+rand*(#2)}}
  +(0:\len pt)
  \foreach \a in {10,20,...,350}{
    \pgfextra {\pgfmathsetmacro\len{(#1)+rand*(#2)}}
    -- +(\a:\len pt)
  } -- cycle
}


% IF POSSIBLE, CAN WE USE THIS SETTING JUST FOR 13_wagiel?

% \forestset{
% default preamble={for tree={s sep=10mm, inner sep=0, l=0}},
% fairly nice empty nodes/.style={
%         delay={where content={}{shape=coordinate,
%               for siblings={anchor=north}}{}},
% for tree={s sep=4mm}},
% }

   %% hyphenation points for line breaks
%% Normally, automatic hyphenation in LaTeX is very good
%% If a word is mis-hyphenated, add it to this file
%%
%% add information to TeX file before \begin{document} with:
%% %% hyphenation points for line breaks
%% Normally, automatic hyphenation in LaTeX is very good
%% If a word is mis-hyphenated, add it to this file
%%
%% add information to TeX file before \begin{document} with:
%% %% hyphenation points for line breaks
%% Normally, automatic hyphenation in LaTeX is very good
%% If a word is mis-hyphenated, add it to this file
%%
%% add information to TeX file before \begin{document} with:
%% \include{localhyphenation}
\hyphenation{
    par-a-digm
    Cart-wright %filip
    Truswell %docekal
    Shcher-bi-na %simik
    Mün-chen %gehrke
    Mo-ni-ka %gehrke
    spi-sov-né %gehrke
    li-te-ra-tur-no-go %gehrke
    russ-koj %gehrke
    so-ver-šen-no-go %gehrke
    russ-kom %gehrke
    sla-vjan-sko-mu %gehrke
    Zeit-schrift %gehrke
    Na-ta-sha %gehrke
    Zu-stands-pas-siv %gehrke
    Um-gangs-spra-che %gehrke
    Bra-so-vea-nu %geist
}

\hyphenation{
    par-a-digm
    Cart-wright %filip
    Truswell %docekal
    Shcher-bi-na %simik
    Mün-chen %gehrke
    Mo-ni-ka %gehrke
    spi-sov-né %gehrke
    li-te-ra-tur-no-go %gehrke
    russ-koj %gehrke
    so-ver-šen-no-go %gehrke
    russ-kom %gehrke
    sla-vjan-sko-mu %gehrke
    Zeit-schrift %gehrke
    Na-ta-sha %gehrke
    Zu-stands-pas-siv %gehrke
    Um-gangs-spra-che %gehrke
    Bra-so-vea-nu %geist
}

\hyphenation{
    par-a-digm
    Cart-wright %filip
    Truswell %docekal
    Shcher-bi-na %simik
    Mün-chen %gehrke
    Mo-ni-ka %gehrke
    spi-sov-né %gehrke
    li-te-ra-tur-no-go %gehrke
    russ-koj %gehrke
    so-ver-šen-no-go %gehrke
    russ-kom %gehrke
    sla-vjan-sko-mu %gehrke
    Zeit-schrift %gehrke
    Na-ta-sha %gehrke
    Zu-stands-pas-siv %gehrke
    Um-gangs-spra-che %gehrke
    Bra-so-vea-nu %geist
}

   \boolfalse{bookcompile}
   \togglepaper[4]%%chapternumber
}{}

\begin{document}
\maketitle

% \input{template-osl.tex}
% Just comment out the line above (place % at its beginning) and start writing your own document
% WRITE DIRECTLY HERE, INTO main.tex, DO NOT START YOUR OWN tex FILE!

% START WRITING HERE

\section{Main data and questions} \label{fil:sec:introduction}




%Klimek-Jankowska Imperfective_and_Perfective_Habituals_in

%https://plato.stanford.edu/entries/two-dimensional-semantics/#Role2D


Generic sentences are sentences like \textit{Dogs bark, Sugar dissolves in water, 
Sam walks to school on Mondays}. This chapter focuses on their encoding in Slavic languages that have a dedicated marker on the verb 
enforcing only a generic (or habitual) reading of a sentence, besides 
using imperfective, and also perfective, verbs for this purpose in an  appropriate context that triggers their generic (or habitual) interpretation. Its instantiation in Czech (West Slavic language) 
is treated as a suffix in standard reference grammars 
(see \citealt[184]{petr1986mluvnice},
\citealt{dahl1995}, 
\citealt{Kosek2014}, \citealt{nuebler2017}, for instance); 
some Bohemists treat it as an infix (e.g., \citealt{kucera1981aspect}, \citealt{bily1986iterative}).
In \citet[23]{Genis2021}, it is treated as an infix across different Slavic languages.\footnote{\citet[23]{Genis2021}: ``[T]he most common infix in this [habitual, frequentative/iterative, HF] function across languages seems to be one that might be rendered -/va/-."} 
Henceforth, I will use \textit{-va-} as its citation form (as is also common in Bohemistics), and refer to it as the \textsc{morpheme} \textit{-va-} or \textsc{marker} \textit{-va-}, 
given that trying to resolve the suffix-infix controversy would  
go well beyond the scope of this contribution.

This marker is attested in all Slavic languages (see \citealt{Genis2021} for examples from different Slavic languages). However, 
it is productive only in West Slavic languages, specifically in Czech, Slovak and to a lesser extent in Polish, and commonly used in both spoken and written language (see e.g., \citealt[62]{sirokova1963}, \citealt[27]{Comrie1976}, \citealt{kucera1981aspect}, \citealt{bily1985preliminaries}, \citealt{bily1986iterative,petr1986mluvnice,Genis2021}, among others).\footnote{The productivity of \textit{-va-} in Czech is also evident in the observation that children use verbs on which it occurs early in their language acquisition, even before they master the grammatical perfective/imperfective distinction \citep{chmelickova2005}.}
In East and South Slavic languages it is also attested, but 
it is reported to be largely unproductive, at least in their standard varieties.\footnote{In standard spoken and written Russian, the generic marker only occurs on some verbs that are taken to be marginal in use (\citealt[fn.~1 on p.~27]{Comrie1976}, who cites \citealt[405--407]{isacenko1962}). They are attested in fixed collocations (e.g., \textit{V žizni tak byvaet} [gloss: in life so be.\textsc{gen.3.sg.prs}] `That's how life goes. / That's life.'). \citet[87]{Dickey2000}
claims that in standard Serbian explicitly ``habitual'' (in his terms, generic in my terms)
verbs can be formed with the morpheme \textit{-va-} (which is an infix, according to him) from biaspectual verbs, e.g.: \textit{kršćavati} (from \textit{krstiti}\textsuperscript{PFV/IPFV} `baptize’); \textit{noćivati}
(from \textit{noćiti}\textsuperscript{PFV/IPFV} `spend the night’); \textit{ručavati} (from \textit{ručati}\textsuperscript{PFV/IPFV} `eat lunch’);
\textit{večeravati} (from \textit{večerati}\textsuperscript{PFV/IPFV} `eat dinner’). However, Dickey's claim is contested in \citet[26]{Genis2021}: ``[A]ccording to our
observations, these verbs do not always have a habitual interpretation but can be
used in a durative meaning, as well; other than their base verbs they are explicitly
imperfective."}

Here, the main focus is on the instantiation of this Slavic marker in Czech. The constraints on the length of this chapter prohibit me from
adducing relevant examples from other Slavic languages. To the extent that the relevant generic verb forms exist in other Slavic languages, e.g., \textit{pisyvat'} in Russian, meaning `to write repeatedly, sporadically, as a rule' (albeit judged as archaic),
it is reasonable to expect that the arguments made here for 
the corresponding Czech forms, e.g., \textit{psávat}, should apply to them, as well. 

Let me start with a few examples, given below:  

\ea  %\jambox*{Czech (West Slavic)}
\ea\label{psi.IPF}
\gll Psi \textit{štěkají}.\\
dogs bark \\ \hfill  \textsc{ipfv}     
\glt `Dogs bark.' [Intended interpretation: `Dogs as a kind have the characterizing property of barking.'] 
\ex\label{psi.PF}
\gll Psi se \textit{poštěkají}, že ani nevíte, jak k tomu došlo. \\ 
dogs \textsc{refl} \textsc{pref}.bark that even \textsc{neg}.know how to it came\\   \hfill  \textsc{pfv}
\glt `Dogs (will) bark (out) at each other, and one hardly knows why.' 
\ex\label{psi.secondIMPFV}
\gll Psi \textit{poštěkávají} 
ve spánku, vrčí a cukají packami. \\ 
dogs \textsc{pref}.bark.\textsc{ipfv} in sleep, growl and twitch paws\\ \hfill \textsc{ipfv}
\glt `Dogs bark now and then while sleeping, growl and twitch their paws.' [i.e., bark a little/on and off] 
\ex\label{psi.GEN}
\gll Psi \textit{štěkávají} na ty, které neznají.    \\
dogs bark.\textsc{gen} at those whom \textsc{neg}.know \\  \hfill \textsc{gen}
\glt `Dogs bark at those whom they don't know.'		
\z
\z
\noindent
As the English translations of the four Czech generic sentences above show, one and the same English verb form \textit{bark} (3rd person plural, present tense) corresponds to four different verb forms in Czech.
All four share the same root \textit{-štěk-} `bark'. (\ref{psi.IPF}) is headed by
a primary (root) imperfective verb, 
formally unmarked for imperfectivity, while 
(\ref{psi.secondIMPFV}) by a secondary imperfective verb,
which is derived by means of the imperfectivizing suffix
(glossed with \textsc{ipfv}) from a perfective base, realized in the perfective verb in (\ref{psi.PF}). 
In languages with the perfective/imperfective distinction like Slavic or Romance languages, the semantics of imperfectivity is 
taken to have close affinities to genericity (and habituality), based on the observation that imperfective verbs have common contextually determined generic readings, as we see in (\ref{psi.IPF}) and (\ref{psi.secondIMPFV}). 

However, there is no incompatibility between the semantics of perfectivity and genericity (and habituality). Perfective verbs are also used in generic sentences, and this use is not as unusual or rare as some
assume
(\textit{pace} \citealt{Comrie1976}, \citealt{dahl1995}, i.a.). For instance, 
in Romance languages, and in West Slavic languages (see e.g., \citealt{filip1997suigenva}), perfective 
verbs have common contextually determined generic readings, as is illustrated by (\ref{psi.PF}).
Here, we have a perfective verb derived by means of the prefix \textit{po-} from the stem of the primary (root) imperfective verb  \textit{štěkat} `to (be) bark(ing)' (used in (\ref{psi.IPF})).\footnote{For a recent study on the use of present tense perfective verbs in Polish generic sentences see \citet{wiemer2025}.}

Most importantly, (\ref{psi.GEN}) has an obligatory generic interpretation in all contexts, and unlike 
the (im)perfective verb forms in (\ref{psi.IPF}) and (\ref{psi.secondIMPFV}), (\ref{psi.GEN}) can never be used 
in episodic contexts. This is due to the morpheme 
\textit{-va-} on its main verb. It constitutes a sufficient, but not a necessary, condition for the 
encoding of generic statements, given that, as we have 
just seen, such statements can also be expressed by verbs without this morpheme, in an appropriate context. 

One of the main goals of this chapter is to argue that the 
morpheme\linebreak \textit{-va-} (as in (\ref{psi.GEN})) is best viewed as a generic morpheme sui generis (also building on some previous arguments for Czech in \citet{filip1993berkeleyva,filip1994mitva,filip1997suigenva, Mendia2018,filip2021cologneva}, i.a.).
This argument implies a rejection of the prevalent traditional view on which this morpheme is a marker of a (possibly unbounded) plurality of situations of the same type (see \citealt{Genis2021}, and specifically for Czech see e.g., \citealt{petr1986mluvnice,dahl1995}, \citealt{nuebler2017}, and references therein), 
and as such, it is treated as a marker
that falls under imperfective aspect and delimits an (iterative) \textit{Aktionsart} (\citealt{poldauf1964}, for instance), or
as a kind of pluralization operator deriving a predicate of plural situations. 
While the notion of a plurality of situations is relevant to 
the meaning of the morpheme \textit{-va-}, it is neither sufficient nor necessary to ground its meaning, as I will show. Given the erroneous idea that this marker 
encodes a mere plurality of situations,
it is traditionally referred to as the ``iterative'' marker, and also the ``multiplicative'', ``frequentative'' or ``repetitive'' marker.
That its most common label ``iterative'' is a misnomer at best, and a category error at worst, is evident in the simple fact that this morpheme is incompatible with iterative adverbials like `three times' or `several times' (see section \sectref{semprag}).
It will also be
shown that the marker \textit{-va-} does not have a meaning akin to the adverb of quantification 
`usually' or the determiner quantifier `most' (\textit{pace} \citealt{dahl1995}, i.a.), as in `in most situations'. 

Related to the main argument to be made here, 
let me make two clarifications.  
First, terminology for morphemes and grammatical constructions that are used in sentences that express regularities, patterns, or habits -- in short generalizations of various kinds -- is not standardized and, as suggested in the previous paragraph, often misleading.  
Iterativity, multiplicativity,  frequentativity, repetitivity, usuality, pluractionality, and the like are often used interchangeably,
and also confounded with habituality and genericity. 

When it comes to habituality, it
is either viewed as separate from genericity, equated with it (e.g., \citealt{Comrie1976}), or subsumed under genericity as its special case (e.g., \citealt{krifka1995genericity}).\footnote{For instance, \citet{Comrie1976}
uses ``habituality'' for 
all the uses of imperfective forms that are not episodic.
He divides the semantic domain of imperfectivity into a
``habitual'' subdomain and an episodic subdomain. However, 
he confusingly refers to the episodic subdomain as ``continuous'',
and divides it into two further subdomains: namely, ``nonprogressive'' and ``progressive''. The category ``continuous'' in \citep{Comrie1976} is notionally hardly coherent, and 
formally speaking, there seems to be no language with formal markers for it. For further criticism of Comrie's classification see \citet{filip1997suigenva}.}
%The status of both `habituality' and `genericity' with respect to the categories of the TAM system is a matter of different opinions. 
%Judging from Comrie's examples, `habituality' in his sense is co-extensional with `(characterizing) genericity' here, that is, it covers the meaning of (characterizing) generic sentences that are analyzed by means of \cnst{Gen} in \citet{krifka1995genericity} . That is, it covers the three main types of generalizations, given in (i)-(iii) above, which are distinguished in the genericity literature. 

In its widest sense, the term ``habitual'' 
is used for all verbs, VPs or sentences that are not episodic.
Episodic sentences are true relative to a given world and reference time, they do not describe regularities. In this widest sense, habituality has been popularized by  \citet{Comrie1976}, 
and is commonly used for one subdomain of imperfectivity, 
as opposed to its other, namely, episodic subdomain in aspect studies and 
traditional/descriptive
linguistics. Habituality in this sense is, however, problematic, because it subsumes uses of sentences describing what is not generally viewed as a matter of habit 
at all (\citealt[716]{Lyons1977}, i.a.), i.e., habit in the sense of regularities of action of typically human, or at least animate, individuals, as in  \textit{Fred smokes after dinner}. Some examples are: \textit{Glass breaks easily} (a disposition), \textit{Soap is used to remove dirt} (a function), \textit{Bishops move diagonally} (a rule of a game),  \textit{Kim works for the government} (an occupation). Therefore, a better term to cover all types of sentences that are not episodic is \textit{generic}, which is prevalent in formal 
semantics studies that are informed by philosophy, computer science and psychology. 

The Czech marker \textit{-va-}, as in (\ref{psi.GEN}),
despite being occasionally labeled ``habitual'' (e.g., \citealt{dahl1995,danaher2003semantics,karlik2017}),
can formally mark not only sentences that describe habits proper
(as in \textit{Fred smokes after dinner}), 
but also generic sentences with kind reference, as in our initial example (\ref{psi.GEN}), and also in other types of generic sentences that describe what is not a matter of habit (see e.g., (\ref{sea.va}) below). 

Therefore, here, for the sake of simplicity, I will use the term \textit{generic} as a cover term for all statements that are \textit{not} episodic, while habituality 
will be treated as a special case of genericity (as in e.g., \citealt{krifka1995genericity}). The term \textit{habitual sentence} 
will be reserved for a subclass of generic sentences that concern habits proper, namely regularities of action of typically human, or at least animate, individuals (e.g., \textit{Fred smokes after dinner}). 

The second point of clarification concerns the confounding of the Slavic generic marker \textit{-va-} (as in (\ref{psi.GEN})) with the imperfectivizing suffix (as in (\ref{psi.secondIMPFV})). 
Their confounding has to do with the fact that
some of their allomorphic variants seem 
virtually indistinguishable, and both are used in generic (and habitual) contexts.\footnote{Occasionally this confounding may be also motivated
and reinforced by speculations about the diachronic 
development of one of the morphemes
from the other, which, however, seem unfounded. For instance, \citet{Kuznetsova2016} assume, without any evidence and only
with reference to some assumptions by \citet{kuznecov1953},
that habituality or iterativity is diachronically the original meaning of the Russian secondary imperfectivizing suffix, standardly cited as \textit{-yva/iva-}. According to \citet[259]{kuznecov1953},
it developed as the result of a reanalysis of the 
the habitual or ``iterative'' morpheme \textit{-va-} of the verb \textit{byvati} `to be usually, typically, normally' or `to tend to be':
\textit{by-va-ti > b-yva-ti}. Hence, according to him, the imperfectivizing suffix \textit{-yva/iva-} originally only had the habitual or ``iterative'' meaning component.}
While the generic marker \textit{-va-} enforces a generic reading in all its occurrences, the imperfectivizing suffix
on secondary imperfectives is compatible with their
episodic or generic reading. 
For instance, only the imperfectivizing suffix, but not the generic \textit{-va-}, can be used in 
sentences with a contextually determined progressive (episodic) reading. 
I will address the differences between these two
morphemes in detail in \sectref{basicmorphology},
and touch on their differences also elsewhere in this chapter. 

With these basic clarifications and caveats out of the way, the main question to be pursued in this chapter is as follows:\\

\noindent
\textsc{Question 1}:  What properties of the marker \textit{-va-} justify its status as a marker of genericity? \\

\noindent
The marker \textit{-va-}
constitutes a sufficient, but not a necessary, condition for the  expression of generic (and habitual) statements. Such markers
are attested in a number of typologically distinct languages. A short list, which also includes Czech, is 
given in \citet [421,~fn.~8]{dahl1995}.\footnote{\citet[421, fn.~8]{dahl1995} gives the following list:  
Arabic (Classical), Akan, Catalan, Czech, Didinga, German, Guarani, Hungarian, Kammu, Limouzi, Montagnais, Sotho, Spanish, Swedish, Swedish Sign Language, Yucatec Maya, Zulu.} In 
fact, Dahl 
uses the Czech marker \textit{-va-} as a paradigm example of this class of markers, which, according to him, fail to qualify as generic markers. Instead, they are to be treated as what he calls 
``habitual tense-aspect markers'' that function as a kind of quantifier over situations with the semantics of `most' or `usually'. 
For \citet{dahl1995}, this claim bolsters his overarching argument concerning the reflexes of the episodic/generic distinction in the grammar of natural languages: namely, if all such verb markers like the Czech \textit{-va-} are not generic markers, but rather habitual markers (`habitual' in the sense of \citealt{dahl1995}) that fall either under tense or imperfective aspect,
% (see also
% e.g.,  \citealt[26ff.]{Comrie1976}, \citealt[40]{comrie1985tense}, \citealt{Dahl1985}, i.a.) 
then ``however 
sad this may be for those of us to whom the episodic/generic distinction is dear, this distinction more often than not is only indirectly reflected in speakers' choices between grammatical markers" of aspect or tense \citep[425]{dahl1995}.

These are the claims that will be contested here.
Contrary to \citet{dahl1995}, I will show that the Czech marker \textit{-va-} 
has properties that preempt its subsumption under 
tense or imperfective aspect (also \textit{pace} \citealt{cable2022}).
Its  meaning
cannot be reduced to a quantifier over situations with 
the meaning amounting to `most', or `usually', 
nor any other single expression of quantity or quantification,
also contrary to \citet{dahl1995}.

Having established what the marker \textit{-va-} is not, I will show that it derives
verbs denoting predicates that 
exhibit three hallmark properties of 
generic predicates (see e.g., \citealt{carlson1977unified,krifka1995genericity}): namely, 
(i) they are aspectually stative,
(ii) have a modal (intensional) import which is tied to the quintessentially predictive force of generics, and as such their evaluation requires access to other possible worlds than just the actual one, 
and (iii) they express generalizations that crucially involve reasoning with exceptions.
This then leads me to the next two questions:\\

\noindent 
\textsc{Question 2}: What kind of generic sentence does the generic marker \textit{-va-} delimit? \\

\noindent
\textsc{Question 3}: What is the relation of the  generic marker \textit{-va-}
to the phonologically null generic \cnst{Gen} operator?\\

\noindent
Given that there is no comprehensive semantic analysis of genericity that is commonly accepted, when it comes to the foundational assumptions and terminology in the domain of genericity (and habituality), 
I will rely 
on the theoretical background
of the quantificational theory of genericity (see e.g., \citealt{krifka1995genericity}). 
Its key hypothesis 
is that the semantics of generic sentences can be represented
by means of a logical (tripartite) structure, where 
the generic import is captured by means of the phonologically null, modal generic \cnst{Gen} operator, which is akin to overt adverbs of quantification (Q-adverbs) (see  \sectref{CharGenericity}). 
The quantificational theory of genericity
will here serve as a useful point of reference for examining theoretical and empirical issues raised by the generic morpheme \textit{-va-}. It has the advantage that it has a wide interdisciplinary scope, drawing on insights not only from semantics and pragmatics, but also on insights regarding the nature of 
generalizations, law-like statements and reasoning with exceptions 
in philosophy, psychology, computer science and other areas of cognitive science. 

As I will show, just as \cnst{Gen},
the generic morpheme \textit{-va-} functions as a quantifier over individuals and/or situations, and exhibits variable binding properties that are akin to those of \cnst{Gen}. Moreover, the meaning of both \cnst{Gen}
and \textit{-va-} cannot be reduced to any single quantifier, e.g.,  `usually',  `most', or any other quantifier or quantity expression. At the same time, the Slavic generic morpheme 
differs from \cnst{Gen} in so far
as it has two uncancellable inferences: namely, that 
(i) there are verifying instances in the actual world that count as evidence for the truth of the generalization, and (ii) there are actual or at least potential exceptions or counterexamples to the generalization. The former inference is shared with what are taken to be so-called `habitual'  markers in other languages (see e.g., \citealt{cable2022} on the realization inference of Tlingit habitual morphology), and it is also known as the actualization entailment in the genericity literature (see e.g., \citealt{green2000aspectual}, \citealt{hacquardDISS}). 

The latter inference
is clearly manifested in the observation that 
\textit{-va-} is semantically anomalous or false in generic sentences that have no exceptions in all the possible worlds (such as universal laws of nature), and it must be used in generic sentences that have known counterexamples
and would be false without it, all else being equal (see \sectref{mustbeused}). 
Generally, as I will show, 
the truth conditions and felicity of generic sentences that are formally marked with \textit{-va-} crucially depend on the speaker's reasoning with exceptions, which reflect the 
speaker's commitment
to the strength of the generalization calibrated in terms of exceptions that (the speaker believes) it has or might have.
In effect, generic sentences formally marked with \textit{-va-} are 
endowed with two layers of modality (intensionality): one that is tied to their generic, law-like, (g)nomic force, similarly to \cnst{Gen},
and the other to reasoning with exceptions, which 
may be characterized in terms of the speaker-oriented epistemic and doxastic modality (exploiting general pragmatic principles of interpretation). 

%The Slavic generic marker \textit{-va-} carves out a subdomain from the whole domain of generic sentences that the null generic operator \cnst{Gen} covers.
The presence of the generic marker \textit{-va-}, and its contrastive absence, in generic sentences differentiates two main types of generic sentences in Slavic languages.
Some have argued that different linguistic forms for the encoding of generic statements differentiate different subtypes of generic sentences, each distinguished by different clusterings of formal properties, and possibly requiring different semantic/ontological commitments, and separate semantic/ontological models for their interpretation
(see e.g., \citealt{greenberg2003manifestations,pelletier-equal2009,krifka2012definitional, boneh2013hab, boneh2008habituality}). 
Hence, the properties of the Slavic generic marker 
directly bear on the fundamental question in the domain of genericity 
(\citealt{carlson1977unified,carlson1995truth}): \\

\noindent
\textsc{Question 4}: Can we provide a unified semantic analysis for all generic sentences? \\

\noindent 
A conclusive answer to the above question is still outstanding. A unified analysis for all generic sentences is desirable, and hypothesizing that
generic sentences might constitute a single class, 
markers that are dedicated to the encoding of generic statements, such as the generic marker \textit{-va}-, provide some of the best empirical data for evaluating the claims that address the above 
question.
%If a unified analysis turns out not viable for the whole domain of generic sentences, given the properties of overt generic markers, we may have to split it into subdomains in order to provide adequate semantic analyses  for all the relevant data. 

The proposal advanced here that the morpheme \textit{-va-} is a generic marker sui generis, as well as the empirical and theoretical arguments in its support, 
have potentially significant consequences for the theory of genericity. First, if the morpheme \textit{-va-} in Czech, 
in other Slavic languages where it is attested, 
is a generic marker, as I argue, it would shed at least some doubts on the widespread belief in linguistics, philosophy and psychology that natural languages have no dedicated markers of genericity (see e.g.,  \citealt{leslie2008}, \citealt{Johnston2012}, \citealt{LeslieGelman2012}, \citealt{Rhodes2018}, \citealt{cohen2022}, \citealt{Schiller2023}, to give just a few more recent examples),
or only a negligible number \citep{dahl1995}.\footnote{For instance, \citet[541]{Schiller2023} observes: ``Generics are present in every known language, but in no language are generic sentences `marked' with an explicit generic operator. Generics are studied (by philosophers, linguists, and developmental psychologists) because of their abstruse formal properties and the important role they play in cognitive development.''} 
Second, it would support the claims independently made elsewhere (see e.g., \citealt{carlson1977unified}, \citealt{schubert1989pelletier}, \citealt{filip1997suigenva}, \citealt{pelletier1997generics}, i.a.),
and mostly on semantic grounds, 
that the episodic/generic 
distinction is relevant for the grammar of natural language
(\textit{pace} e.g., \citealt{dahl1995}), 
and bolster the view on which genericity is a category in its own right, 
and specifically independent from the tense, modality and aspect (TMA) system. 



%This view is compatible withthe undeniable fact that there are affinities between genericity and speakers' choices between grammatical markers of tense and aspect in the formulation of generic sentences (e.g., the simple present tense in English preferably has a generic interpretation, the imperfective form is preferred to the perfective one in Russian in generic sentences).


This chapter is structured as follows. \sectref{CharGenericity}
introduces the theoretical background of the quantificational theory of genericity.
%, including the key properties of generic sentences, which will here serve as a framework for examining theoretical and empirical issues raised by the Slavic generic morpheme \textit{-va-}. 
In \sectref{basicmorphology}, the formation of generic 
verbs by means of the generic morpheme \textit{-va-}, and its place in the Czech verb system will be outlined in a theory neutral
way. Its properties will be also shown not to be aligned with markers of tense or imperfective aspect. 
\sectref{semprag}
focuses on its main semantic and pragmatic properties, 
which together with \sectref{basicmorphology} invalidate the grounds for the claim made by \citet{dahl1995}
that the Czech morpheme \textit{-va-} fails to qualify as a generic marker.
\sectref{proposal} addresses the semantic similarities and differences between 
the phonologically null generic operator \cnst{Gen} and 
the Czech generic marker \textit{-va-}. 
The marker \textit{-va-} introduces a 
modal (intensional) generic operator \cnst{va} into the logical representation 
that differs from \cnst{Gen}  in so far as it contributes the uncancellable inference regarding the  existence/possibility of exceptions or counterexamples to the generalization which, moreover, must have been realized. The use of the generic \textit{-va-} form, when the corresponding non-generic form might have also been used, gives rise to a rank order (see \citealt{hornabbott2012,horn2017}, and references therein), where the non-generic form is the strong alternative (as it allows no exceptions when used in generic sentences) and the generic form the weak one (given that it requires that there be actual or possible exceptions). The rank order generates either the speaker-oriented certainty or ignorance inference concerning the likelihood of exceptions (a kind of quantity-based implicature), which may be characterized in terms of epistemic and doxastic modality, exploiting general pragmatic principles of interpretation.


\section{Genericity: Kind reference and generic sentences} \label{CharGenericity}

The domain of genericity is divided into two independent, 
but related, phenomena: namely, \textsc{kind reference} and \textsc{characterizing generic sentences} (\citealt{krifka1995genericity}, \citealt{carlpell1995generic}). Some examples are given below:\\


\ea
\ea \textit{Ravens} are widespread.      	\jambox*{\textsc{kind reference}}			\ex {\label{dodos} \textit{Dodos} became extinct in the 17th century.}	
\ex \textit{Bronze} was invented as early as 3000 B.C.		
\ex {\label{the.dodo} \textit{The dodo} / \textit{*A dodo} became extinct in the 17th century.}	
\z
\z

\ea
\ea {\label{ravens.black} \textit{Ravens} are black.}    \jambox*{\textsc{characterizing generic sentences}}
\ex {\label{dogs.bark} \textit{Dogs} bark / \textit{A dog} barks.} 
\ex {\label{paul.pc.prs} Paul plays chess on Saturdays.} 
\z
\z

\noindent
There are both formal and semantic reasons for this
division. The category of \textsc{kind reference} is motivated by the existence of kind-level predicates (KLPs) like \textit{extinct, widespread, invent} that directly select \textsc{kind-denoting} terms for one of their arguments (in italics in the above examples).\footnote{A given language may impose 
specific requirements on the form of a kind-denoting term. In English, for instance, a singular indefinite is often unacceptable.} 
KLPs attribute a property to a whole kind and one which
does not hold of its individual members. For instance, in (\ref{dodos}), the property of being extinct holds of the whole kind \textsc{dodo}, and it makes no sense to attribute it to any individual dodo or a closed set of dodos.	
 
The main topic of this chapter concerns \textsc{characterizing generic sentences}, or 
\textsc{generic sentences} for short. 
Their generic import often arises from the combined meanings of the subject NP/DP and VP, and possibly also in interaction with pragmatic, pro\-s\-odic (focus), and discourse factors. Generic sentences with kind-denoting terms, as in (\ref{ravens.black}) and (\ref{dogs.bark}), express a generalization that holds of a kind 
as well as of its individual members. 
Generic sentences like (\ref{paul.pc.prs}) describe a regularity of action ascribed to an ordinary individual, i.e., a habit in the proper sense of this term. While in \citet{krifka1995genericity} or \citet{carlson2011}, for instance, they constitute a subclass of generic sentences, some take them to be less typical or marginal members of generic sentences (see e.g. \citealt{nickel2017}).

%\textsc{kind reference} and \textsc{generic sentences} are independent of each other. There are generic sentences without kind reference like (\ref{paul.pc.prs}),and kind reference realized in episodic sentences, as in \textit{Marconi invented the radio, The horse arrived in the New World around 1500}.Kind reference also occurs in generic sentences, as we have just seen in (\ref{ravens.black}) and (\ref{dogs.bark}).
%\footnote{Such examples of kind-reference in episodic sentences  are referred to having an `avant-garde' reading, see e.g., \citet{krifka1995genericity,carlson2006meaningful}.}.

All agree that \textsc{generic sentences} are fundamentally opposed to \textsc{episodic sentences}. 
The latter refer to particular singular situations, as in (\ref{Fido.prog}), 
or to a closed set of plural of situations, as described in 
the iterative sentence (\ref{paul.iter}):

\ea {\label{Fido.prog} Fido is barking at the mailman.} \jambox*{\textsc{episodic} sentences}  
\ex {\label{paul.iter} Paul played  three games of chess last Saturday.}  
\z

%Episodic sentences describe particular situations, conditions or facts, which may serve as the base for the truth conditions of related generic sentences, e.g., witnessing what (\ref{Fido.prog}) and (\ref{paul.iter}) describe is part of inductive evidence for the truth of (\ref{dogs.bark}) and (\ref{paul.pc.prs}), respectively.
 
\noindent
However, there is no general agreement on the positive criteria that delimit all and only generic sentences (see e.g., \citealt{Dahl1985,dahl1995,nickel2008,nickel2016,
pelletier-equal2009,pelletier-intro2009,carlson2013}). There is some agreement that prototypical examples are ``bare'' generics like \textit{Birds fly} 
that involve kind reference in a generic sentence
and have no 
restrictors like Q-adverbs (e.g., \textit{always, generally, usually}), 
overt determiner quantifiers (e.g., \textit{some, many, most, all}), 
modal expressions, or qualifiers that restrict the relevant cases based on which the generalization is made.

Moreover, at least since \citet{carlson1977unified}, 
there has been some agreement
that generic sentences, in contrast to episodic sentences, have three main properties: 

\eanoraggedright \textsc{Properties of generic sentences}
\eanoraggedright All are aspectually stative.
\ex All have an intensional import, which is tied to their predictive\linebreak force, with implications for other accessible worlds than just the (actual) world of evaluation, and perhaps all possible worlds.
\ex Most admit exceptions, but may be compatible with no exceptions. 
\z 
\z 

\noindent
All generic sentences are uniformly stative, 
be they headed by (i) generic predicates  that are derived from stage-level predicates like \textit{bark}, as in \textit{Dogs bark} (on its generic interpretation),
or by (ii) individual-level predicates like \textit{(be) intelligent, know (French)}, which are lexically stative predicates (e.g.,  \citealt{chierchia1995}). One standard test for stativity is the lack of reference to particular situations which are locatable 
at specific spatio-temporal locations. 
For instance, this is evident in the incompatibility or oddity with adverbials referring to specific time-points and locations:
\ea \ea [?] {\label{stative1}Dogs bark right now.} 
\ex[?] {\label{stative2}John was intelligent on his porch at 4 p.m.}
\z
\z
\noindent 
The reason for the oddity of (\ref{stative1}) 
and (\ref{stative2}) is that they involve characterizing properties of individuals that hold of them during their whole life-time or at relatively long intervals as well as at any moments of such intervals. Therefore, it is pragmatically odd to assert that they hold only at one specific moment of time or location.


Second, all generic sentences have an intensional (modal) meaning component
(first observed by \citealt{lawler1973,dahl1975}),
which follows from the nature of generalizations. 
They transcend our immediate experiences of 
the world, particular situations, 
states of affairs, or facts, in so far as 
they specify what is possible, likely, i.e., they have a predictive, law-like force. This may be based on what 
actually obtains at given worlds and times, and 
we might have observed. 
However, some generics only 
concern what is \textit{merely} 
possible and may even never ever be realized, as in 
(\ref{speaker}) and (\ref{glass}). They provide the most compelling evidence for intensionality of generics:
\ea
\ea \label{speaker} The Speaker of the House succeeds the Vice President.\jambox*{\citep{carlson1995truth}}
(Context: The rule of the United States Constitution that has never been enacted.)
\ex \label{glass} This knife cuts cheese.\\
(Context: The knife has just been delivered in the mail, but its blade got irreparably damaged in shipping.)
\z
\z
%Generics like \textit{John drinks beer}, `habituals' in the proper sense of the word describe regularities of action that have been instantiated with some regularity (weak, descriptive generalizations). The fact that they actualized realizations does \textit{not} mean that they lack an intensional meaning component, as some seem to assume.  In fact, they do. \textit{John drinks beer} may be understood as meaning that John has a disposition to drink beer, and supports the counterfactual inference that if John were offered beer, he would most likely or probably drink it. 

%Intensionality sets generic, and also habitual, sentences clearly apart  from \textsc{extensional} quantified sentences. For instance, quantified sentences like \textit{All/Most dogs bark} on their extensional interpretation make a quantitative claim  about a closed set of existing dogs in some domain of quantification and relative to a particular possible world. In contrast, for the truth of \textit{Dogs bark} (on its generic interpretation) the actual extension of its subject at a particular possible world  may not matter at all (\citealt[p.~1132-3]{pelletier1997generics}),  because it makes a claim about an open set of (realistically) possible dogs also in other accessible worlds, and perhaps all possible worlds (\citealt{carlson1995truth,pelletier1997generics}).  \textit{Dogs bark} on its generic interpretation may be true even if at a particular world of evaluation no dogs happen to bark. \textit{Dogs bark} is true based not just on inductive evidence (e.g., having witnessed barking dogs on a number of occasions), but also on evidence which may not be directly accessible to us, like some underlying causes in the world (such as regularities governing the biology of living organisms) which ultimately ground observable behavior like dog barking. 

The intensional component of generic sentences may be grounded in some underlying causes, principles or factors, or some kind of law that we may surmise behind the pattern, which we may know, but also about which we may be agnostic or ignorant. For instance,
\textit{Dogs bark} is true based not just on inductive evidence (e.g., having witnessed barking dogs on a number of occasions), but also on evidence which may not be directly accessible to us, like some underlying causes in the world (such as regularities governing the biology of living organisms) which ultimately ground observable behavior like dog barking. 
%Most importantly, the kind of intensionality that generic sentences exhibit cannot be reduced to the modality (intensionality) independently proposed for the analysis of progressive aspect (e.g, pace \citealt{ferreira2004, ferreira2016, hacquardDISS}, i.a.), and neither does it fit deontic and epistemic modality, nor other known modal notions like necessity or essence  (\citealt{nickel2013,nickel2017}, and elsewhere). 

The third key property of generic sentences is their exception-tolerance. 
There are generics that admit of no exceptions in all the known conditions, e.g., universal laws of nature 
like \textit{The Earth revolves around the Sun} 
or \textit{Water consists of two hydrogen atoms and one oxygen atom}. However,
most tolerate exceptions, but may also be compatible with no exceptions whatsoever.
For instance, we judge \textit{Dogs bark} to be true,
because most dogs bark, even if, as is generally known, 
there are non-barking dog species like Basenjis and Greyhounds, or individual 
dogs that do not bark due to injury or illness.
These are exceptions that we tolerate.
As an extensional consequence, it follows that
we cannot validly infer for any particular, arbitrary individual dog that it will have the property of barking. Put differently, that dogs bark is a defeasible fact about dogs. 
%From the fact that Fido is a dog, I may infer that Fido can bark, but I may retract this conclusion at a later time if I learn that Fido is a Greyhound. Therefore, for the representation of generic and habitual sentences non-monotonic logics and generally non-monotonic reasoning approaches seem useful, as they focus on exception-tolerating generalizations and reasoning with exceptions.
This makes generic sentences particularly amenable to 
analyses in terms of non-monotonic logics. 


What makes any analysis of exception-tolerance
so intriguing and challenging is 
that different types of generic sentences admit a different number and kind of exceptions 
or non-confirming cases, while still remaining true. 
Some we judge to be true based on majority satisfaction, 
such as \textit{Dogs bark}.  However, majority 
satisfaction is neither sufficient nor necessary for the truth of all exception-tolerating generics. It is not sufficient, because, for instance, (\ref{printed.books})
is false, despite the majority of printed books being paperbacks. Majority satisfaction is not necessary, because 
what is characteristic of a kind need not be prevalent among its members. (\ref{lions.mane}) holds of only a minority subkind, male adult lions.
What is characteristic of a kind may also be characteristic of only
a minute minority, as 
(\ref{mosquitoesEnglish}) shows. It is judged true, even if a 
vast majority of mosquitoes (\qty{99}{\percent}) fail to carry the virus.

\ea
\ea
 \label{printed.books} Printed books are paperbacks. \hfill \textsc{false}
\ex {\label{lions.mane} Lions have a mane.} \hfill \textsc{true}
\ex {\label{mosquitoesEnglish}Mosquitoes carry the West Nile virus (though \qty{99}{\percent} do not).} \hfill \textsc{true}\\ 
\z
\z

\noindent How many and what kind of exceptions a given 
generic sentence will admit  depends on a variety of
factors which may include not only objective facts  of the world, but also our 
kind-specific expectations, which may be influenced
by our subjective judgments about what we 
evaluate as significant, salient or striking in some way for that kind,
like a lion's mane or a virus that is highly dangerous to us.
%We judge (\ref{lions.mane}) to be true, because a lion's mane isparticularly salient, noticeable, and striking so that we associate it with our kind-specific expectations about lions in general.The West Nile virus is for us potentially highly dangerous, which makes it highly significant, striking, in short subjectively impressive (see \citealt{krifka1995genericity}, \citealt{leslie2007generics,leslie2008}). Moreover, given that we cannot validly infer for any arbitrary mosquito whether it carries the virus, it may be thought as advantageous and prudent for us to treat it as a characterizing property of the whole kind. While subjective, affective-evaluative factors that interact with the truth (and falsity) of generic sentences have been discussed (e.g., \citealt{krifka1995genericity, leslie2007generics,leslie2008, Lazaridou-Chatzigoga2019}), they largely remain at a pretheoretical level, and it  is unclear how they can be precisely analyzed and formally represented (see \citealt{nickel2017}).
Compatibility of generics with a different number and 
kind of exceptions raise the 
fundamental question about what we
view as ``characteristic'', ``typical'' or ``normal'', on the one hand, and what we view 
as failing to conform to a given generic statement and why, on the other hand.

The above examples and observations should suffice to illustrate
that intensionality and exception-tolerance foil any attempt to analyze 
all generic sentences in 
terms of a \textit{single} extensional quantifier (over episodic formulas containing free variables)
(e.g., \textit{all, most} or \textit{some}), 
a \textit{single} quantity expression (e.g., \textit{in a majority of cases} or 
\textit{in a significant number of cases}), quantity-based
criterion, or statistical criterion/correlation, 
no matter how contextually constrained and/or modalized they might be, or possibly also 
defined in vague or probabilistic terms
(\citealt{pelletier1997generics,carlson2013,nickel2013,nickel2017}, i.a.). 

A number of promising analyses have been proposed for exception-tolerance of generics and related issues, while intensionality has garnered much less attention (see summaries in \citealt{krifka1995genericity,pelletier1997generics,MariBayssadeDelPrete2013genericity,nickel2017}, i.a.). None of the proposed analyses has offered a satisfactory comprehensive semantic 
framework for  
all generic sentences, and none so far has emerged as widely accepted. 

Therefore, when it comes to representational assumptions, as a useful point of reference
and an organizing principle for examining meanings of generic 
(and habitual) sentences here, 
I will at various points rely on a logical representation in terms of a tripartite structure with a phonologically null, modalized dyadic operator \cnst{Gen}, akin to overt Q-adverbs, used in quantificational theory of genericity. In its original formulation,
it has the following
general form (adapted from \citealt[32, ex.~(57)]{krifka1995genericity}, which in turn relies on suggestions in \citealt{krifka1987nominal}):\footnote{\cnst{Gen} and overt Q-adverbs are not unselective binders, in contrast to their early treatments, 
as in \citet{carlpell1995generic}, for instance. Which variables can be bound in a given logical structure depends on a number of factors, including context and topic-focus structure.}

\ea $\cnst{Gen} [x_1\, \dots\, x_i ; y_1\, \dots\, y_j](\cnst{Restrictor} [x_1\, \dots\, x_i]; \cnst{Matrix} [\{x_1\}\, \dots\, \{x_i\}, y_1\, \dots\, y_j])$\\\label{tripartite}	
$x_1\, \dots\, x_i$: variables bound by \cnst{Gen}\\	$y_1\, ...\, y_j$: variables bound existentially, with scope just in matrix \\		 
$\{x_1\}\, \dots\, \{x_i\}$: 	means $x_1 \dots x_i$ may or may not occur in matrix 
\z

\noindent
In the tripartite structure (\ref{tripartite}),
just like an overt Q-adverb, \cnst{Gen} relates two arguments, a Restrictor and a Matrix  which are separated by `;'.  \cnst{Gen} stands for a certain characterizing relation between the conditions stated in the Restrictor (which states the restricting cases) and 
Matrix (which makes the main assertion of the sentence):  intuitively, it sets conditions on the number or proportion of individuals
that may satisfy the formula, and if the conditions stated in the Restrictor hold, other conditions stated in the Matrix  will also hold.


We may distinguish three main types of instances or particulars 
that can serve as the base for the generalization (\citealt{carlson2011}): (i) instances of particular \textsc{individuals} (as in \textit{Ravens are black}), 
%where it is an instance of a particular raven that is in the state of being black), 
(ii) instances of particular \textsc{situations} involving 
(stages of) a single individual of an ordinary sort (as in \textit{Paul plays chess}), or (iii) instances of particular \textsc{situations} involving instances of different \textsc{individuals}
that are members of a given kind (as in \textit{Dogs bark}).\footnote{The notion of 
\textsc{situations} corresponds to occasions or cases, as \citet{lawler1972} originally proposed, and are akin to ``cases'' of \citet{lewis1975} \citep[30]{krifka1995genericity}.
It plays a role similar to that of \textsc{stages} in \citeauthor{carlson1977unified} (\citeyear{carlson1977unified}, and elsewhere) which is tied to Carlson’s distinction between stage-level (episodic) predicates (SLPs) and individual-level (stative) predicates (ILPs).}$^,$\footnote{\label{fn:double}Generics that fall under (iii)  involve a ``double generalization'' (in \citealt{carlson2011}; see also \citealt[1131]{pelletier1997generics}): 
\textit{Dogs bark} is true by virtue of (a) \textit{individual dogs} that realize the kind \textsc{dog}  having the characterizing property of barking, and also by virtue of (b) particular \textit{situations} of barking by a stage of an individual dog.} 

%It is important to emphasize that the generic morpheme in Slavic languages can be used in sentences that express generalizations over any of the three cases (i)-(iii) as their base. 
The cases (ii) and (iii) involve quantification over situations, which in \citet{krifka1995genericity}
is a defining feature of what they refer to as \textsc{habitual sentences}, a subcase of \textsc{generic sentences}:

\ea \label{habitual.def} A sentence is \textsc{habitual} iff its semantic representation is of the form\\
$\cnst{Gen}[ \ldots  s \ldots ; \ldots ] (\cnst{Restrictor} [ \ldots s \ldots ]; \cnst{Matrix} [ \ldots  s  \ldots ])$, \\
where $s$ is a situation variable\hfill\citep[32, ex.~(56)]{krifka1995genericity} 
\z
%It is important to observe that this is one sense in which `habitual(ity)' is used, and that `habitual(ity)' is not consistently used for the same phenomena in the relevant literature,  as already observed above.Henceforth, I will use the terms `habitual' and `a habitual sentence' in the sense of (\ref{habitual.def}), i.e., habituality, defined in terms of quantification over situations, will be taken as a special case of (characterizing) genericity. 
% $\cnst{Gen} [s,x;]  (x = \textrm{Paul} \wedge x\, "\textrm{in}\, s \wedge  \textsc{on Saturday}(s); \textsc{plays chess}(x,s))$\\

\newpage
\noindent Here,  I reserve the term \textit{habitual sentence} for only generic sentences that express habits in the proper sense of the word,
as also observed above (\sectref{fil:sec:introduction}). Most importantly, 
and as we will see below in \sectref{semprag}, 
the generic morpheme \textit{-va-} is not constrained to conveying habits proper, i.e., regularities of action of animate, mostly, human agents (the case (ii) above).  
Neither does it require that its domain contains situations
(the cases (ii) and (iii) above), but it may quantify
over individuals only (the case (i) above). In short, 
it functions as a quantifier over all the three main types of 
particulars mentioned above, and in this respect it
is akin to \cnst{Gen}. 

For any theory of genericity
the main burden of the analysis is on specifying the relation between the generalization and
what it is based on, its basis for generalization, 
which in many cases (but not all) are some particulars, instances, and particular facts. 
How this relation comes about is at the core of 
what it means for a sentence to express a generalization,
rule, pattern and the like. We have a rich variety of competing proposals regarding this relation, and by the same token regarding the truth conditions of generic sentences that underpin \cnst{Gen}. One useful way of sorting them out was proposed by \citet{carlson1995truth}, who suggests that they can be divided into two main perspectives, and mixed proposals combining some elements from both
can also be found.
One perspective is characterized in terms of an \textsc{inductive model}, on which 
generic sentences express inductive generalizations that are based on some observed (or unobserved) set of instances or particulars in the actual world from which we infer a general rule. Paradigm cases are sentences like \textit{Paul plays chess} that express generalizations over situations of the kind described by their corresponding episodic sentences like \textit{Paul was playing chess when I called him yesterday, Today during the lunch break, Paul played chess with Erwin, Last week, Paul played chess with grandpa and won each and every time}.  On this perspective, generics would (mainly) express descriptive (weak) generalizations. 
%One disadvantage of the inductive model is that it inherits all  the problems associated with induction, of course. 

The second main perspective corresponds to what \citet{carlson1995truth} dubs a \textsc{rules-and-regulations model}. On this perspective, generics express generalizations that do not (necessarily) depend on particular episodic conditions in the actual world, but rather on some causal structure or underlying forces in the world, which give rise to episodic instances (if there are any), and motivate the patterns that we recognize in the actual world. 
Paradigm cases are generic sentences expressing various rules, such as game rules that we may learn directly, rather than inferring them from episodic conditions in the actual world,
for instance, by reading
the relevant rules, manuals or guidelines, such as the rules of chess like \textit{Bishops move diagonally}, or following instructions, customs and the like.
 
In contrast to the inductive model, the rules-and-regulations model can
also handle generic sentences  that are independent of particular instances or situations, because they have no actualized verifying instances, such as (\ref{speaker}).
%, or like \textit{Sam is a bachelor}, for which it seems impossible, or perhaps nonsensical, to try to identify once and for all what counts as  a verifying instance in the actual world. 
On the rules-and-regulations model of genericity, we judge generic sentences to be true or false with respect to a set of rules (or a finite list of propositions), which are considered to be irreducible entities, supporting the truth of possibly an unlimited number of propositions. 

\citet{carlson1995truth} proposes that the rules-and-regulations model is a better candidate than the inductive one 
for a unified semantic analysis of generic sentences, and a model
which might ground the semantics of \cnst{Gen}.
However, it is not a good fit for generics
like \textit{Paul plays chess}, ``which lend the most prima facie plausibility to the inductive model'' \citep[237]{carlson1995truth}.
A rules-and-regulations model requires us to identify some rule, causal factor or underlying condition that gives rise to or ``regulates'' Paul's chess playing behavior, which seems rather far-fetched. 

These two main perspectives raise an intriguing question which of them might best fit generics marked with the generic morpheme \textit{-va-} in Czech, and in generic sentences that are marked with the cognate marker in other Slavic languages.  Given that the the marker \textit{-va-} requires that there be realized instantiations that count as evidence for the truth of generic sentences marked by it, as I will show, generic sentences
marked with \textit{-va-} at first blush fit the inductive model of genericity.
At the same time, as it will be shown, \textit{-va-} can also be used in generics that express rules, regulations, often with the goal of ensuring plausible deniability, or safeguarding a rule, a regulation against refutation by states of affairs 
of which its author is ignorant. 
The exploration of the properties of the generic morpheme \textit{-va-} is highly relevant to the answer to the questions posed at the outset: What kind of generic sentence does the generic morpheme \textit{-va-} delimit? Can we provide a uniform analysis for all generic sentences?

\section{Morphological background} \label{basicmorphology}

\subsection{The perfective/imperfective and the generic/non-generic distinctions} 
\largerpage[-1]
The generic marker \textit{-va-}, here glossed as \textsc{gen}, is productive in West Slavic languages: namely, in Czech, Slovak (e.g., \citealt{kucera1981aspect,bily1986iterative}, i.a.), and to a lesser degree also in Polish (\citealt{bily1986iterative}, i.a.). 
In South and East Slavic languages, its occurrences are also
attested, but they may be treated as a part of a lexicalized combination with a verb base.\footnote{In Russian, 
verbs with this marker are taken to be a feature of a colloquial (``substandard'') speech (\citealt[405--407]{isacenko1962}, \citealt[28 and 168--171]{Forsyth1970}, \citealt[27]{Comrie1976},  \citealt[177]{kucera1981aspect}, \citealt{bily1985preliminaries}, \citealt[413--414]{vinogradov1986russian}, \citealt{petr1986mluvnice}, \citealt{bily1986iterative}),
while some of them are in fact quite frequent: e.g., \textit{byvat’} (< \textit{byt’} `to be'), \textit{znavat’} (< \textit{znat’}  `to know'), \textit{govarivat'} (< \textit{govorit'} `to speak'), \textit{pivat'}
(< \textit{pit'} `to drink'), \textit{siživat'} (< \textit{sidet'} `to sit') and \textit{xaživat'} (< \textit{xodit'} `to go'). Interestingly, the generic marker similar to that in West Slavic languages is productive in some Northern Russian dialects \citep{barnetova1956nasobenost}.}

\begin{figure}[t]
%FK 19 Aug: I entered a forest version of the tree and put everything in a Figure environment. I also sent a e-mail to Radek and Berit about it.
%\includegraphics[width=13cm]{5VerbForms.jpg}\\
\begin{forest}
[Slavic verb forms 
    [perfective (pfv)\\non-generic,font=\scshape
        [primary \textsc{pfv}\\(a),name=a] 
  		[prefixed \textsc{pfv}\\(b),name=b]
	]
	[imperfective (ipfv),font=\scshape,calign=child,calign child=1
        [\textsc{non-generic}\\“general imperfective”\\
             \citep{Comrie1976}
    		 [primary \textsc{ipfv}\\(c),name=c]
    		 [secondary \textsc{ipfv}\\marked with\\the \textsc{ipfv} suffix\\
    		  (d),tier=lowest,name=d
    		 ]
        ]
		[\textsc{generic}\\mainly West Slavic,edge=dashed
			[marked with\\the generic -\textit{va}-\\(e),tier=lowest,name=e]
			% no edge?
		]
	]
]
\begin{scope}[>={Triangle[]}]	
\draw[->] (a.south) -- ++(0pt,-7\baselineskip) -| (d.270); % 270 = south
\draw[->] (b.255) -- ++(0pt,-6.5\baselineskip) -| (d.260);
\draw[->, very thick] (d.280) -- ++(0pt,-1.5\baselineskip) -| (e.260); 
\draw[->, very thick] (c.280) -- ++(0pt,-4\baselineskip) -| (e.280);
\draw[->] (c.260) -- ++(0pt,-1.5\baselineskip) -| (b.280);
\end{scope}
\end{forest}
    \caption{Common formation patterns of non-generic verb forms (perfective and imperfective) and generic verb forms (following e.g.  \citealt[180]{petr1986mluvnice}). The dashed line: The traditional subsumption of generic (also known as ``iterative'') verbs under imperfective aspect is not uncontroversial (see e.g. \citealt{kopecny1948}, i.a.) and may best be refuted.} 
    \label{fig:filip:1}
\end{figure}

Taking Czech as our case in point, 
in traditional reference grammars, the verb system is structured by two main 
distinctions: (i) the perfective/imperfective distinction, and (ii) the generic/non-generic distinction.
%In order to explain the place of generic forms built with the generic morpheme in Slavic verb systems, 
This is sketched in \figref{fig:filip:1}, 
which builds on the rudimentary schema given in  \citet[180]{petr1986mluvnice}. \tabref{filip:tab:examples}
contains examples illustrating the different types of verb forms and their morphological relations. 

In Bohemistics, the morpheme \textit{-va-} (here glossed as \textsc{gen}) is traditionally labeled ``iterative'' (\textit{iterativní} in Czech) or ``multiplicative'' (\textit{násobená} in Czech). It is clearly separated from the imperfectivizing suffix (see e.g., \citealt{petr1986mluvnice}, \citealt{lamprecht1986historicka}, \citealt{Kosek2014}, \citealt{biskup2024}; \citealt{Genis2021} also separate the cognate generic morpheme from the imperfectivizing suffix in other Slavic languages).  However, at the same time,
this morpheme is subsumed under imperfective aspect, treated as a kind of imperfective marker, with one notable exception, namely \citet{kopecny1948,kopecny1966}.

Notice that \figref{fig:filip:1} has a dashed line connecting the \textsc{imperfective (ipfv)}  with the \textsc{generic} node. 
It indicates that the traditional subsumption of verbs with 
generic morpheme \textit{-va-} under imperfective 
aspect is not uncontroversial.  
It is one of the main goals of this chapter to show 
that verbs marked with the generic morpheme \textit{-va-} do not neatly fit the imperfective paradigm in Slavic languages, and imperfectivity in natural language generally.
They have a number of semantic and pragmatic properties that differ from both primary and secondary imperfectives in Slavic languages. This in turn conforms with independent arguments made elsewhere that genericity is is best viewed as independent of imperfectivity, and other categories of the TMA system (see e.g., \citealt{filip1997suigenva}).

\begin{table}
\caption{Examples illustrating the morphological relations included in the schema in \figref{fig:filip:1}}
\label{filip:tab:examples}
\begin{tabularx}{\textwidth}{l@{~~}X@{}cl@{~~}X}
\lsptoprule
    (a)&dát\textsuperscript{PFV primary}&$\rightarrow$&(d)&dávat\textsuperscript{IPFV secondary}\\
    &give.\textsc{inf}&&&give.\textsc{ipfv.inf}\\
    &`to give'&&&`to give'/`to be giving'\\[2ex]
    (d)&dávat\textsuperscript{IPFV secondary}&$\rightarrow$&(e)&dávávat\textsuperscript{generic}\\
    &give.\textsc{ipfv.inf}&&&give.\textsc{ipfv}.\textsc{gen}.\textsc{inf}\\
    &`to give'/`to be giving'&&&`to give regularly, seldom,\\
    &&&&often, usually, \dots\\[2ex]
    (c)&štěkat\textsuperscript{IPFV primary}&$\rightarrow$&(b)&poštěkat\textsuperscript{PFV prefixed}\\
    &bark.\textsc{inf}&&&\textsc{pref}.bark.\textsc{inf}\\
    &`to (be) bark(ing)'&&&`to bark (a little)';\\
    &&&&metaphorically:\\
    &&&&`to bicker', `to squabble'\\[2ex]
    (b)&poštěkat\textsuperscript{PFV prefixed}&$\rightarrow$&(d)&poštěkávat\textsuperscript{IPFV secondary}\\
    &\textsc{pref}.bark.\textsc{inf}&&&\textsc{pref}.bark.\textsc{ipfv}.\textsc{inf}\\
    &`to bark (a little)';&&&`to (be) bark(ing) (a little)';\\
    &metaphorically: `to bicker', \dots&&&metaph.:`to (be) bicker(ing)', \dots\\[2ex]
    (d)&poštěkávat\textsuperscript{IPFV secondary}&$\rightarrow$&(e)&poštěkávávat\textsuperscript{generic}\\
    &\textsc{pref}.bark.\textsc{ipfv}.\textsc{inf}&&&bark.\textsc{ipfv}.\textsc{gen}.\textsc{inf}\\
    &`to (be) bark(ing) (a little)';&&&`to (tend to) bark/bicker a little,\\
    &`to (be) bicker(ing)', \dots&&&regularly, seldom, often,\\
    &&&&usually, \dots\\[2ex]
    (c)&štěkat\textsuperscript{IPFV primary}&$\rightarrow$&(e)&štěkávat\textsuperscript{generic}\\
    &bark.\textsc{inf}&&&bark.\textsc{gen}.\textsc{inf}\\
    &`to (be) bark(ing)'&&&`to (tend to) bark regularly,\\
    &&&&seldom, often, usually, \dots\\
    \lspbottomrule
\end{tabularx}
\end{table}

%% THE CODE BELOW IS REDONE BY THE TABULARX ABOVE
% \begin{multicols}{2}
% \exdisplay
%     \begingl
%     \gla(a)\hspace{0.2cm}dát\textsuperscript{PFV primary}  \hspace{1.8cm}  $\rightarrow$ //
%     \glb\hspace{0.65cm}give.\textsc{inf} //
%     \glft\hspace{0.65cm}`to give'//
%     \endgl
%     \xe
%    \exdisplay
%     \begingl
%     \gla(d)\hspace{0.2cm}dávat\textsuperscript{IPFV secondary} //
%     \glb\hspace{0.65cm}give.\textsc{ipfv.inf} //
%     \glft\hspace{0.65cm}`to give'/`to be giving'//
%     \endgl
%     \xe  
% \end{multicols}
% \vspace{-0.5cm}
% \begin{multicols}{2}
% \exdisplay
%     \begingl
%       \gla(d)\hspace{0.2cm}dávat\textsuperscript{IPFV secondary}  \hspace{1.8cm}  $\rightarrow$ //
%     \glb\hspace{0.65cm}give.\textsc{ipfv.inf} //
%     \glft\hspace{0.65cm}`to give'/`to be giving'//
%     \endgl
%     \xe
%    \exdisplay
%     \begingl
%       \gla(e)\hspace{0.2cm}dávávat\textsuperscript{generic} // 
%      \glb\hspace{0.65cm}give.\textsc{ipfv}.\textsc{gen}.\textsc{inf} //
%     \glft\hspace{0.65cm}`to give regularly, seldom, //
%     \hspace{0.65cm}often, usually, ...
%     \endgl
%     \xe  
% \end{multicols}
% \vspace{-0.5cm}
% \begin{multicols}{2}
% \exdisplay
%     \begingl
%     \gla(c)\hspace{0.2cm}štěkat\textsuperscript{IPFV primary}  \hspace{1.8cm}  $\rightarrow$ //
%     \glb\hspace{0.65cm}bark.\textsc{inf} //
%     \glft\hspace{0.65cm}`to (be) bark(ing)' //
%     \endgl
%     \xe
%    \exdisplay
%     \begingl
%     \gla(b)\hspace{0.2cm}po.štěkat\textsuperscript{PFV prefixed} //
%     \glb\hspace{0.65cm}\textsc{pref}.bark.\textsc{inf} //
%     \glft\hspace{0.65cm}`to bark (a little)'; metaphorically: //
%     \hspace{0.65cm}`to bicker', `to squabble'
%     \endgl
%     \xe  
% \end{multicols}
% \vspace{-0.5cm}
% \begin{multicols}{2}
% \exdisplay
%     \begingl
%     \gla(b)\hspace{0.2cm}po.štěkat\textsuperscript{PFV prefixed}  \hspace{1.5cm}  $\rightarrow$ //
%     \glb\hspace{0.65cm}\textsc{pref}.bark.\textsc{inf} //
%     \glft\hspace{0.65cm}`to bark (a little)';//
%     \hspace{0.65cm}metaphorically: `to bicker', ...
%     \endgl
%     \xe
%     \exdisplay
%     \begingl
%     \gla(d)\hspace{0.2cm}po.štěkávat\textsuperscript{IPFV secondary} //
%     \glb\hspace{0.65cm}\textsc{pref}.bark.\textsc{ipfv}.\textsc{inf} //
%     \glft\hspace{0.65cm}`to (be) bark(ing) (a little)'; //
%     \hspace{0.65cm}metaph.:`to (be) bicker(ing)', ... 
%     \endgl
%     \xe   
% \end{multicols}
% \vspace{-0.9cm}
% \begin{multicols}{2}
% \setlength{\parindent}{0pt}
% \exdisplay
%     \begingl
%     \gla(d)\hspace{0.2cm}po.štěkávat\textsuperscript{IPFV secondary}  \hspace{0.9cm}  $\rightarrow$ //
%     \glb\hspace{0.65cm}\textsc{pref}.bark.\textsc{ipfv}.\textsc{inf} //
%     \glft\hspace{0.65cm}`to (be) bark(ing) (a little)'; //
%     \hspace{0.65cm}metaph.:`to (be) bicker(ing)', ... 
%     \endgl
%     \xe
%     \exdisplay
%     \begingl
%     \gla(e)\hspace{0.1cm}po.štěkávávat\textsuperscript{generic} //
%     \glb\hspace{0.65cm}bark.\textsc{ipfv}.\textsc{gen}.\textsc{inf} //
%     \glft\hspace{0.65cm}`to (tend to) bark/bicker a little, //
%     \hspace{0.65cm}regularly, seldom, often, usually, ...
%     \endgl
%     \xe
% \end{multicols}
% \vspace{-0.9cm}
% \begin{multicols}{2}
% \exdisplay
%     \begingl
%     \gla(c)\hspace{0.2cm}štěkat\textsuperscript{IPFV primary}  \hspace{2cm}  $\rightarrow$ //
%     \glb\hspace{0.65cm}bark.\textsc{inf} //
%     \glft\hspace{0.65cm}`to (be) bark(ing)'//
%     \endgl
%     \xe
%     \exdisplay
%     \begingl
%     \gla(e)\hspace{0.2cm}štěkávat\textsuperscript{generic}//
%     \glb\hspace{0.65cm}bark.\textsc{gen}.\textsc{inf} //
%     \glft\hspace{0.65cm}`to (tend to) bark regularly, seldom,//
%     \hspace{0.65cm}often, usually, ...
%     \endgl
%     \xe
% \end{multicols}

In Slavic languages, all the forms without the generic morpheme \textit{-va-} that are inherently episodic alternate between episodic and generic readings depending on context.
Hence, they are referred to as non-generic in 
Figure \ref{fig:filip:1}.
Slavic non-generic forms are either imperfective or perfective, and members of both classes can be either simple (root, primary) or derived, mostly by means of derivational affixation. Primary and secondary imperfectives
correspond to ``general imperfectives'' \citep{Comrie1976}, meaning that they
cover the whole semantic domain of imperfectivity,
both its episodic and generic subdomains. 

Perfective verbs in Slavic languages, all of which are non-generic (i.e., have either episodic or generic uses depending on context), are for the most part derived. Prefixation operating on primary imperfective bases is the most common way of forming perfective verbs. This corresponds to the derivational relation (c) $\rightarrow$ (b) in Figure \ref{fig:filip:1}.
Slavic verb prefixes 
have hallmark properties of lexical-derivational, rather than inflectional, morphology, and hence are not grammatical markers of perfectivity akin to bona fide inflectional morphology like tense markers (\textit{pace} \citealt{borer2005}, i.a., for arguments see \citealt{filip1992,filip2000quantization, Filip2026a}, and elsewhere).\footnote{The argument that Slavic verb prefixes are not grammatical markers of perfective aspect in Slavic languages provides a key support for the view that the Slavic perfective/imperfective distinction is predominantly of lexical-derivational, rather than inflectional, nature. This is also reflected in the common Slavic lexicographic practice of marking lexical entries of verbs as either perfective or imperfective, and some as bi-aspectual.} 
Slavic prefixation derives new verbs, typically with a different lexical meaning and argument structure than their derivational base. In the relevant examples in \tabref{filip:tab:examples}, the prefix \textit{po-} not only derives a perfective verb from the primary imperfective base, under (c) in Figure \ref{fig:filip:1}, but also contributes
an attenuative meaning to the resulting prefixed verb 
which roughly amounts to `bark (a little)',
and this prefixed verb also has a metaphoric sense of `to bicker', `to squabble'. Prefixes are also applied to already prefixed perfective bases, so that more than one prefix can occur on one and the same verb. 

Secondary imperfectives (under (d) in Figure \ref{fig:filip:1}) are derived by means of the imperfectivizing suffix, glossed as \textsc{ipfv}. 
In all Slavic languages, the imperfectivizing suffix applies to perfective bases and derives their imperfective counterparts, typically with the
same lexical meaning and argument structure. The perfective
bases to which the imperfectivizing suffix is applied is
either primary perfective, as in (a) $\rightarrow$ (d), or prefixed, as in (b) $\rightarrow$ (d).\footnote{The imperfectivizing suffix may generate so-called ``aspectual pairs'', i.e., two verbs with the same lexical meaning that only differ in grammatical perfective/imperfective aspect. Therefore, some classify it as an inflectional suffix (\citealt{Isacenko1960,filip1993phd,Filip1999}).
It is also often taken to be 
the only marker in Slavic languages that is solely dedicated to the encoding of aspect (\citealt{Isacenko1960,filip1993phd,Filip1999}).}

%and are attested in a number of languages (e.g., Romance
%languages, Modern Greek, Georgian, ibid.). 
 
%I.e., they are non-generic, both formally and semantically unmarked for genericity,and alternate between episodic and generic readings depending on context. 

Most importantly, in contrast to the Slavic
imperfectivizing suffix, the generic morpheme applies to an imperfective base realized in verb forms that can have either an episodic or a generic reading and derives a form that is unambiguously generic, i.e., it excludes any episodic interpretation.\footnote{In Czech, the generic morpheme \textit{-va-} applies to past tense imperfective stems \citep[184]{petr1986mluvnice}.} The imperfective
base to which the generic morpheme \textit{-va-} applies
is either primary, as in (c) $\rightarrow$ (e), or secondary, as in (d) $\rightarrow$ (e).
When the generic morpheme is applied to a secondary imperfective base, as in (d) $\rightarrow$ (e), 
the resultant generic verb contains both the imperfectivizing suffix (\textsc{ipfv}) and the generic morpheme (\textsc{gen}).

%So there are two types of general imperfective, non-generic imperfective forms to which the generic morpheme applies: Primary (root) imperfective 
%(formally unmarked for imperfectivity (c)), and secondary imperfective (formally marked for imperfectivity with the \textsc{ipfv} suffix (d)). 
%Some examples are given below the Figure \ref{fig:filip:1}. 
%It is important to emphasize that The secondary imperfective formally marked with \textsc{ipfv}, (d) in Figure \ref{fig:filip:1}, is a general imperfective, non-generic form, and it freely occurs in episodic contexts. In contrast, the generic form (e) in Figure \ref{fig:filip:1} that containsboth \textsc{ipfv} and \textsc{gen} cannot ever be used in episodic contexts. 


%. Further below, I will show that they have a number of properties that sets them apart from Slavic primary and secondary imperfective verbs.

%Continuing with the forms in Figure \ref{fig:filip:1}, secondary imperfectives are formed by means of the imperfectivizing suffix (\textsc{ipfv}).
 

In Czech, the derivation of specifically generic forms from imperfective non-generic ones
is transparent, regular and productive,
and verbs with the generic morpheme \textit{-va-}  are common in all styles of speech (\citealt[177]{kucera1981aspect}, \citealt{petr1986mluvnice}).\footnote{There are also a few irregular generic verbs: e.g., in Czech \textit{číst > čítat} `to read', \textit{vidět > vídat} `to see', \textit{slyšet > slýchat} `to hear', \textit{ležet > léhat} `to lie (down)'.} 
The Czech generic morpheme \textit{-va-} can be reduplicated, which also confirms its productivity: e.g., \textit{mít}\textsuperscript{\textsc{ipfv}} `to have' > \textit{mí-vá{\textasciitilde}vá{\textasciitilde}va-t} [have-\textsc{gen{\textasciitilde}gen{\textasciitilde}gen-inf}] `to tend to have, very sporadically/very often', `to have once in a great while'.
Attested reduplication examples are fairly 
easy to find, and two are given below:

\ea
\ea {\label{redup.battle}
\gll \textit{Říkávává} se, že po bitvě je každý generál.\\
say.\textsc{gen.gen.3sg.prs} \textsc{refl} that after battle is everybody general \\}
\glt `As the saying goes, after the battle, everybody is a general.' 
\ex {\label{redup.custom} 
\gll Bude to z jiného konce, než \textit{bývává} zvykem.\\
will.be it from different end than be.\textsc{gen.gen.3sg.prs} customary\\}
\glt `It will be approached from a different angle than is customary.'
\z
\z
\noindent
Reduplication is often associated with the speaker's subjective judgment concerning the regularity of instances that form the base for the generalization.
It may affectively connote their strikingly high or low frequency.\footnote{(\ref{redup.battle}) retrieved at \url{https://www.vidlakovykydy.cz/clanky/maginotova-linie-ceskoslovenske-opevneni} accessed 13/12/2021; 
(\ref{redup.custom}) from a native Czech speaker (p.c.).}

The formation of generic verbs by means of the morpheme
\textit{-va-}, as just illustrated with some Czech examples,
follows a cross-linguistic pattern that is attested
in a number of typologically unrelated languages that have specifically 
generic forms: namely, generic forms are derived from non-generic ones; they are the marked case relative to non-generic forms, which are morphologically or syntactically simpler and alternate between episodic and generic interpretations. We do not find generic forms that are basic and episodic forms derived from them (\citealt{carlson1995truth,carlson2005}, and elsewhere). 

Examples of languages that have verb markers 
enforcing only a generic reading of sentences, but
are not a necessary condition for the expression of generic statements, can be found in 
\citet{Dahl1985, dahl1995}. The Czech marker \textit{-va-} serves as his example par excellence for this class of markers.
However, \citet{dahl1995} also argues that 
they are not generic markers, but rather habitual markers falling either under
tense and/or aspect; they are habitual markers, because,  
according to him, a bona fide
generic marker must constitute not only a sufficient but also a necessary condition on the expression of all and only generic statements.
Only 3 languages in Dahl's data base of 76 languages have 
such a generic marker: namely, Wolof (Niger-Kordofanian), Iṣẹkiri (Niger-Kordofanian), and Maori (Austronesian) \citep[420,~424]{dahl1995}. Such observations constitute 
one of Dahl's empirical arguments 
in support of his main  conclusion that the episodic/generic distinction has no direct relevance for the grammar of natural language; rather, it is generally only indirectly reflected in speakers' choices between grammatical markers \citep[425]{dahl1995}.

To this it might be immediately objected that Dahl's claim
here relies on an overly restrictive, unwarranted
requirement of what it means for a given marker to count as a grammatical marker of a given category. As is well-known, bona fide grammatical markers (e.g., the English past tense morpheme \textit{-ed}) do not neatly delimit a single semantic domain, and absolute obligatoriness is not a necessary criterion for grammaticalization, as “[s]omething is obligatory relative to the context; i.e., it may be obligatory in one context, optional in another and impossible in a third context” (\citealt[12]{lehmann1995}; see also \citealt{heine2002kuteva,heine2007kuteva}, and others, for similar observations).

\citet[421]{dahl1995} also observes that ``[a] relatively safe generalization is that habituals in general exhibit a low degree of grammaticalization. They are [\dots] possible sources of a grammaticalization path yielding what would be the primary candidates for a gram that systematically marks genericity."  In West Slavic languages at least, and as we have just seen in Czech, the generic marker is productive, and systematically in all its occurrences only contributes a generic interpretation to the meaning of a verb, in a completely transparent and compositional way. This would seem to indicate that it exhibits a high degree of grammaticalization.
%Grams systematically marking genericity, according to \citet{bybee1989dahl}, have their source in forms which at an earlier stage had progressive meaning only, subsequently developed into present or imperfective markings that can be used to express generic statements. 

Having sketched some basic morphological facts, the next section will address the independence of the Slavic generic morpheme from tense and imperfective aspect, and specifically its independence 
from the imperfectivizing suffix. We will also see that it does not neatly fit into the Slavic imperfective paradigm, contrary to standard assumptions in Bohemistics, and also contrary to \citet{dahl1995}.
%Hence, this strongly suggests that it has the status of a gram that systematically marks genericity, independently of imperfectivity.




\subsection  {Independence of the generic morpheme from tense and imperfective aspect} 

The generic morpheme \textit{-va-} is independent of tense. It occurs on non-finite forms 
(including the infinitive and imperative, the relevant examples are given in \citealt{filip1997suigenva}), and with all the 
overt tense markers (past, present, and future) on the same verb or in the same verb
complex, as we see in the following Czech examples:\footnote{In the past tense, generic sentences formally marked with \textit{-va-} 
have a remote past time reference, whereby its reduplication entails a high degree in the past temporal remoteness \citep{kucera1981aspect,kucera1983semantic,bily1986iterative}: e.g., \textit{Vídávával} [see.\textsc{gen.gen}.1\textsc{sg.pst}] 	\textit{jsem 	ho 	tam 	ve čtvrtky} `I used to see him there on Thursdays once in a while / regularly … a very long time ago.' \citep{chmelickova2005}. Generally, 
\textit{-va-} (possibly also reduplicated) combined with the past tense triggers a conversational implicature that the described situation no longer holds, and
can also refer to single extended situations in the past.}


\ea\label{va-tenses}
\ea{\label{va.pasttense} \gll Karel		hrával			hokej.\\ 
Charles	play.\textsc{gen.pst}	hockey	\\} \jambox*{\textsc{past tense}}
\glt `Charles used to play hockey.’ [remote past] 
\ex {\label{va.presenttense} 
\gll Karel		hrává		hokej.	 \\ 
Charles	play.\textsc{gen.prs}		hockey	\\} \jambox*{\textsc{present tense}}	
\glt `Charles tends to play hockey.’ [i.e., regularly, often, seldom,  ... ] 
\ex {\label{va.futuretense} 
\gll Karel		bude			hrávat			hokej. \\	
Charles	be.\textsc{aux.fut}  play.\textsc{gen.inf} 	hockey \\} \jambox*{\textsc{future tense}}  
\glt `Charles will tend to play hockey.’ [i.e., regularly, often, seldom,  ... ]
\z
\z

%Corresponding data are attested in other Slavic  languages, especially in West Slavic, and theyshow that the generic morpheme is independent of tense  (see also \citealt{filip1997suigenva}). 
\noindent
The independence of the generic morpheme \textit{-va-} from tense is not surprising, given that generally
markers of genericity and tense are orthogonal to one another in languages that have overt markers for genericity and tense. 
There are tensed languages with generic markers (e.g., Czech, Polish, Slovak), without generic markers (e.g., English), and tenseless languages with arguably generic markers (e.g., American Sign Language) and without them (e.g., Chinese), in which case they convey generic statements by other means (e.g., in Dyirbal and Burmese by means of a modal distinction between realis and irrealis; see \citealt[51]{comrie1985tense}). 

Also notionally, genericity (and habituality) and tense are independent of each other.  Tense is a deictic category, while genericity (and habituality)
are non-deictic. In tensed languages like English,
generic sentences can be in any tense form, and refer to present, future or past time: e.g., \textit{Dinosaurs ate kelp; Mary walks to school along the river; When he retires, he will spend more time with his family.} In English, the simple (non-progressive) present tense, as opposed to the present progressive, preferably has a generic interpretation.  

Generic sentences are compatible with any sufficiently large temporal interval, contextually located in the present, past, or future. They may also express generalizations that contain overt or covert spatial/temporal restrictors related to the present, past, or future, as shown in \REF{va-tenses2}.
% (e.g., \textit{\textsc{současný} prezident jídává} (\textsc{gen}) \textit{hamburgery} `The current President eats hamburgers', \textit{Dinosauři jídávali} (\textsc{gen}) kelp `Dinosaurs (usually) ate kelp' (the generalization is restricted to the distant past by virtue of general world knowledge), \textit{\textsc{od pondělka} se bude tato kancelář otvírávát} (\textsc{gen}) \textit{až ve dvě hodiny} `Starting next Monday this office will open only at 2 p.m.'

\ea\label{va-tenses2}
\ea\gll \textit{Současný} prezident jídává hamburgery.\\
current president eat.\textsc{gen.prs} hamburgers\\
\glt `The current president eats hamburgers.'
\ex\gll Dinosauři jídávali kelp.\\
dinosaurs eat.\textsc{gen.pst} kelp\\
\glt `Dinosaurs (usually) ate kelp.' [The generalization is restricted to the distant past by virtue of the general world knowledge about the dinosaur kind.]
\ex\gll \textit{Od} \textit{pondělka} se bude tato kancelář otvírávat až ve 2 hodiny.\\
from Monday \textsc{refl} be.\textsc{aux.fut} this office open.\textsc{gen.ipfv} only at 2 o'clock\\
\glt `Starting next Monday this office will open only at 2 p.m.'
\z\z


Let me now turn to the independence of the generic marker from the imperfectivizing suffix. 
%, notice that \figref{fig:filip:1} has a dashed line connecting the \textsc{imperfective (ipfv)}  with the \textsc{generic} node. \figref{fig:filip:1} is an adaptation of the schema given in \citet[180]{petr1986mluvnice} and reflects the traditional view on which the morpheme that is here glossed as \textsc{gen} and traditionally labeled ``iterative'' (\textit{iterativní} in Czech) or ``multiplicative'' (\textit{násobená} in Czech)   is  taken to be clearly separate from the imperfectivizing suffix (see e.g., recently \citealt{biskup2024} and references therein), on the one hand, but on the other hand, it is still subsumed under imperfective aspect (see e.g., \citealt{Genis2021}, \citealt{biskup2024}, and also Czech reference grammars, including \citealt{kopecny1948}, \citealt[Vol.~2,~p.~180]{petr1986mluvnice}, \citealt{lamprecht1986historicka}, \citealt{Kosek2014}).  I use the dashed line here to indicate that its subsumption under imperfective aspect is problematic, given that, as I argue in this paper, the class of verbs on which it occurs does not neatly fit the imperfective paradigm.
It is not uncommon 
for the generic marker \textit{-va-} 
and the imperfectivizing suffix to be confounded, both notionally and formally, given that their surface allomorphic variants may often seem alike, and both 
occur in generic sentences. 
One compelling piece of evidence for their formal independence 
stems from the fact that they may co-occur on the same verb (see \citealt{filip1997suigenva}). Consider the following attested example:\footnote{The example is from \url{https://abecedazdravi.cz/diskuse/psychicke-problemy-a-psychologie/deprese-i-nova-diskuze/89}; accessed on July 9, 2023.}

\ea{\label {bouquet.va} 
\gll Bývalý mi taky \textit{dávával}  nádherně a originálně vázané kytice, přímo umělecká díla, \textellipsis \\
ex me.\textsc{dat} also give.\textsc{ipfv.gen.pst} wonderfully and originally bound bouquets directly art works\\}
\glt `My ex also used to give me wonderful bouquets in original arrangements, real works of art, \ldots'
\z 

\noindent
In  (\ref{bouquet.va}), the morpheme sequence which is here glossed with \textsc{ipfv.gen} is the result of the derivation pattern (a) > (d) > (e), given
in \figref{fig:filip:1}. In the first step (a) > (d) the imperfectivizing suffix (\textsc{ipfv}) applies to a (primary, or root) perfective base and derives an imperfective form, which alternates between an episodic and a generic reading depending on context. Hence, it is unmarked for genericity, i.e., non-generic or  a ``general imperfective'' form
(in the sense of \citealt{Comrie1976}).
In the second step (d) > (e), the generic \textit{-va-} `\textsc{gen}' is added ``on top of'' the imperfectivizing suffix (\textsc{ipfv}), and yields a verb that only has a generic reading.  
%Again, while this derivational pattern is here illustrated with the examples from Czech, similar examples can also be found in Slovak and Polish, and are attested in other Slavic languages, as well. 

Following the standard assumption about what it means to be a ``formal member of a category system'', given that the Slavic imperfectivizing suffix (\textsc{ipfv}) is a marker of imperfectivity, as all agree,
the  morpheme \textit{-va-} `\textsc{gen}' cannot be a different marker of the same imperfective category.
A given formal member of a category system stands in a complementary distribution to other formal members of the same category, morphologically and syntactically. For instance, the present and past tense morphemes in English are formal members of the same tense category, hence are in complementary distribution: e.g., *\textit{she laugh-s-ed} or *\textit{she laugh-ed-s}. The markers of tense and progressive aspect in English can co-occur, because they are \textit{not} formal members of the same grammatical category: e.g., \textit{she was/is/will be laughing}.
If we accept this line of reasoning for the present and past tense morphemes in English, we should be ready to accept that
the fact that the generic marker \textit{-va-} and the imperfectivizing suffix do not stand in a paradigmatic opposition to each other indicates that they are not only two separate morphemes, but also that they cannot be two different markers of the same grammatical category system, and specifically both marking imperfectivity.

To the above argument it could be objected that
the morpheme sequence which is here glossed with \textsc{ipfv.gen} (as in (\ref{bouquet.va})) instantiates two occurrences
of one and the same imperfectivizing suffix: `\textsc{ipfv} + \textsc{ipfv}'. Hence, what we have
in (\ref{bouquet.va}) is something like a doubly aspectualized imperfective verb associated with a generic (or habitual) reading.
Implicit behind this idea would seem to be
some encoding  economy principle, namely that it is
superfluous to apply twice one and the same imperfectivizing suffix with one and the same imperfective function. Therefore, the second application of \textsc{ipfv} must
have a different function, i.e., affect a truth-conditional change from a non-generic, or general imperfective, to a generic imperfective meaning.

One immediate prima facie counterargument against the
`\textsc{ipfv} + \textsc{ipfv}' double aspectualization idea has to do with verbs that only have a generic interpretation, but
contain a \textit{single} occurrence of the overt
alleged  \textsc{ipfv} suffix. This is a subclass of generic verbs derived following the pattern (c) > (e) in \figref{fig:filip:1}.
In order to rescue the double aspectualization idea it might be proposed that simplex (primary) imperfective verbs like \textit{štěkat} `to (be) bark(ing)',
in \figref{fig:filip:1} under (c), contain a phonologically null imperfectivizing suffix
$\varnothing$\textsuperscript{\textsc{ipfv}},
%i.e., on the `lower imperfective-form of the verb',
and so generic verbs that are formed following the pattern (c) > (e) contain one covert and
one overt instance of 
% the morpheme sequence glossed as `\textsc{ipfv.gen}'.
the same \textsc{ipfv} suffix: $\varnothing$\textsuperscript{\textsc{ipfv}} + \textsc{ipfv}.\footnote{That Slavic primary imperfectives introduce a null imperfective suffix is
proposed in syntactic accounts of Slavic aspect, see e.g. \textcitetv{chapters/07_gronn}.}
However, such a move is hardly empirically motivated at all, theoretically costly, and generally supports a proliferation of empty categories for purely theory-internal reasons.

Reduplication without a change in the truth conditions is common in natural languages, it may serve to add affective connotations, intensification, as we have seen above in (\ref{redup.battle}) and (\ref{redup.custom}).
We would have to ask: Just why should the reduplication of
two overt \textsc{ipfv} morphemes have a truth-conditional consequence, and enforce the generic interpretation, when
a truth-conditional change is not a necessary consequence of all reduplication processes?
Similarly, we may ask: Just why
should a double aspectualization consisting of one overt  \textsc{ipfv} morpheme and one covert $\varnothing$\textsuperscript{\textsc{ipfv}} morpheme
have a truth-conditional consequence such that the second overt occurrence of this morpheme restricts the verb to only a generic interpretation?

Moreover, the \textsc{ipfv} morpheme occurs on verbs
that are general imperfective verbs, because
they can  alternate
between an episodic (including progressive) and a generic interpretation depending
on context (as
imperfective morphemes in natural language do,
see e.g., \citealt{Comrie1976}).
Presumably the second overt occurrence of the alleged \textsc{ipfv} should be also inherently amenable to either the episodic or generic
interpretation, but then it is unclear just why the second overt occurrence of \textsc{ipfv} should enforce the
generic interpretation, and not the episodic
interpretation of its output.
For instance, it is unclear why it does not constrain its output to only
the episodic progressive interpretation.
The specifically generic properties of verbs containing the
sequence \textsc{ipfv}+\textsc{ipfv} or
the sequence $\varnothing$\textsuperscript{\textsc{ipfv}}+\textsc{ipfv} would have to be postulated, because they do not follow from the semantics of either \textsc{ipfv} or $\varnothing$\textsuperscript{\textsc{ipfv}}.

A more sophisticated version of the reduplication or double aspectualization idea is
proposed by \citet{cable2022}: namely, the Czech ``habitual'' (in his terms, generic in my terms) morpheme, as explored in \citet{Filip2018incable}, is, according to Cable ``etymologically an instance of imperfective-aspect
combining with another, lower imperfective-form of the verb, a kind of doubly aspectualized verb \citep{Filip2018incable}.
One might wonder, then, whether that higher imperfective-morphology might
at all synchronically be a realization of a (bound) T-head, while the lower aspectual morphology is a realization of (non-modal) [IMPRV\textsubscript{OG}]" (\citealt{cable2022}: 44--45) [where OG stands for ``ongoing'', HF].
This suggestion builds on Cable's analysis of habitual morphology
in Tlingit (Na-Dene; Alaska, British Columbia, Yukon),
where a verb may be realized in the ``habitual perfective-mode'' formed with the habitual suffix \textit{-ch}, or in the ``habitual imperfective-mode'' consisting of the imperfective form followed by the habitual particle \textit{nooch}.
The main thesis is that there are two different morphosyntactic structures that underpin habituality both within and across languages. One where
the habitual semantics is directly contributed by imperfective aspect, which has the
meaning of a modal operator quantifying over alternative worlds/situations (see e.g., \citealt{ferreira2016}).
In Tlingit, it is only verbs bearing imperfective-mode that can be used to describe pure dispositions, non-actualized regularities.
The other structure underpinning habituality corresponds to the bi-partite morphosyntactic structure, where
the habitual morphology is the realization of
a quantificationally bound T(emporal)-head, licensed when it is bound by a (possibly covert) Q-adverb,
which combines with the lower perfective or imperfective Asp-head to yield a complex form with a habitual meaning.\footnote{This licensing strategy resembles that of \citet{chierchia1995}, who assumes that genericity is subsumed
under aspect, and ``all languages have a distinctive habitual morpheme (say, Hab) which
can take diverse overt realizations" \citep[197]{chierchia1995}. Chierchia's Hab is a functional head in an (imperfective) aspectual projection, and
either realized by explicit aspectual morphemes, or, as in English, or it is covert. Hab
carries the agreement feature [$+$Q]  which in its specifier requires some overt Q-adverb \citep[197--198]{chierchia1995} or the \textbf{Gen} operator \citep[210]{chierchia1995}. The idea that both stand in a local licensing relation to Hab\textsubscript{[$+$Q]} is inspired by the analysis of NPIs (\citealt{ladusaw1979polarity}), which effectively makes genericity a polarity
phenomenon \citep[212ff.]{chierchia1995}. The main innovation is the
analysis of ILPs as inherent generics,
capitalizing on the assumption that polarity phenomena are rooted in the lexicon (with effects on sentences). The lexical entries of ILPs have an abstract habitual morpheme
Hab\textsubscript{[$+$Q]}, whereby its agreement feature [$+$Q]
can be locally licensed \textit{only} by \textbf{Gen}, and never by overt Q-adverbs.
ILPs are also taken to involve generic quantification over situations.
Their Davidsonian situation variable $s$ is bound by \textbf{Gen}
making it unavailable for binding by overt Q-adverbs.
Aside from the problems with \citeauthor{chierchia1995}'s (\citeyear{chierchia1995}) proposal,
some of which are detailed in \citet{fernald2000}, the Czech morpheme \textit{-va-} cannot be the spell-out of the covert functional head Hab\textsubscript{[$+$Q]}. If it were, it should have
the same semantic properties as its covert counterpart in the logical representation of primary and secondary
imperfectives on their generic interpretation. However, it does not. Moreover, if
\textit{-va-} were the spell out of Hab\textsubscript{[$+$Q]}, it should be uniformly incompatible with all ILPs, because their Hab\textsubscript{[$+$Q]}
can be locally licensed \textit{only} by \textbf{Gen}, and never by overt Q-adverbs.
However, \textit{-va-} is compatible with ILPs provided they allow for an ``interruption'' construal, which also sanctions their co-occurrence
with overt Q-adverbs (see \sectref{mustnotbeused}). }
The Tlingit habitual
morphology is a kind of quantificationally dependent tense,
where it is the (potentially covert) Q-adverb
that is the source of the habitual meaning.
It is analyzed as a temporal quantifier, e.g., `always', `sometimes', which quantifies strictly over times in the actual world, it does not introduce modal quantification over other possible worlds, because, according to Cable, it comes with the actualization entailment.

\citeposst{cable2022} double aspectualization proposal for Czech does not work for three main reasons. First, it raises the questions posed above about just why the verb ``doubly aspectualized'' with the same imperfectivizing suffix should
have its interpretation constrained to the generic (Cable's ``habitual'') reading, when the imperfectivizing suffix covers the whole imperfective domain, both its episodic and generic subdomains. Why does not the ``double imperfective aspectualization'' exclude the generic interpretation, and
enforces only an episodic (e.g., progressive) interpretation?

Second, \citet{cable2022} proposes that habitual
morphology, within and across languages, is a kind of quantificationally dependent tense. However, generally, the semantics of generic sentences and habitual sentences as their special case cannot be analyzed in terms of temporal notions, or reduced to the semantics of Q-Adverbs (see \sectref{CharGenericity}, and references therein).
As has been shown above, genericity, and habituality as its special case, is independent from tense, both notionally
and formally. Specifically, the Czech generic morpheme \textit{-va-} (a habitual morpheme in terms of \citealt{cable2022}) is
fundamentally independent of tense.

Third, as many others,
\citet{cable2022} correlates the so-called actualization entailment of habitual sentences with the lack of intensionality (modality). However, all generic sentences, including habitual sentences (i.e., concerning habits in the proper sense), be they formally unmarked or marked by means of dedicated markers like the Czech \textit{-va-}, have modal import and predictive force, precisely because
they express some habit or regularity. As such they have implications for other accessible worlds, than just the actual world, and perhaps all possible worlds. The modality of habitual sentences is not only compatible with their actualization entailment, or better realization inference, but is their \textit{sine qua non} semantic property.
Hence, a purely temporal analysis in terms of quantifiers
that strictly quantify over times in the actual world will not do to capture the meaning of habitual sentences.

Finally, there is another problem faced by Cable's proposal. \citet{cable2022} assumes
that any modality of habituals must be contributed by imperfective morphology, or some other material in the sentence, which has the meaning of a modal operator quantifying over alternative worlds other than the actual one, but never by perfective morphology, which is assumed to be non-modal. However, not only is there no notional and formal incompatibility between perfectivity and genericity as well as habituality (specifically for Slavic languages see \citealt{filip1993berkeleyva, filip1994mitva, filip1997suigenva}, and elsewhere), which Cable does not deny, but also perfective forms when used to
express generic and habitual statements possess a modal semantics, and so necessitate quantification over alternative worlds other than the actual one. For instance, the English generic sentence (\ref{speaker}),
which has a purely intensional interpretation,
and is true despite having no actualized instances that count as evidence for its truth, and may never have any, is
naturally rendered with a perfective verb in Czech:

\ea \label{speaker.Czech}
\gll Po viceprezidentovi \textit{nastoupí}\textsuperscript{PFV} na místo prezidenta předseda sněmovny.\\
after vice.president ascend.\textsc{prs.3sg} on place president chairman assembly\\
\glt `The Speaker of the House succeeds the Vice President.'
\z

\noindent This leads me to make some observations on the use of perfective forms in generic (and habitual) sentences,
as a conclusion to this section.
Some examples of generic sentences with naturally occurring
perfective forms from Slavic languages, Russian, Serbo-Croatian and Czech,
as well as Romance (Italian) are given below:\footnote{The Czech example (\ref{czech.pf1}) is taken from \url{https://www.kosmas.cz/oko/ukazky/430121/o-panovnici-elite-prezidentu-povysenci-a-polarnikovi-v-nervoznim-pohybu/}, accessed July 9, 2023.}

\ea \label{russian}
\ea
\gll Nastojaščij		drug		vsegda	pomožet\textsuperscript{PFV}.\\
real				friend	always	helps \\
\glt `A real friend will always help you.’\hfill (a Russian proverb)
\ex \label{serbo-croatian}
\gll Svako jutro			popijem\textsuperscript{PFV} 	čašu rakije. \\
every morning	\textsc{pref.}drink			glass brandy		\\
\glt `Every morning I drink a glass of brandy.’\\\hfill (Serbo-Croatian; \citealt[62]{monnesland1984})
\ex \label{czech.pf1}
\gll To se jeden nalítá\textsuperscript{PFV} v úřadovnách prezidentského paláce.\\
it \textsc{refl} one \textsc{pref}.fly in offices president palace\\
\glt `What a runaround one must take in the offices of the presidential palace.'\hfill (Czech)
\ex { \label{italian}\gll Sempre, quando 	mi vide\textsuperscript{PFV}, il custode  aprì\textsuperscript{PFV}  la porta.  \\
always 	when me saw the janitor 	opened 	the door \\}
\glt `Always when he saw me, the janitor opened the door.’\\\hfill (Italian; \citealt[508]{Bonomi1997})
\z\z

\noindent
\citet{dahl1995} claims that
the use of perfective forms in generic (and habitual) sentences in Slavic languages is another ``case of deviant behavior'' \citep[420]{dahl1995} of their aspect systems.
One must wonder about the justification for this claim, given
that perfective forms are commonly used in generic sentences in Romance languages, for instance, as the above example from Italian shows. What is noteworthy though is that perfective forms are
commonly used in generics and habituals in West Slavic languages.\footnote{For a recent study on the uses of the perfective present in Polish generic (and habitual) sentences see \citet{wiemer2025}.}
In East and South Slavic languages they are less common \citep{Forsyth1970,fortuin2015typology,wiemer2017diachrony},
but perhaps not as rarely used as is generally thought \citep{wiemer2023}.\footnote{In West Slavic language the use of perfective forms for the expression of  generic and habitual statements is common, though their frequency as well as the functional range depends on the type of generalization to be expressed, and lexical idiosyncracies of the relevant (im)perfective forms. For instance, in Czech, perfective verbs are common or even preferred in cooking recipes, assembly instructions, game rules like rules of chess, proverbs, to name just a few types of generic statements.}

Not only is perfectivity notionally compatible with genericity (and habituality),
but also their respective overt markers are compatible,
because they can co-occur on the same verb form: e.g., in Lithuanian ``frequentative past'' \citep{dambriunas-etal-1966}, or in Awa (New Guinea) \citep{loving1964mckaughan}.

\newpage
In sum, genericity (and habituality) is orthogonal to the imperfective and perfective distinction (see also \citealt{filip1997suigenva}).
Acknowledging this separation of grammatical
aspect from genericity (and habituality) is not to deny that there
are close semantic affinities between imperfectivity
and genericity, as well as habituality,
as its special case.\footnote{For instance, aspectually stative predicates in Slavic languages are expressed by imperfective verbs, and aspectual stativity is a key semantic property of
generic predicates. Moreover, from the grammaticalization point of view, the expression of genericity (habituality) is tied to markers that become grammaticalized on the path from (present) progressive markers to imperfective ones \citep{bybee1989dahl}.}
Moreover, acknowledging that genericity (and habituality) is orthogonal to the imperfective and perfective distinction
means that markers like the Slavic morpheme \textit{-va-}
that systematically enforce
only a generic (or a habitual) interpretation of sentences
need not be subsumed under imperfective aspect.

In support of this point, in the next section, I
show that the Czech generic morpheme \textit{-va-}
contributes semantic and pragmatic properties
to the meaning of generic sentences
that their corresponding sentences headed by non-generic imperfective verbs, both primary and secondary, do not have. This also supports my conclusion that generic
\textit{-va-}marked verbs do not neatly fit the Slavic imperfective paradigm,
and ``general imperfective'' paradigm in natural language (in the sense of e.g., \citealt{Comrie1976}).




\section{Semantic properties of the Slavic generic morpheme}\label{semprag}


\subsection{When the generic morpheme must, must not and may be used}\label{delimitingVA}


This section identifies cases in which
verbs marked with the generic morpheme \textit{may},
\textit{must} and \textit{must not}
be used in contrast to imperfective verbs, which are not formally marked for genericity with this morpheme.
The examples come from Czech. However, if a given Slavic language
has a verb with the generic marker and could be used in the same context as that of a corresponding Czech verb, it is predicted that the same observations
should hold for it that are made for Czech.


\subsubsection{When the generic morpheme \textsc{must not} be used}\label{mustnotbeused}

\noindent
The generic morpheme \textit{-va-} must not be used in episodic contexts, where the reference is (i) to a particular on-going situation (progressive) or
(ii) to a closed set of plural situations at a given world and time; neither can it be used
(iii) for the expression of generalizations that permit no exceptions, nor (iv)
for the expression of generalizations with unrealized particulars in the actual world, i.e., in only dispositional, purely hypothetical generalizations.


\subsubsubsection{Lack of reference to a particular (on-going) situation}

One of the well-known tests tapping into the semantics of the Slavic
perfective/imperfective distinction is the possibility
of reference to a particular episodic situation that is on-going at a specific moment of time,  i.e.,  the possibility of the progressive interpretation. All imperfective verbs that denote stage-level predicates pass this test, but those denoting individual-level (stative) predicates do not, and neither do perfective verbs.

Most importantly, all verbs marked with the generic morpheme \textit{-va-} also fail this test.
In traditional Bohemistics, this property is well-known as their hallmark property,
which is referred to as the ``non-actuality'' or ``atemporality'' property at least since \citet{kopecny1948}. (See also \citealt{kopecny1962,kopecny1965}, \citealt[259]{kopecny1966}.)
The example below, which is adapted from \citet{kopecny1948},
shows that Czech verbs marked with \textit{-va-} cannot be felicitously used in answers to questions about what is \textit{actually} going on `now' at the speech time. In contrast,
their imperfective counterparts without the morpheme \textit{-va-} can.
It is important to emphasize
that the evaluation with respect to ``extended-now'' that covers extended intervals of time is irrelevant for this test:

\ea \textit{Context}: `What are you doing right now?'
\ea[\#]{{\label{kopecny.dialogue} \gll Psávám 	jeden	naléhavý	dopis. \\
write.\textsc{gen.1sg.prs}  one   well-known		letter\\}
\glt `I write one urgent letter.' (as a rule/as a habit/often/rarely/usually)}
\ex[]{{\label{bla}
\gll  Píšu\textsuperscript{IPFV} 	jeden		naléhavý 	dopis. \\
write.\textsc{1sg.prs}	one 		urgent		letter\\}
\glt `I am writing one urgent letter.'}
\z
\z
\noindent
The generic morpheme \textit{-va-} if it is attested on a verb in other Slavic languages will also give rise to the oddity that we see in (\ref{kopecny.dialogue}).

The  semantic property of ``non-actuality'' or ``atemporality'' contributed by
the Czech morpheme \textit{-va-}
along with its productivity and compositional transparency  (see also \sectref{basicmorphology}) lead
\citet{kopecny1948} to propose  that it is a grammatical marker of a ``third aspect'', namely an ``atemporal aspect'' that is independent of the Slavic perfective/imperfective distinction (see also \citealt{Kosek2014}).
As a wide-reaching consequence, \citet{kopecny1948} rejects the claim made by  \citet{koschmieder1929} that Indo-European languages lack a grammatical category of ``atemporality''.\footnote{\citet{kopecny1948}: ``Není tedy zcela správné tvrzení Koschmiedrovo v jeho Nauce o aspektech (33), že indoevropské jazyky nemají [153] gramatickou kategorii atemporálnosti (mimočasovosti, ``pozaczasowości").'' [\citeauthor{koschmieder1929}'s claim that Indo-European languages have no grammatical category of atemporality (extra-temporality, `pozaczasowość' [Polish]) is not entirely correct. (my own translation)]}
On Kopečný's third ``atemporal aspect'' hypothesis, the Czech verb system
might be schematically represented as in \figref{fig:filip:3},
albeit recast in terms of the non-generic/generic distinction, as it
is used in the contemporary literature and in this chapter.

\begin{figure}
%\includegraphics[width=11.5cm]{KopecnyThirdAspect.comp.png}\\
\begin{forest}
[Classification of Czech verbs, calign=child, calign child=2
    [perfective aspect,name=perf]
    [imperfective aspect,name=imperf]
    [“atemporal aspect”\\grammatical marker -\textit{va}-,name=atemp
    ]
]
\node [below=1mm of atemp,font=\scshape] (generic) {generic};
\path let \p1 = ($ (perf) !0.5! (imperf) $), \p2 = (generic)
  in node [font=\scshape] at (\x1,\y2) {non-generic};
\draw [line width=1pt, decorate, decoration={calligraphic brace, amplitude=5pt, mirror, raise=5pt}]
  (perf.south west) -- (imperf.south east);
\end{forest}
    \caption{Classification of Czech verbs (following \citealt{kopecny1948} and elsewhere)}
    \label{fig:filip:3}
\end{figure}

Aligned with Kopečný's ``third atemporal aspect'' hypothesis,
\citet[325]{lamprecht1986historicka} states that Czech verb forms
that are built with the generic morpheme \textit{-va-} (`generic' in my terms)
are clearly not imperfective forms, because imperfectives are formed by other morphologically means.
Also \citet[152]{Kosek2014} observes that in Czech the ``third aspect'' with its dedicated formal expression by means of the morpheme \textit{-va-} has been established since the 16th century,
when it became independent of the perfective/imperfective distinction. However, Kosek, unlike Kopečný, erroneously views it as a means for the expression of iterativity, or frequentativity.

%Similarly, \citet{biskup2024}, also building on \citet{filip1997suigenva}, treats the Slavic generic morpheme as separate from the imperfectivizing suffix, on both formal and notional grounds.


\citeposst{kopecny1948} ``third atemporal aspect'' hypothesis is of interest here in so far as it highlights the hallmark ``atemporal'' or ``non-actuality'' property of the Czech morpheme \textit{-va-}, i.e., that it formally marks sentences that lack reference to a specific situation in all its occurrences.  As observed above (see \sectref{CharGenericity}), it is one of the key properties
of generic predicates in natural languages generally, which makes them aspectually stative. Kopečný's account of the Czech morpheme \textit{-va-} is a notable improvement on its traditional
treatment as an imperfective marker carrying the  iterative or multiplicative feature, which it quite clearly lacks
(as will be shown below).

In sum, based on the above observations, the first key property of Czech verbs that are formed with the generic morpheme can be stated as follows:

\eanoraggedright\label{ex:semprop1}
\textsc{Semantic Property 1} (to be revised): The generic morpheme \textit{-va-} derives verbs that systematically lack reference to particular situations. It is a property at the core of stativity,  and one of the properties of generic predicates in natural languages.
\z

\noindent
Semantic Property 1 in (\ref{ex:semprop1}) clearly separates the generic morpheme  from the imperfectivizing suffix in Slavic languages, which formally marks secondary imperfectives that, as a whole class, are not  aspectually stative. Just like primary imperfectives, such as \textit{píšu} `I am writing' in (\ref{kopecny.dialogue}), secondary imperfectives freely alternate between episodic and generic readings, and can be used with reference to a specific situation that is on-going at some given reference time, as the following attested example shows:\footnote{The example is from \url{https://www.praha11.cz/redakce/tisk.php?lanG=cs&clanek=9892&slozka=7155&as4uOriginalDomain=www.praha11.cz&as4u_protocol=https&}; accessed on August 15, 2025.}

\ea \label{sec.IPFV.prog}
\gll V křoví právě \textit{přepočítával}\textsuperscript{IPFV} bankovky v peněžence, strážníky na sebe upozornil světlem z mobilu.\\
in shrubs just \textsc{pref}.count.\textsc{ipfv.pst.3.sg} banknotes in wallet policeman\textsc{pl.acc}  on himself make.aware\textsc{pst.3.sg} light from mobile\\
\glt `At that moment he was in the shrubs counting money in his wallet, and the light of his mobile phone alerted the police to him.'
\z

\subsubsubsection{Incompatibility with iterativity, or a mere plurality of situations}

In general, the lack of reference to a specific situation, which is a hallmark property of generic predicates (related to  their stativity), is closely tied to their incompatibility
with iterative adverbials (or cardinal count adverbials) like `three times' or `several times'. Consider the following example:

\ea[\#]{{\label{chessva3times}
\gll Včera Pavel 	 hrával    třikrát 		šachy s dědou. \\
yesterday  Paul   play.\textsc{gen.pst}
three.times	chess with grandpa \\}
\glt \#`Paul used to play chess with grandpa three times yesterday.'}
\z

\noindent The oddity of the above sentence must be due to the clash between the generic marker \textit{-va-} and the iterative adverbial
`three times' which counts three episodes that occurred within a single interval of time `yesterday'. That is, `yesterday' enforces an episodic reading of a sentence, which must be evaluated with respect to a particular world/time,
but the generic morpheme \textit{-va-} requires evaluation at alternative possible worlds than just one (actual) world.

The generic morpheme \textit{-va-} can scope over the iterative adverbial, if the context does not indicate otherwise, as we see \REF{filip:ex:write-three}. Here the closed sequence of iterated three eventualities forms the basis of the generalization which is a part of a larger pattern
that holds every day, and across different possible worlds.

\ea\label{filip:ex:write-three}\gll Psával mi aspoň třikrát každý den. \\
write.\textsc{gen.pst} me at.least three.times every day\\
\glt `He used to write to me at least three times every day.'
\z

\noindent In (\ref{chessva3times}), however, `yesterday'
restricts the reference to a single interval/world which excludes a generic interpretation of (\ref{chessva3times}), and consequently the use
of the generic morpheme is odd.
The corresponding non-generic sentence (\ref{chess3times})  without \textit{-va-} is perfectly acceptable with the iterative adverbial `three times'.
The main verb here is a non-generic primary imperfective verb:

\ea{\label{chess3times}
\gll Včera Pavel 	hrál\textsuperscript{IPFV}		třikrát 	šachy s dědou.\\
yesterday Paul   	play.\textsc{pst} 	three.times	chess with grandpa \\}
\glt `Paul played  chess with grandpa three times yesterday.'
\z

\noindent Secondary imperfective verbs formed with
the imperfectivizing suffix (\textsc{ipfv}),
which are also non-generic like primary imperfectives,
are compatible with iterative adverbials like `three times':

\ea {\label{win3times}
\gll Sice jsem   již	      třikrát  vyhrávala\textsuperscript{IPFV}, ale  nikdy jsem   nevyhrála\textsuperscript{PFV}.\\
while \textsc{aux}  already  three.times \textsc{pref}.play.\textsc{ipfv.pst}   but never \textsc{aux}	   \textsc{neg}.\textsc{pref}.play.\textsc{pst}\\}
\glt `Although I was already winning / was about to win three times, I never won.'
\z

\noindent
The incompatibility with iterative adverbials is a test that generally separates verbs with the generic morpheme \textit{-va-}, on the one hand,
from primary and secondary imperfective verbs, on the other hand;
only the latter can co-occur with iterative adverbials, modulo their lexical semantics and context in sentences that describe episodic situations that are not a part of a larger pattern. This also means that verbs
with the  generic morpheme \textit{-va-}, on the one hand, and secondary imperfective verbs with the imperfectivizing suffix (\textsc{ipfv})
along with primary imperfectives, on the other hand, make different contributions to the meaning of a sentence. This provides one piece of evidence in support of separating verbs with the generic morpheme \textit{-va-} from primary and secondary imperfectives.

Generally, the incompatibility of a morpheme with iterative  adverbials  means that  it cannot be assimilated to markers in natural languages that  are variously labeled as iterative, multiplicative, frequentative, repetitive, pluractional, and the like, of the kind that are used to encode a mere plurality of situations
(i.e., a closed series of particular episodes or situations of the same type at a given world and time).
The incompatibility of the generic morpheme \textit{-va-} with iterative adverbials
clearly shows that its common labels ``iterative''  or  ``multiplicative'' and the like
are misnomers at best (see also \citealt{bily1986iterative}), or a category error at worst (\textit{pace} standard Czech reference grammars, e.g., \citealt{petr1986mluvnice}, \citealt{Kosek2014}, \citealt{lamprecht1986historicka}, \citealt{nuebler2017} or  \citealt{dahl1995}).

Based on the above data and observations, we  may update Semantic Property 1 in (\ref{ex:semprop1}) as follows:

\eanoraggedright\label{ex:semprop1final}\textsc{Semantic Property 1} (final version): The generic morpheme \textit{-va-} derives verbs that lack reference to particular situations. It is a property at the core of stativity, and one of the properties of generic predicates in natural languages.\smallskip\\
\textit{Corrolary}:  The  generic morpheme \textit{-va-} derives verbs
that are incompatible with iterative adverbials like `three times', `several times' used to denote a cardinality of a closed set of situations at a given time interval and world (i.e., a set that is not a part of a larger pattern, or generalization).
\z

\noindent Semantic Property 1 is related to independent arguments made elsewhere that genericity (and habituality) is clearly separate from iterativity, and the two must not be conflated and confounded, both on formal and notional
grounds (see e.g., \citealt{CarlsonHabGen2012}).\footnote{See also e.g. \citet{BertinettoLenci2012} for an overview of notions related to iterativity and their delimitation from genericity and habituality.}
Iterativity involves a closed set of a plurality of
events, eventualities, situations, or cases of the same type (or
a plurality of time intervals associated with them),
but this has virtually nothing to do with habituality and genericity, i.e., with what gives rise to a generalization, a regularity, or a pattern (see e.g. \citealt{krifka1995genericity,filip1993berkeleyva,carlson2005,carlson2008patterns, CarlsonHabGen2012}, i.a).
It is neither sufficient nor necessary for the truth of habitual and generic statements. For instance, a mere plurality of situations is not sufficient, because it can be a part of the meaning of
what clearly are non-generic statements, be it a sequence of temporally
close situations that constitute a single complex situation (e.g.,
\textit{Last night, he knocked three times on the door}), or multiple separate situations (e.g., \textit{During our train ride to Rochester, the conductor came to our car three times to check our tickets}), which may also be spread over a large interval (e.g., \textit{During his lifetime Franz met Milena only twice}).
A closed set of a plurality of situations of the same type is not necessary for the truth of generic statements, because some have a purely
intensional interpretation with no actualized verifying particulars, as in (\ref{speaker}). Moreover, there are generic statements with ILPs whose truth is not inferrable from any array of a particular situations at all, such as \textit{Bob owns a yacht} that presupposes one situation of acquisition like buying \citep[38]{krifka1995genericity}, or \textit{Bob is a bachelor} for which it makes no sense to try to
identify particular situations that count as evidence for its truth
\citep{carlson1995truth}.  Such observations also indicate
that genericity and habituality should not be confounded in analyses of phenomena referred to as iterative aspect, event-internal or event-external iterativity, frequentativity, pluractionality, multiplicativity, iterated semelfactivity and the like (\textit{pace} e.g., \citealt{nef1986semantique, geenhoven2001noun,geenhoven2004adverbials, scheiner2003temporal}).

\largerpage
The conflation of iterativity (i.e., a mere plurality of situations) with habituality and/or genericity
is also common in the analyses of habitual and generic readings of imperfective sentences.
%that (implicitly or explicitly) subsume them under imperfective aspect in the spirit \citet{Comrie1976}. (Recall that Comrie's `habituality' corresponds to `genericity' as it is used here, given the wide range of examples that Comrie covers by his `habituality', see \sectref{CharGenericity}).
%As is well-known, and observed above, Slavic non-generic imperfectives (see \ref{fig:filip:1}), and generally `general imperfective' forms in the sense of \citealt{Comrie1976}, cover the whole semantic domain of imperfectivity, i.e., both its episodic (progressive) and `habitual' (generic) subdomains.
A common strategy is to add to the imperfective operator a parameter capturing the number requirement on its application.
%, yielding a predicate of singular eventualities in the case of the episodic progressive interpretation, and a predicate of plural eventualities in the case of habitual (or generic) interpretation.
%Apart from this number difference, the logical forms for both the episodic and habitual/generic readings of imperfective sentences are the same, but may include additional elaborations of their respective temporal and/or modality components.
Examples of this strategy can be found
in \citet{ferreira2016} (also \citealt{ferreira2004,ferreira2005}, and elsewhere), and \citet{hacquardDISS}, for instance.
Imperfective forms  \textit{either} denote singular eventualities under their episodic progressive interpretation, \textit{or} a plurality thereof
under their habitual/generic interpretation.\footnote{Assuming that VPs denote
sets of time intervals, and VPs like NPs can be singular or plural, as e.g. \citet[358]{ferreira2016} proposes,
the progressive (episodic) interpretation of a VP amounts to
singular time intervals at which it holds and the ``habitual'' (generic) interpretation to plural time intervals.}
For the latter, \citet[29]{ferreira2005} assumes a phonologically null ``habitual/generic operator \textit{hab}" when they
occur without an overt Q-adverb or restrictor.
Both the episodic progressive and habitual/generic interpretations share the same temporal component, namely existential quantification over time intervals, and also the same modal component with modal quantification over possible worlds,
adapted from the analyses of the progressive aspect (as developed in \citealt{Dowty1979,Landman92,portner1998progressive}). In short,
the habitual/generic  operator \textit{hab} is the plural counterpart of the progressive operator
Prog, where Prog is defined on singular (atomic) eventualities.
The Prog truth conditions
are stated with respect to a pair of an interval and a world $\stb{i,w}$,
iff there is a singular eventuality \textit{e} in \textit{w} and modal continuations in possible worlds in which it terminates or is completed.
A habitual (generic) sentence like \textit{John plays soccer (regularly)} is true in all the worlds in which there were plural eventualities of John playing soccer before the utterance time, and there is an ongoing plural eventuality of John playing soccer at some reference time, whereby all potential external obstacles to the continuation of this plural eventuality are removed.

Setting aside that a proper characterization of the modal continuation component of Prog is still outstanding, and that the meaning of habitual/generic sentences cannot be analyzed in terms of a plurality of events (or plurality of time intervals), one obvious problem with this proposal is that it does not cover sentences with a purely dispositional reading like
\textit{Mary handles the mail from Antarctica},
which may be true even if Mary never handled any mail from Antarctica before the utterance time, and such mail may never ever arrive in the future. Neither can it handle all the relevant aspects of the meaning of bare habitual sentences like \textit{John plays soccer}.
This sentence may mean that John is not averse to playing soccer, he has an ability or disposition for playing soccer, and I can truthfully assert it,
based on my experience of having seen him play soccer once in the past.
This allows me to conclude that if John were in a situation to play soccer, he just might likely play soccer. \textit{John plays soccer} is true, even if there are no multiple occasions of John playing soccer at some reference interval including the reference time,
and John may never play soccer in the future.
Moreover, suppose you asserted that John does not engage in any sports activity. I can contradict you by uttering \textit{John plays soccer},
even if I have only seen John play soccer once in the past.
But this cannot be captured if we assume that ``habituals are about plural events" \citep[353]{ferreira2016} combined with the modality that has been associated with or is akin to Prog.


\subsubsubsection{Generalizations permitting no exceptions}\label{sec:no-exceptions}

One of the key properties of generic sentences is that most are exception-tolera\-ting, but
are also compatible with no exceptions (see \sectref{CharGenericity}).
%An exception-tolerating generalization separates a given domain into its confirming cases and non-confirming cases or exceptions.
The generic morpheme, as I argue, has a meaning
that combines a generalization assertion with an exception or a counterexample inference
that there are or may be exceptions to the generalization.
\textit{Exceptions} to the generalization are compatible with its truth, while \textit{counterexamples} are incompatible with its truth (see e.g. \citealt{nickel2010}). This property was already discussed by \citet{filip1993berkeleyva,filip1994mitva,filip2019harvardva, filip2021cologneva}, and here, it is captured as follows:

\eanoraggedright\label{ex:semprop2} \textsc{Semantic Property 2}: The generic morpheme \textit{-va-} contributes the uncancellable inference (as part of its truth-conditional content) that there are actual or possible exceptions or counterexamples to the generalization.\smallskip\\
\textit{Corollary:} The generic morpheme \textit{-va-} is incompatible with universal generalizations (e.g., universal laws of nature), or any generalization that is commonly known to permit \textsc{no exceptions}, possibly in \textsc{all the known conditions}.
\z
\noindent
The clearest examples of Semantic Property 2 (in (\ref{ex:semprop2})) are sentences
in which
the  generic morpheme \textit{-va-} clashes with universal laws of nature (as in (\ref{earthva})) or definitional statements that necessarily hold for all the members of a given kind (as in (\ref{waltz})). Such exceptionless regularities call for  verb forms without the generic morpheme, which, in
the Czech examples below are primary imperfective verbs.
\ea
\ea  \label{earthva}
\gll Země se	 \minsp{\{} točí /  \minsp{\#} točívá\}	kolem	slunce.	 \\
 earth 	\textsc{refl} {}	revolve.\textsc{3sg.prs}  {} {}   revolve.\textsc{gen.3sg.prs}	around	sun\\
\glt %(\textit{Intended}:)
(Intended:) `The Earth revolves around the Sun.'
\ex \label{waltz}
\gll Valčík  \minsp{\{} je /  \minsp{{\#}} bývá\} ve tříčtvrtečním taktu.	\\
 waltz {} be.\textsc{3sg.prs} {} {}    be.\textsc{gen.3sg.prs}	in three.quarter time \\
\glt %(\textit{Intended}:)
(Intended:) `A waltz is in three quarter time.'
\z
\z
%\ea  \label{earthva}
%\gll Země se		\minsp{\{} točí \(\)/ \minsp{\#} točívá\}	kolem	slunce.	 \\
% earth 	\textsc{refl} {}	revolve.\textsc{3sg.prs}  {} \(\)/  revolve.\textsc{gen.3sg.prs}	around	sun\\
%\glt %(\textit{Intended}:)
%`The Earth revolves / \minsp{\#} tends to revolve around the Sun.'
%\ex \label{waltz}
%\gll Valčík  \minsp{\{} je \(\)/
%\minsp{\#}bývá\} ve tříčtvrtečním taktu.	\\
% waltz {} be.\textsc{3sg.prs}  {}\(\)/   be.\textsc{gen.3sg.prs}	in three.quarter time \\
%\glt %(\textit{Intended}:)
%`A waltz is / \minsp{\#} tends to be in three quarter time.'
%\z
%\z

\noindent Semantic Property 2 is also
clearly manifested in the incompatibility of
the generic morpheme \textit{-va-} with an overt universal quantifier in the same clause:
%\ea
%\ea[\#] {\label{every.Sat.gen}
%\gll \minsp{\{} Každou 	sobotu 	/ vždycky \} 	sedává  	Honza	v hospodě.	\\
%{} each Saturday {} always {} 	sit.\textsc{gen.3sg.prs} Honza in pub\\}
%\glt \textit{Intended}: `Every Saturday / always  Honza usually tends to sit in a pub.'
%\ex {\label{every.Sat}
%\gll \minsp{\{} Každou 	sobotu / vždycky \} 	sedí Honza	v hospodě.	\\
%{} each Saturday {} always {} 	sit.\textsc{3sg.prs} Honza in pub\\}
%\glt `Every Saturday / always Honza sits in a pub.'
%\z
%\z

\ea
\ea{\label{every.Sat.gen}
\gll \minsp{\{{\#}} Každou 	sobotu 	/ \minsp{\#} vždycky\}  	sedává  	Honza	v hospodě.	\\
{} each Saturday {} {}  always  	sit.\textsc{gen.3sg.prs} Honza in pub\\
\glt `\{\#Every Saturday / \#always\}  Honza usually tends to sit in a pub.'}
\ex {\label{every.Sat}
\gll \minsp{\{} Každou 	sobotu / vždycky\} 	sedí Honza	v hospodě.	\\
{} each Saturday {} always sit.\textsc{3sg.prs} Honza in pub\\
\glt `\{Every Saturday / always\} Honza sits in a pub.'}
\z
\z

\noindent
(\ref{every.Sat.gen}) is at least odd, if not
involving a contradiction (but see below analogous examples (\ref{capek}) and (\ref{suplik}), where there is no oddity), assuming that the universal Q-adverbs, \textit{každou sobotu} `every Saturday' and \textit{vždycky} `always', are used with their universal quantificational force.
If so, they require that
the sentence be true just in case Honza's sitting in a pub occurs on each and every Saturday with no exception, or on each appropriate contextually determined occasion. However, the generic morpheme \textit{-va-} contributes the uncancellable inference that
this precisely is \textit{not} necessarily the case.
That the oddity of (\ref{every.Sat.gen}) must be due to the generic morpheme \textit{-va-} is evidenced by
the fact that (\ref{every.Sat}) is perfectly acceptable and true, because it minimally differs from  (\ref{every.Sat.gen}) by lacking \textit{-va-}.

However, the generic morpheme \textit{-va-} may co-occur with universal Q-adverbs, if they are
used for ``intensification of the strength''
of the expressed regularity  (\citealt [45]{danaher2003semantics}),
rather than with their inherent universal quantificational force and so allowing no exceptions. The examples below are taken from  \citet [45]{danaher2003semantics}:
% The generic verbs and the Q-adverbs are in bold-face, glosses are omitted for the sake of brevity and simplicity:

\ea
\ea \label{capek} \gll Mládež ve Vídni se zabývala Hebblem -- já jsem \textit{vždycky} \textit{býval} skeptický k takovým módním proudům.\\
youth in Vienna \textsc{refl} concerned Hebbel {} I \textsc{aux} always was.\textsc{gen} skeptical to such fashion streams\\
\glt `Viennese youth were all reading Hebbel -- I was usually always skeptical about these fashionable influences.'\\\hfill (from K. Čapek: \textit{Hovory s TGM})
% \citep[p.~57]{capek1990hovory}%Čapek 1990, p.57\\
\ex	\label{suplik} \gll \minsp{„} Je to divný,“ pokračovala pak rychlým a věcným šepotem, \minsp{„} jeden šuplík má zamčenej, a \textit{nikdy} ho \textit{nemívá} zamčenej. A nepasuje mi do něj žádnej klíč.“\\
{} is it strange continued then quick and matter.of.fact whisper {} one drawer has locked and never.\textsc{nci} it \textsc{neg}.has.\textsc{gen} locked and \textsc{neg}.fit me into it no.\textsc{nci} key\\
\\
\glt `\.{“}It’s strange,'' she continued in a quick and matter-of-fact whisper, ``one of his desk drawers is locked and he never has it locked. And none of my keys fit the lock.''\.'\hfill (from H. Bělohradská: \textit{Poslední večeře})
% \citep{belohradska1992posledni}
\z
\z

\noindent The felicitous co-occurrence of
the  generic morpheme \textit{-va-} with universal Q-adverbs
is similar to the felicitous use of non-universal Q-adverbs together with universal Q-adverbs in combinations like \textit{usually always} or \textit{usually never} in English, where the universal Q-adverbs also arguably lack their inherent universal quantificational force,
and function as intensifiers of a regularity that is known to have exceptions:

\ea
\ea I am usually always happy, but today I feel really depressed.
\ex I am usually never neurotic about being messy and keeping things tidy, but I can't seem to go to sleep if clothes are hanging up to dry in my room. \hfill (Barbara Partee, p.c.)
\z
\z
%As we have just seen, the Semantic Property 2 in (\ref{ex:semprop2}) predicts that a generic verb formally marked with the generic morpheme  cannot be used for  the encoding of universal generalizations and definitional statements that necessarily hold for all the members of a given kind.

\noindent Semantic Property 2 in (\ref{ex:semprop2}) also predicts
that the generic morpheme \textit{-va-} will be unacceptable, or at least strongly dispreferred, for the encoding of properties that are
denoted by inherent individual-level predicates like \textit{know (French), believe, have long arms, be intelligent}
that are stable for an individual, in so far as they spread over a relatively large interval of that individual's life-time,
and possibly over most/all of its life-time, and including all of its subintervals and moments.  One way of thinking about this property is that
individual-level predicates involve a kind of implicit universal quantification over intervals and moments at which they are true of individuals.
%Take \textit{know French}, for instance.
%If you know French, your competence spreads over all the moments of your existence, with no gaps or interruptions (even if your actual reading/writing performance may fluctuate).
This motivates why the  generic morpheme \textit{-va-} cannot (freely
and straightforwardly) apply to imperfectives that denote individual-level predicates: precisely because
it introduces the uncancellable inference that there are or may be exceptions or counterexamples, i.e.,
that the characterizing property does not stably hold of an individual at all
the moments of its life-time.
This is illustrated by the following Czech example:

\ea{\label {know.va}
\gll Je rodilá Francouzska, a proto \minsp{\#} znává francouzsky.\\
is born French.woman  and therefore {} know.\textsc{gen.3sg} French\\}
\glt (Intended:) `She is a French native, and so knows French.'
\z

%Just as the Czech verb `know (French)' is odd with the generic morpheme, as in (\ref{know.va})
\noindent Similarly, \textit{know (French)} in English is odd
with a Q-adverb like \textit{sometimes}, as we see in (\ref{know.FrenchE}).
In contrast, \textit{sometimes} is compatible with the individual-level predicate \textit{be a California resident}, as it can be construed as describing
a property holding with interruptions,
on and off at different subintervals
of evaluation, sometimes it holds, and sometimes it does not.
This fluctuation in truth values generates ``gaps'' or ``interruptions''
in their application and hence also a plurality of separate cases,
which is a prerequisite for the application of the Q-adverb (see e.g. \citealt{fernald2000} and related suggestions by \citealt{deswart1991adverbs}, i.a.):

\ea
\ea[\#]{\label{know.FrenchE} Mary sometimes knows French.}
\ex[]{\label{maxcalifornia} Max is sometimes a California resident. \hfill \citep{fernald2000}}
\z \z

\noindent The possibility of an ``interruption'' construal
also makes the combination of the  generic morpheme \textit{-va-} with individual-level predicates acceptable,
and so Semantic Property 2 in (\ref{ex:semprop2}) is satisfied.
Some Czech examples are given below, where
(\ref{karelva1}) is taken from
\textit{Seven poems written on a mirror} by Karel Kryl (1976):

\ea\label{ex:genILP}
\ea myslit (si) `to think' $\rightarrow$ myslívat (si) `(to tend) to think' (see (\ref{thinkva1}))
\ex patřit `to belong (to)' $\rightarrow$  patřívat	 `(to tend) to belong to' (see (\ref{karelva1}))
\ex věřit `to believe'  $\rightarrow$ věřívat	`(to tend) to believe'
\ex mít	`to have' $\rightarrow$ mívat  `(to tend) to have'
\z
\z

\iffalse
\noindent
\textit{Figure 4}: The generic morpheme \textit{-va-} applied to \textsc{ipfv} verb stems denoting \textsc{ilp}s\\

\begin{tabular}{ p{3.8cm} p{0.5cm} p{9cm}  }
\textbf{\textsc{ipf: episodic/generic}} & & \textbf{\textsc{generic}} \\
myslit (si) `to think' & $\rightarrow$ & myslívat (si) `(to tend) to think' (see (\ref{thinkva1})), \\
patřit `to belong (to)' & $\rightarrow$ & patřívat	 `(to tend) to belong to' (see (\ref{karelva1}))	\\
věřit `to believe'		 & $\rightarrow$ &	 věřívat	`(to tend) to believe'	 \\
mít	`to have'			 & $\rightarrow$ &	 mívat  `(to tend) to have'		\\
\vspace{0.5cm}
\end{tabular}
\fi
\ea
\ea  {\label{thinkva1}
\gll Každopádně 		nejste 		sám, 	kdo 	si 		to 	\textit{někdy} 		\textit{myslívá}.\\
 in.any.case	\textsc{neg}.be.\textsc{2pl}	alone who \textsc{refl} it 	sometimes 	think.\textsc{gen.3sg} \\}
\glt `In any case, you are not the only one who occasionaly thinks like this.’
\ex {\label{karelva1}
\gll {\dots} \textit{občas} 		i 	k 	slovesu milovat	\textit{patřívá}  zájmeno zvratné.\\
{} at.times even to verb love belong.\textsc{gen.3sg} pronoun reflexive \\
\glt `{\dots} the reflexive pronoun occasionally belongs even to the verb \textit{love}'}
\z
\z

\noindent The ``interruption'' construal of individual-level predicates,  which facilitates
their use in the scope of non-universal Q-adverbs in generic sentences,
might look like a kind of coercion of individual-level predicates into a stage-level predicate interpretation.
However, if this were a matter of coerced stage-level predicate interpretations,
then they should be also available for individual-level predicates in episodic contexts,
for instance, triggered by adverbials denoting particular spatio-temporal locations. However, such coerced stage-level predicate interpretations of individual-level predicates are not available:

\ea\#Max is a California resident right now in his garden.
\z

\noindent Such considerations lead \citet{carlson1977unified}, \citet{kratzer1989stage}, and \citet{fernald2000} to propose that individual-level predicates in generic sentences like (\ref{maxcalifornia})
are not to be treated as having first undergone a shift into a stage-level predicate interpretation, but rather still have arguments that are of the individual sort, and denote individual-level stative properties,
albeit having truth values that vary across evaluation times, are true on and off.  \textit{Mutatis mutandis} the same holds for individual-level predicates in the scope of the
generic morpheme \textit{-va-}.

Another manifestation of Semantic Property 2 in (\ref{ex:semprop2})
is the oddity, or even unacceptability, of the use of
the  generic morpheme \textit{-va-}
in generic contexts in which what is at stake are
characterizing generic properties, such as
inherent dispositions, abilities, and acquired professions, that are understood as holding without interruptions for
a given individual over an extended period of that individual's life-time,
and all its subintervals and moments. Consider the following contrast:
%between a generic verb with the generic morpheme \textit{-va-} and its corresponding counterpart without \textit{-va-} under its generic interpretation:

\ea \textit{Question}:  %Jaké má Marie povolání? -
`What is Mary's profession?'
\ea[]{\label{teacher}
\gll  Učí 		na střední škole. \\
teach.\textsc{3sg.prs} 	on middle school \\
\glt `She teaches at high school.'  \\
Inference: `Mary is a high-school teacher.'}
\ex[\#]{\label{teaches}
\gll Učívá           		na střední škole.  \\
teach.\textsc{gen.3sg.prs}      on middle school	\\
\glt Intended: `She teaches at high school on and off /  sometimes / often\ldots'\\
Inference: `Being a high-school teacher is not Mary's profession.'}
\z
\z

\noindent
The question about Mary's profession is best answered with (\ref{teacher}), which is headed by the imperfective non-generic verb \textit{učí} `she teaches', here under its generic interpretation, and is paraphrasable with `Mary is a high-school teacher.' Teaching as one's profession
typically presupposes having acquired the requisite training
and accreditation, which then holds for an extended interval of one's life,
and all its subintervals, without interruptions, including times off work (weekends, holidays, or absence due to illness).

In contrast, (\ref{teaches}) which is formally marked with the generic \textit{-va-}, would be odd or infelicitous in this context, because it conveys the inference that
teaching only holds with interruptions of Mary. This invites the inference that
Mary only teaches on and off, because she
may be a temporary or a substitute teacher, and may also have other alternative jobs or occupations. In short, it invites the inference that Mary is not a teacher by profession.


\subsubsubsection{Generalizations with unrealized particulars in the actual world}


The generic morpheme \textit{-va-} cannot be used
in generic sentences that concern unrealized rules, dispositions, functions, instructions and the like.\footnote{In the studies on genericity and habituality this meaning component is commonly (though not uncontroversially) referred to as the ``actuality entailment", which was originally coined by \citet{bhatt1999covert} to describe the implicative inference that arises with ability modals in perfective aspect. However, in the context of
habitual/generic sentences it is perhaps not appropriate, because the requirement that there be actual instantiations of the generalization tends to survive in negated habitual/generic sentences. Therefore, henceforth, I will use the term \textit{realization inference}.}
Consider the following example from Czech:

\eanoraggedright\label{crushoranges}
\textit{Scenario}: We have just bought a new machine for crushing oranges, but
we have not yet used it, and it is still packed in its original box.\smallskip\\
\gll Tento stroj   drtí  / \minsp{\#} drtívá  	pomeranče.\\
this machine  crush.\textsc{3sg.prs} {} {} crush.\textsc{gen.3sg.prs}	oranges\\
\glt (Intended:) `This machine crushes oranges.'
\z

\noindent
What is at stake under this scenario is the intended function
of the specific machine token in the domain of discourse, which is as yet not realized.
Here, only a verb that sanctions a purely intensional (dispositional) interpretation can be used. In
Slavic languages, it is a non-generic verb,
which, in the example above, is the primary imperfective verb \textit{drtí} `it crushes'.
The generic verb marked with the generic morpheme makes the whole sentence false under the given scenario. The reason for it is that it
requires that this particular machine have already been put to use, and so
there have been (observed) instances of it crushing of oranges in the actual world from which the expressed generalization is inductively projected.
We may formulate this realization constraint on the use of the Slavic generic morpheme as follows:

\eanoraggedright\label{ex:semprop3} \textsc{Semantic Property 3}: The generic morpheme \textit{-va}- carries the uncancellable inference (as part of its truth-conditional content) that there are verifying instances in the actual world that count as evidence for the truth of generic sentences it marks.
\z

\noindent
The realization inference that is often taken to be a signature property of so-called ``habitual'' sentences is also mistakenly correlated with their lack of modality (intensionality), and used as an argument for not treating ``habitual'' sentences as a special case of generic sentences which are inherently modal, as all agree.
However, as many semanticists also agree, any sentence that
describes a habit, regularity, pattern,  generalization of any kind
has a modal (intensional) import (see also \sectref{CharGenericity} above, and references therein). This is because its meaning transcends what we observe in the actual world; they are not only about what was realized, but also predict what is possible or likely.
By virtue of having this predictive, law-like force, habitual sentences also have
implications for other accessible alternate worlds, rather than just the actual world, and perhaps all possible worlds. Consequently, so-called ``habitual'' markers also introduce quantification over alternate possible worlds.\footnote{Similar observations are also made by  \citet{bonehdoron2010} in connection
with the Modern Hebrew periphrastic ``habitual'' construction with
\textit{haya} (`was')+participle.} If the alleged lack of this modal component is the reason for not treating so-called ``habitual'' markers as falling under genericity, or as not qualifying as generic markers, then
this exclusion is in error.

We can illustrate this point with the following example from Czech:

\ea {\label{children}
\gll Děti v této škole mívají dobré známky. \\
children in this school have.\textsc{gen.3pl.prs} good grades \\}
\glt `Children in this school tend to have good grades.'
\z

\noindent The truth of (\ref{children}), which contains
the generic marker \textit{-va-}, is not just based on a mere
plurality of situations of various children having already actually received
good grades, but (\ref{children})  also
gives rise to counterfactual inferences that transcend observable real world conditions we have access to:
namely, for instance, if a child were to attend this school, that child too would most likely or probably get good grades.
The reason for this is that the truth of (\ref{children}) also
depends on there being some underlying reason(s) or cause(s) for this, such as the quality of instruction, parents' attention, their innate predispositions, etc.  While we may not know the particular reason(s) or cause(s), whatever they might be they also ground the truth of (\ref{children}), besides our observations of children in this school getting good grades.

The same considerations that motivate the intensionality of (\ref{children}), which has an overt marker of genericity, also motivate the intensionality of \textit{Children in this school have good grades} in English, which lack any overt generic marker (see e.g., \citealt{pelletier1997generics} for similar examples),
and which are associated with a logical representation
with the intensional generic operator \cnst{Gen} (in \citealt{pelletier1997generics} or \citealt{krifka1995genericity}, i.a.).  If one accepts that the English sentence
\textit{Children in this school have good grades} has an intensional (modal) import, as most do, then one should acknowledge that (\ref{children}) does, as well.

That the realization inference of generic and habitual sentences, including those that have overt so-called ``habitual'' markers
(such as we find in Tlingit or Modern Hebrew),
is compatible with the intensionality (modality) which is characteristic of
sentences expressing regularities,
should not be surprising given that there are other modal operators that introduce the realization inference (see \citealt{bhatt1999covert} and \citealt{hacquardDISS}).


\subsubsection{When \emph{-va-} \textsc{must} be used: Positive counterinstances}\label{mustbeused}

Due to its exceptive component, here captured by Semantic Property 2 in (\ref{ex:semprop2}), the addition of the generic marker \textit{-va-} to a generic sentence may reverse its truth conditions.
This also supports its treatment as part of the lexical meaning of \textit{-va-}.
%In interaction with a given generalization and context, one of the effects of this property is to delimit the counterexamples in the domain of that generalization that would falsify it, all else being equal.
A case in point is its use in generics that have counterexamples
that count as positive counterinstances (in the sense of \citealt{leslie2008}) to the generalization.
%In order to set up what is meant by positive counterinstances to the generalization, let memake make a few observations on exceptions and counterexamples to generic sentences. Universal generalizations like \textit{Water consists of two hydrogen and one oxygen atoms} (a universal law of nature) permit no exceptions in all the known conditions, hence the issue of exceptions does not arise. Any exception to an extensional universally quantified statement renders it false, e.g., \textit{Every table has four legs} is falsified by the existence of a single table that does not have four legs in the universal quantifier's domain.
Generalizations that allow for exceptions
partition the domain of their applicability into confirming cases and non-confirming cases. However, not all non-confirming cases are alike, not all render such generalizations false.
As \citet{leslie2008} shows, for a generalization to be true (or false) it matters how exactly non-confirming cases to it are judged to fail it. Consider the following contrast:

\ea
\ea \label{birds.neg.ci} Birds fly.\jambox*{\textsc{true}}
\ex \label{books.false} Books are paperback.\jambox*{\textsc{false}}
\z
\z

\noindent
(\ref{birds.neg.ci}) is an example of an exception-tolerating, non-universal generalization that is true
by virtue of the majority of birds having the \textsc{bird}-kind characterizing property of flying,
and it is not falsified by exceptional non-flying bird subkinds like penguins, ostriches or emus, or individual birds that do not fly because of some accidental
defect. In contrast, (\ref{books.false}) is false, even if the majority of
books are paperbacks.
%The exceptions that falsify (\ref{books.false}) are books that come in different formats. So why should books that are not paperbacks make (\ref{books.false}) false, but non-flying birds do not falsify (\ref{birds.neg.ci})?

As \citet{leslie2008} argues, we are likely to judge
a generic sentence to be true, if exceptions to it constitute ``negative counterinstances''.
Non-flying birds constitute negative counterinstances to the positive
property of flying characterizing the kind \textsc{bird}.
They simply lack this positive property without providing any
alternative positive property that  would be equally (subjectively) salient or striking as flying and that we could pit against it.
This makes such negative counterinstances
easy for us to ignore, or set aside, when judging
a generalization like (\ref{birds.neg.ci}) as true.


In contrast, if exceptions to a generalization are viewed as
positive counterinstances, we are likely to judge that generalization
as false.
(\ref{books.false}) is judged false, because exceptions to it, even if they are fewer than the confirming paperback members of the kind \textsc{book},
do not simply fail to be paperbacks, but possess a
positive distinctive feature, such as being hardback books.
I.e., they have an alternative positive property that is subjectively for us
just as salient, memorable, significant,
or striking as having the positive property of paperback.
Positive counterinstances
are cognitively salient in our kind-related expectations, which makes them
particularly hard for us to
ignore when judging the truth of generics like (\ref{books.false}).
Other examples of generics that are judged false
due to positive counterinstances to them are: \textit{Nurses are female} (there are also male nurses),
\textit{People have brown eyes} (there are people with blue, green, grey eye-color, for instance).

In order for a generalization with positive counterinstances to be true, it must contain some overt expression qualifying the
generalization that signals that  not all the members of the
kind share the property that can be truthfully predicated of that kind. For instance, \textit{typically} and \textit{usually} do this qualifying job, as   \citet{leslie2008} observes:

\ea \label{books.paperback.true}
Typically/Usually, books are paperback. \jambox*{\textsc{true}}
\z

\noindent
Other overt qualifiers of generic statements include
generic non-universal Q-adverbs like \textit{usually, rarely, seldom}, quantity expressions
like \textit{in the majority of cases, in some cases},
spatio-temporal modifiers like \textit{(books) on sale in this bookstore}, and also exceptives
like \textit{but}, prepositions \textit{except (for)}, and ceteris paribus clauses.

In Czech, Slovak, and also to a lesser extent in Polish, in addition to the above restrictors, we can also
use the generic morpheme \textit{-va-} to qualify a given generalization and so guarantee its truth, which otherwise
would be false.  Taking examples from Czech,
this can be illustrated by
the contrast between (\ref{books.paperback.falseCzech}), on the one hand, and (\ref{books.paperback.true.adv.nova.Czech}) and (\ref{books.paperback.true.va}), on other hand:
\ea\label{books.paperback-toplevel}
\ea {\label{books.paperback.falseCzech}	\gll Knihy	jsou\textsuperscript{IPFV} 	bro\v{z}ovan\'{e}.\\
books	be.3\textsc{pl.prs}	paperback \\}  \jambox*{\textsc{false}}
\glt `Books are paperback.'
\ex {\label{books.paperback.true.adv.nova.Czech}
\gll Knihy jsou\textsuperscript{IPFV}   \minsp{\{} oby\v{c}ejn\v{e} / nej\v{c}ast\v{e}ji\} 	bro\v{z}ovan\'{e}.\\
books	be.\textsc{3\textsc{pl.prs}} {} usually {} most.often	paperback \\} \jambox*{\textsc{true}}
\glt `Books are \{usually / most often\} paperback.'
\ex {\label{books.paperback.true.va}
\gll  Knihy	b\'{y}vaj\'{i}	bro\v{z}ovan\'{e}.	\\
books	be.\textsc{gen}.3\textsc{pl.prs}	paperback \\} \jambox*{\textsc{true}}
\glt `Books tend to be paperback.'
\ex {\label{books.paperback.true.adv.Czech}
\gll  Knihy b\'{y}vaj\'{i} \minsp{\{} oby\v{c}ejn\v{e} / nej\v{c}ast\v{e}ji\} 	bro\v{z}ovan\'{e}.		\\
books be.\textsc{gen}.3\textsc{pl.prs} {} usually {} most.often paperback\\} \jambox*{\textsc{true}}
\glt `Books are \{usually / most often\} paperback.'
\z
\z
\noindent
(\ref{books.paperback.falseCzech}) contains
the primary imperfective verb \textit{jsou} meaning `they are', and it is judged false based on the same reasoning about exceptions qua positive counterinstances proposed by \citet{leslie2008} for the English generic sentence (\ref{books.false}).
Adding the Q-adverbs `usually' or `most often' in (\ref{books.paperback.true.adv.nova.Czech})
or the generic morpheme \textit{-va-} in \REF{books.paperback.true.va},
reverses the truth value to true.
Q-adverbs and the generic morpheme can co-occur in the same clause, as \REF{books.paperback.true.adv.Czech} shows (and I will return to this point further below).

The contrast between (\ref{books.paperback.falseCzech})
and (\ref{books.paperback.true.va}) shows that Slavic imperfective verbs have no exceptive meaning component that could function as a
safeguard for the truth of the generalization, and the presence of \textit{-va-}
is sufficient to make a generic sentence true.
This is yet another difference between
verbs with the generic \textit{-va-} and imperfective verbs.
In (\ref{books.paperback.falseCzech}), we have a primary imperfective verb.
Just as a primary imperfective so a secondary imperfective with
the imperfectivizing suffix is insufficient to guarantee the truth of
``bare'' generics with positive counterexamples.
Generally, an exceptive meaning component is not a part of the semantics of imperfective forms in natural languages. In this respect
verbs with the generic morpheme \textit{-va-} do not neatly fit the imperfective paradigm.





\subsubsection{When \textit{-va-} \textsc{may} be used: Pragmatic inferences based on quantity} \label{maybeused}

The exceptive component of the generic
marker \textit{-va-}, as given in (\ref{ex:semprop2}),
also has pragmatic consequences in contexts when it \textit{may} be used.
Consider the following pair of sentences which only differ
in their main verbs, the generic one marked with \textit{-va-} and
the related morphologically simpler non-generic verb:

\ea\label{classes.toplevel}
\ea {\label{classes}
\gll Výuka začíná\textsuperscript{IPFV} od 8.00 a končí v poledne.\\
teaching start.3\textsc{sg.prs} from 8:00 and ends in noon\\}
\glt `Classes start at 8 a.m. and end at noon.'
\ex {\label{classes.gen}
\gll  Výuka začínává  od 8.00 a končí v poledne. \\
teaching starts.\textsc{gen.3sg.prs} from 8:00 and ends in noon \\}
\glt `Classes tend to start at 8 a.m. and end at noon.' [as a rule, usually]
\z
\z

\noindent Suppose I registered for a course at school where classes are scheduled every day, and
I ask the school administrator when classes start.
Under this scenario, either (\ref{classes}) with the non-generic verb \textit{začíná} here meaning `it starts' or (\ref{classes.gen}) with the generic verb
\textit{začínává} meaning
`it tends to start', `it usually / often / as a rule / {\dots} starts' can be used as an answer.
Suppose the administrator responds with (\ref{classes}). The
next day I arrive at school shortly before
8 a.m., but my classes on this day in fact
start only at 10 a.m., and so I
waste two hours waiting. And yet, I could not blame
the administrator for saying something blatantly false,
because generic sentences like (\ref{classes}) with a verb formally unmarked for genericity, on their generic interpretation, describe regularities that
allow for exceptions, i.e., are exception-tolerating like most generic/habitual sentences (see \sectref{CharGenericity}).
Indeed, (\ref{classes}) can be continued without contradiction with  \textit{{\dots} ale ne vždycky} `{\dots} but not always' or
\textit{{\dots} s výjimkou pátku} `{\dots} except for Friday'.

While I could not blame
the administrator for saying something blatantly false by using (\ref{classes}), I could blame them for being misleading at least.
Here is why.  (\ref{classes}) is also compatible with no exceptions, as many generic sentences are (see \sectref{CharGenericity}).
Given that (\ref{classes}) can be used to express an exceptionless generalization, i.e.,
the strongest generalization,
from the point of view of the generalization strength, (\ref{classes})
is the stronger alternative than
(\ref{classes.gen}), which requires that there be actual or possible exceptions (Semantic Property 2).
From the administrator's choice of (\ref{classes}),
the stronger alternative,
when the weaker alternative (\ref{classes.gen})
with the generic verb could have also been used, I may infer the denial/rejection of the weaker alternative that requires that
there be actual or possible exceptions.
That is, (\ref{classes}) can be understood as involving pragmatic strengthening to an exceptionless regularity, namely that classes always start at 8 a.m., without exception.
%This follows, assuming that from the choice of the stronger alternative we may infer the denial of the weaker one. (I will return to this  ranking idea in more detail in \sectref{proposal}.)

In contrast, in the described scenario,
if the administrator utters (\ref{classes.gen}),
I cannot not blame them
for being misleading, or even
giving me wrong information about when classes start. The reason for this is that the generic marker, besides its generic import, requires that there be actual or possible exceptions to the generalization, which excludes
pragmatic strengthening of (\ref{classes.gen}) to an exceptionless generalization. Moreover, by using (\ref{classes.gen}) instead of
(\ref{classes}), the administrator
implicates that they do not have enough information to be in the position to commit to the stronger alternative (\ref{classes}), which allows a pragmatic strengthening to an exceptionless generalization.

Generally, the choice between the two lexical
lexical alternatives, a generic verb with \textit{-va-} and a corresponding morphologically related non-generic verb without \textit{-va-},
gives rise to one of the two
(defeasible) pragmatic quantity-based inferences which concern either the speaker's \textit{certainty} or \textit{ignorance}
about the existence or possibility of exceptions to it.
In our example, the administrator may use the alternative with the generic morpheme, because they are \textit{certain} that
classes do not always start at 8 a.m.,
or that they are \textit{uncertain} whether exceptions to their starting at 8 a.m.
can be completely excluded.
In the latter case, \textit{-va-} may be used as a hedge to safeguard a generalization against refutation by states of affairs
of which the speaker is ignorant, and so to ensure its plausible deniability.

This mechanism of quantity-based inferences is also exploited in legal discourse,
and in the formulation of various rules, instructions, job descriptions, or customs, which includes generics that have a normative character.
Below are some attested examples taken from Czech,
a game rule (\ref{marbles.va}) and an institutional regulation (\ref{justinian}):\footnote{Example \REF{marbles.va} is from \url{https://cs.wikipedia.org/wiki/Kuličky}, accessed on April 9, 2022. Example \REF{justinian} is from \url{https://is.muni.cz/th/nmkho/Frydek_-_Crimen_maiestatis_-_final.pdf};
accessed on August 16, 2025.}
\ea
\ea {\label{marbles.va}
\gll Zápas \textit{se hrává} na více vítězných her. \\
 match \textsc{refl} play.\textsc{gen.3sg.prs} on more games\\}
\glt `The match tends to be played in several winning sets.'
\ex \label{justinian}
\gll {[\dots]} z komentáře k Iuliovu zákonu de maiestate z Digest císaře Justiniána: [\dots] 4) kdo \textit{svolávává} k {povstání, {[\dots]}}  \\
{} from commentary to Iulius law de maiestate from disgest emperor Justinian: {} 4) who convene.\textsc{gen.3\textsc{sg.prs}} to insurrection\\
\glt `[\dots] from the commentary to Iulius' law de maiestate from the Digest of the Emperor Justinian: [\dots] who tends to call for an insurrection [\dots]'
\z
\z




%\ex {\label{board.va}
%\gll Členské schůze svolávává	výbor aspoň jednou měsíčně.\\
%member meetings convene.\textsc{gen.3\textsc{sg.prs}}  board at.least once monthly\\}
%\glt `The board tends to convene member meetings at least once a month.'

\noindent For such
rules and regulations, using a verb with the generic morpheme is useful as a hedging device.
Due to its incompatibility with exceptionless
generalizations, the generic morpheme explicitly
assures a certain leeway in the application of a rule of law, and hedges against unprecedented or unexpected exceptions to it.
In contrast, rules and regulations, such as those above, if they were formulated with morphologically related non-generic verbs (without the generic morpheme \textit{-va-}),
would be merely exception tolerating; unless the context explicitly indicated otherwise, they would
likely be
pragmatically strengthened to strict, exceptionless rules.

The use of the generic marker \textit{-va-} in
generics describing various rules and norms fits
\citeauthor{carlson1995truth}'s (\citeyear{carlson1995truth}) rules-and-regulations model for the truth conditions of generic sentences (see \sectref{CharGenericity} above). Recall that according to this model, generic sentences are true by virtue of some underlying rule, principle or cause, which we may directly learn, rather than (just) by virtue of real world particulars from which we induce the generalization.
The distribution of the generic marker \textit{-va-} cuts across the
divide between the inductive and rules-and-regulations model, and also across the descriptive/normative distinction. It clearly is not a marker of only inductive generalizations, including habits proper, which are paradigm examples for the inductive model (\citealt{carlson1995truth}).
However, the rules and regulations that are expressed by generic sentences formally marked with \textit{-va-} must have been already instantiated, which follows from its realization
requirement (see Semantic Property 1 in (\ref{ex:semprop1})).


\subsection{The quantificational properties of the generic morpheme\textit{-va-}} \label{quantprop}

\subsubsection{Quantification over situations and/or individuals}\label{quantprop-sit-ind}


One common way of representing the logical structure of
generic sentences was introduced in \sectref{CharGenericity}:
namely, the tripartite structure with \cnst{Gen}, which is a phonologically null, modalized dyadic operator.  The tripartite
structure was originally proposed for the logical structure of
sentences with overt Q-adverbs \citep{lewis1975}.
As \citet{chierchia1995} observes, \cnst{Gen} patterns with overt Q-adverbs when it comes to its variable-binding properties:

\ea \label{varbindingGen} Variable-binding properties of the  generic operator \cnst{Gen} (see e.g., \citealt{chierchia1995}):
\ea a situation variable (in generic sentences without kind terms)
\ex variables provided by singular indefinites and bare plurals
\ex variables provided by kind-denoting definites
\ex more than one variable (in generic sentences with kind terms)
\ex	 relatively freely select the arguments for its binding, modulo context
\z
\z

\noindent
The generic marker \textit{-va-} appears to
share its variable binding properties with the generic operator \cnst{Gen} (see also \citealt{filip2009parisva,filip2021cologneva}),
which serves as one among other arguments given here in support
of its status as a generic marker.
To illustrate this point, I will adduce some examples  from Czech.
Comparable uses of this morpheme in other Slavic languages that
have analogous verbs with the generic marker will exhibit the same binding properties, as can be easily verified.

\largerpage[2]
The representational style of the simplified tripartite structures that are given below follows \citet{krifka1995genericity} and \citet{pelletier1997generics} (see also above in \sectref{CharGenericity}). The contribution of the generic morpheme \textit{-va-} is
represented in the following tripartite structures
by means of the quantifier $\cnst{va}$, a kind of generic quantifier.\footnote{The treatment of the proper name \textit{Tom} in the formula (\ref{tombeer.f}) follows the treatment of the proper name \textit{Fred} in
\textit{Fred always smokes} in \citet [189, ex. (37)]{chierchia1995}.  It is treated as a constant in the logical representation of \textit{Fred always smokes}; the universal Q-adverb \textit{always} only binds the situation variable; its
variable binding properties, according to Chierchia and others,
pattern with \cnst{Gen}. However, strictly speaking, the generic operator \cnst{Gen} in this type of generic
sentence quantifies over different stages of the name bearer
participating in different corresponding particular situations in which
the characterizing property is manifested. This is implied in an alternative tripartite structure associated with this type of sentence in
\citet{krifka1995genericity}, where \textit{Tom} introduces the individual variable \textit{x}, which is satisfied by different stages of Tom at particular situations: \textit{Tom drinks beer after work} $\rightarrow$ $\cnst{Gen}[x,s] (x = \textsc{tom}  \wedge x\, \textrm{in}\, s   \wedge s\, \textsc{after work};
x \, \textsc{drink\_beer}\, \textrm{in}\, s)$.}$^,$\footnote{Slavic languages lack an article system; only South Slavic languages like Bulgarian and Macedonian have
what is taken to be a definite article. The bare subject argument in the following Czech examples can be viewed as corresponding to the English singular definite subject, bare plural subject, and where the predicate does not denote a contingent property also to the English indefinite subject argument.}
%The focus of the simplified logical structures below is on the binding properties of the Slavic generic morpheme, abstracting away from other meaning components of  the example sentences for the sake of simplicity.

\ea \textsc{Binding of situation variables} in generics without kind terms
\ea {\label{tombeer}
\gll Po práci	   si 	   Tom   dávává           pivo.	\\
after work   \textsc{refl}  Tom   give.\textsc{gen.3sg.prs} beer\\}
\glt `After work, Tom has (a) beer.' [i.e., now and then / rarely / often \ldots]
\ex {\label{tombeer.f} $\cnst{va}[s;] (\textsc{Tom} \, \textrm{in} \, s \wedge s\, \textsc{after work}; \textsc{Tom}\,  \textsc{has (a) beer}\, \textrm{in}\, s)$	}\\
`For all appropriate after-work situations \textit{s} such that Tom is in \textit{s}, Tom has (a) beer in \textit{s}.’
\z
\z

\ea \textsc{Binding of variables provided by indefinites}
\ea {\label{chair}
\gll Židle   	\minsp{\{} mívá / mívají\} 		čtyři nohy.\\
chair(s) {} have.\textsc{gen.3sg.prs} {} have.\textsc{gen.3pl.prs}	four legs\\}
\glt `A typical chair/Chairs tend(s) to have four legs.’
\ex $\cnst{va} [x;] (\textsc{chair} (x) ; \textsc{have four legs} (x))$\\
`When a thing has the property of being a chair, it has four legs.’ \\
`For a given thing $x$ such that $x$ is a chair, $x$ has four legs.’
\z
\z

\ea \textsc{Binding of variables provided by kind-denoting terms}
\ea {\label{man.change}
\gll Člověk	se 		k stáru měnívá.\\
man \textsc{refl}	toward old.age	change.\textsc{gen.3sg.prs}\\}
\glt `Man changes as he grows old.'	 \jambox*{(from K. Čapek: \textit{Obyčejný život})}
\ex $\cnst{va}[x,s;] (\textsc{man}(x) \wedge x\, \textrm{in}\, s \wedge s\, \textsc{old age}; x\, \textsc{change}\, \textrm{in}\, s)$	\\
\z
\z

\newpage
\ea \textsc{Binding of more than one individual variable}
\ea {\label{catmouse}
\gll Kočka 	honívá   myš.\\
cat		chase.\textsc{gen.3sg.prs}   	  mouse \\}
\glt `Cats tend to chase mice.’
\ex $\cnst{va}[s,x,y;] (\textsc{cat}(x) \wedge \textsc{mouse}(y) \wedge x\, \textrm{in}\, s  \wedge y\, \textrm{in}\, s ;
x\, \textsc{chase}\ y\ \textrm{in}\, s)$	\\
\z
\z


\noindent
As we see above, the Czech marker \textit{-va-},
which introduces the quantifier $\cnst{va}$ into  the tripartite structure, semantically functions as a quantifier
over (i) situations, as in (\ref{tombeer}), (ii) individuals, as in (\ref{chair}), and (iii) both situations and individuals in sentences with kind denoting terms, as in (\ref{man.change}) and
(\ref{catmouse}),
which is aligned
with the quantificational properties of the null generic operator \cnst{Gen}.

In (\ref{chair}), the Czech marker
\textit{-va-} functions
as a quantifier whose range only contains individual
members of a kind.  This  invalidates its treatment as an iterative marker
denoting a function that returns a set of plural dynamic
situations of the same type, at a given world and time.
Examples like (\ref{chair}) clearly show that the standard
designations of \textit{-va-} as an iterative, multiplicative, or frequentative marker (e.g., \citealt{petr1986mluvnice}, \citealt{Kosek2014}, \citealt{lamprecht1986historicka}, \citealt{dahl1995}) are misleading, if not wrong.

Moreover, examples like (\ref{chair}) invalidate the claim made by  \citet[421]{dahl1995}
that the Czech marker \textit{-va-}  serves ``as a kind of quantifier over situations with, roughly, the semantics of `most'. Following the terminology of \cite{Dahl1985}, let us refer to this kind of tense-aspect markers as habituals.'' Given that this claim
is false, and given that it is one of the two key claims that Dahl uses
to deny that the Czech morpheme \textit{-va-} has a status of a generic marker, there really remains no valid reason, which Dahl adduces, why \textit{-va-} cannot be analyzed as
a generic marker sui generis, given that the other argument, namely
the obligatory occurrence in all and only generic sentences is also
spurious.
Given that Dahl uses the Czech marker \textit{-va-}
as a paradigm example of comparable markers  in other typologically distinct languages (see \citealt[421,~fn.~8]{dahl1995}) that enforce
a generic (and habitual) interpretation, this also raises at least some doubts about their purported status as ``tense-aspect habitual markers''.

The use of the generic morpheme \textit{-va-} in
examples like (\ref{chair}) also shows that it is clearly not a ``habitual'' morpheme in the proper sense of ``habit'', i.e., dedicated to formally marking only sentences that quantify over situations
with animate individuals, typically human agents, as in (\ref{tombeer}).
Another example of a generic sentence
marked with \textit{-va-} that
has nothing to do with habits in the narrow sense is given below:
\ea {\label{sea.va}
\gll Moře tady b\'{y}v\'{a} v	tuto dobu	vyhřáté	na příjemných 28° C. \\
sea here	be.\textsc{gen.prs}	in	this time		warmed-up	on pleasant 28° C\\}
\glt `At this time the sea here tends to be warmed up to a pleasant 28° C.'
\z

\noindent To illustrate the breadth of use of the Czech generic morpheme, we may also mention that it commonly occurs in generic
sentences with dedicated kind-denoting terms like `man', as in (\ref{man.change}).
It also occurs in kind reference sentences (see \sectref{CharGenericity}) that contain kind-level predicates like \textit{be widespread} directly selecting kind-denoting terms for one of their arguments and which assign characterizing properties to kinds
that can never hold of its individual members, as we see
in (\ref{moss}) below:

\ea  {\label{moss}
\gll Rohozub nachový    bývá    rozšířený      u lidských sídlišť.\\
ceratodon purpureus    be.\textsc{gen.3sg.prs}  widespread    at human dwellings\\}
\glt `The ceratodon purpureus tends to be widespread close to human dwellings.’
%\ex $\textsc{va} [x;] (x=\textsc{ceratodon purpureus}; \textsc{widespread at human dwellings} (x))$
%\z
\z

\noindent There is a growing body of studies on
markers that, just like the Slavic generic marker, constitute
a sufficient, but not a necessary, condition
for the expression of various regularities, and often, similarly
as in \citet{dahl1995}, they are
treated as ``habitual'', rather than generic, markers.
Some examples are:
the Tlingit (Na-Dene; Alaska, British Columbia, Yukon) particle \textit{nooch} \citep{cable2022}, the \textit{be}-construction in African American English \citep{green2000aspectual}, or the construction with \textit{haya}+parti\-cip\-le in Modern Hebrew \citep{boneh2008habituality}.
It is striking  that such studies focus on sentences that describe habits in the proper sense, such as (\ref{tombeer}).
%, i.e., they describe regularities of action of animate, mostly, human agents.
It then raises the question whether this empirical focus of studies
on such ``habitual'' markers and on habits
proper is due to a data elicitation bias, and whether
the relevant markers in these languages actually cover
a broader swath of data, similar to those examples that we find in
Czech.  One wonders whether their quantificational properties come close to or perhaps even match those of the null generic operator \cnst{Gen}, or  the Czech generic marker \textit{-va-}, and by the same token,
whether their treatment as merely ``habitual" markers (possibly
subsumed under tense and/or aspect) which are claimed to be separate from genericity might not be reconsidered.

\subsubsection{The Slavic generic morpheme \textit{-va-} is not an overt Q-adverb: Irreducibility to `usually' or to any other single quantifier or quantity expression}\label{noqadv}

In standard reference grammars of Czech, the meaning
of the Czech morpheme \textit{-va-} is often characterized in terms of  ``usuality'' (\textit{uzuálnost} in Czech, see e.g., \citealt[184, i.a.]{petr1986mluvnice}). Hence, \citeauthor{dahl1995}'s (\citeyear{dahl1995}) claim that the Czech morpheme \textit{-va-} functions as a quantifier over situations with the meaning of approximately `most' or `usually' might seem plausible. In \citet{dahl1995}, this semantic claim is based on two rather flimsy data points:
namely, ``the glosses given in \citet{kucera1981aspect}'' (\citealt[421]{dahl1995}), and a cursory observation
(\citealt[420--421]{dahl1995}) regarding the distribution of the Czech \textit{-va-} and similar markers from other languages over his two generic contexts: namely, (i) it tends to be optional in the context of `usually', and (ii) it is typically absent in the context that concerns a job description.\footnote{In \citet[420--421]{dahl1995}, the relevant test context is given as follows: [A: My brother works at an office. B: What kind of work he DO? A:] He WRITE letters.}

While \citet{dahl1995} does not motivate the choice of these two generic contexts, they seem well-chosen.
The first would seem to elicit generics that  are roughly aligned with \citeauthor{carlson1995truth}'s (\citeyear{carlson1995truth}) inductive model of genericity, and concern weak descriptive generalizations,
while the second with \citeauthor{carlson1995truth}'s (\citeyear{carlson1995truth}) rules-and-regulations model of genericity.
However, as shown above (Section \ref{maybeused}), the Czech generic \textit{-va-} can be used not only in generic sentences that describe weak descriptive generalizations, but also in generics stating a variety of rules (e.g., game rules), including the rules of the law.

Now, if the Czech \textit{-va-} had a meaning of approximately `most' or `usually', as \citet{dahl1995} claims, it would not qualify as a generic morpheme.
% The reason is that it would then function as an extensional quantifier, specifying a mere multiplicity of situations, which does not amount to genericity and neither to habituality for that matter (\textit{pace} \citealt{dahl1995}).
However, the Czech \textit{-va-}
contributes no requirement on a specific quantity of situations or instances that count as evidence for the truth of the generalization. This is a hallmark property of  markers of genericity (and habituality as its special case) in natural languages.
In this respect, the generic morpheme \textit{-va-}
patterns with the phonologically null \cnst{Gen} operator. Consequently, just as \cnst{Gen}, it cannot be reduced to `most' or `usually', or the meaning of any other single quantifier or expression of quantity, no matter how modalized
(see \ref{CharGenericity} above).
Let me adduce a few concrete examples to illustrate this point.

First, we observe that the Czech marker \textit{-va-} freely co-occurs with virtually any expression of quantification or quantity (see, for instance, corpus studies of
\citealt[62,~81]{sirokova1963}; \citealt{sirokova1965},  \citealt{danaher2003semantics}).
\citeauthor{danaher2003semantics}'s (\citeyear{danaher2003semantics}) corpus study shows
that \textit{-va-} occurs less often with \textit{obvykle} `usually’ than with Q-adverb(ial)s denoting a low frequency like \textit{občas} `from time to time’, \textit{někdy} `sometimes’, \textit{málokdy} `rarely’, \textit{tu a tam} `here and there’, \textit{vzácně} `rarely’.\footnote{\citeauthor{danaher2003semantics}'s (\citeyear{danaher2003semantics}) corpus comprises 376 attested examples of verbs marked with the marker  \textit{–va-}, collected from a wide variety of text genres (e.g., contemporary literary essays, fiction, memoirs, journalistic prose and scholarly writing).}  Any of
these Q-adverb(ial)s can fill  the  $[$\textsc{q-adverb}$]$ slot of (\ref{thomas.breakfast}):

\ea{\label{thomas.breakfast}
\gll Po snídani Tomáš  \textsc{q-adverb} psává	 dopisy. \\
after breakfast Thomas	\textsc{q-adverb}  write.\textsc{va.prs} letters \\}
\glt `Thomas \textsc{q-adverb}  writes letters after breakfast.'
\z

\noindent Second, \textit{-va-} can be used to formally mark generic sentences
that are true even if \textit{most} relevant instances do \textit{not} satisfy the generically-predicated property, as we see in (\ref{oligarchs}), which makes its reductive analysis to `most' or `usually' entirely implausible.\footnote{This example is tailored on \citeauthor{kucera1981aspect}'s (\citeyear{kucera1981aspect}) example:
\ea
\gll Za Stalina 	        ruští  generálové     umírávali    v mladém věku.	\\
during Stalin    Russian generals    died.\textsc{gen}   	        in young age\\
\glt `In Stalin's times, Russian generals tended to die young.'
\z
\noindent Kučera translates this sentence as `Most generals died young in Stalin's times', but this does not seem to be correct, given that what it describes is factually false; most Russian generals did not in fact die young in Stalin's times, as we know.}

\ea{\label{oligarchs}
\gll Za Putina, ruští oligarchové umírávali pádem z okna.\\
during Putin Russian oligarchs die.\textsc{gen.3pl.pst} fall from window \\}
\glt `In Putin's times, Russian oligarchs tended to die by falling out of windows.'
\z

\noindent
(\ref{oligarchs}) is true even if as a matter of fact most Russian oligarchs who died
in Putin's times did not die by falling out of windows, only a very small number
did. The reason why we judge this sentence as true does not hinge on the number of Russian oligarchs dying by falling out of windows, but rather on
viewing the nature of their death as so highly appalling, striking, unexpected,
or memorable that we treat it as their characterizing property. This kind of subjective evaluation
is based on the same psychological mechanism that motivates why we
judge generic sentences like (\ref{mosquitoesEnglish}) \textit{Mosquitoes carry the West Nile virus} as true,
even if only about \qty{1}{\percent} of mosquitoes carry the virus.

Moreover, adding `most'  or `usually'  to true generic sentences like  (\ref{oligarchs}) that are formally marked with \textit{-va-} results in false sentences, as we see in (\ref{oligarchs.most}) and (\ref{oligarchs.usually}). Given that a generic sentence marked with \textit{-va-} does not entail the corresponding sentence with the added quantifier
`most'  or `usually', \textit{-va-} does not have as part of its meaning
the quantificational import of `most'  or `usually'.
\ea
\ea{\label{oligarchs.most}
\gll Za Putina, \textit{většina}  oligarchů umírávala pádem z okna. \\
during Putin most oligarchs die.\textsc{gen.pst} fall from window \\}
\glt `In Putin's times, most oligarchs died by falling out of windows.'\\
(\ref{oligarchs}): \textsc{true} $\nRightarrow$ (\ref{oligarchs.most}): \textsc{false}
\ex {\label{oligarchs.usually}
\gll Za Putina,  oligarchové \textit{obyčejně} umírávali pádem z okna.\\
during Putin oligarchs usually die.\textsc{gen.pst} fall from window\\}
\glt `In Putin's times,  oligarchs usually died by falling out of windows.'\\
(\ref{oligarchs}): \textsc{true} $\nRightarrow$ (\ref{oligarchs.usually}): \textsc{false}
\z
\z
%Given the above examples, it should also be evident that the generic morpheme \textit{-va-} cannot be analyzed as a kind of overt Q-adverb, say, for instance `usually'. First, (\ref{oligarchs}), the generic sentence marked with \textit{-va-}, is true, but the corresponding sentence (\ref{oligarchs.usually}) with the added overt Q-adverb `usually' is false. Second,

\noindent Neither can (\ref{oligarchs}) with the generic marker \textit{-va-} be paraphrased by a corresponding sentence with an overt Q-adverb like `usually'
and a non-generic imperfective verb without \textit{-va-}. While
(\ref{oligarchs}) is true, (\ref{oligarchs.nova.usually}) is false:
\ea{\label{oligarchs.nova.usually}
\gll Za Putina,  oligarchové \textit{obyčejně} umírali\textsuperscript{IPFV} pádem z okna.\\
during Putin oligarchs usually die.\textsc{pst} fall from window\\}
\glt `In Putin's times,  oligarchs usually died by falling out of windows.'\\
(\ref{oligarchs}): \textsc{true} $\nRightarrow$ (\ref{oligarchs.nova.usually}): \textsc{false}
\z

\noindent In sum, the generic morpheme \textit{-va-} in Czech does not function as a ``habitual'' marker that quantifies over situations only, and its meaning cannot be reduced  to that of an overt Q-adverb like `usually',
or any other quantifier or expression of quantity,
contrary to \citet{dahl1995}. In Slavic languages that have
analogous verbs with the generic marker \textit{-va-}, we can
observe the same general interpretation pattern concerning its
quantificational import.


\subsection{Summary}

%The productivity and use of the Slavic generic morpheme \textit{-va-} varies across different Slavic languages, as observed above.  Here it is explored using the data from Czech, where, as in other West Slavic languages, it is productive. It is plausible to hypothesize that the  properties that have been here identified for the Czech \textit{-va-} also hold for the instantiations of this generic morpheme in other Slavic languages, even if their use may be rather limited, as in South and East Slavic languages.

At the outset, one of the main questions posed
concerned the relation of
the meaning of the generic morpheme \textit{-va-} to the phonologically null generic \cnst{Gen} operator, which in quantificational theories of genericity is hypothesized to stand for the key meaning
component of generic and also habitual sentences
(see (\ref{habitual.def})).
Based on the data and observations so far, and
assuming that the generic morpheme \textit{-va-} introduces a kind of generic operator \cnst{va} into the tripartite logical representation, the similarities and differences between \cnst{va} and \cnst{Gen} can be summarized as follows.

\eanoraggedright
Shared properties of
the generic operator \cnst{va} (in Czech) and the null generic operator \cnst{Gen}:
\eanoraggedright They model the meanings of sentences that only have a generic interpretation (and disallow any episodic interpretation).
\ex They derive complex generic predicates
that are aspectually stative, one of the signature properties of generic predicates (see above \sectref{CharGenericity}, and Semantic Property 1 in (\ref{ex:semprop1final})).
\ex Their analysis cannot be reduced to the meaning of any single quantity expression like \textit{usually}, \textit{generally}, \textit{typically}, \textit{most}, or \textit{in a significant number of cases}, of any extensional quantifier (e.g., universal, existential), or any single, purely quantitative criterion or quantity-based measure, regardless whether they are characterized in modal or probabilistic terms (see  \sectref{noqadv}).
\ex Both are modal generic operators;
their meaning does not (just) depend on a mere plurality of particulars, situations or ``cases'' (in the sense of \citealt{lewis1975}),
but rather they add predictive force to the meaning
of sentences, which has implications for other accessible worlds (perhaps all), than just one possible (actual) world,
and supports counterfactual inferences.
\ex They seem to exhibit the same variable binding properties (akin to those of overt Q-adverbs) (see \sectref{quantprop-sit-ind}).
\z
\z

\noindent
The generic operator \cnst{va} differs
from \cnst{Gen} in so far as it adds two uncancellable inferences (as part of its truth-conditional meaning)
to the meaning of generic sentences. For convenience,
they are repeated here as follows:

\eanoraggedright
Two uncancellable  inferences\label{two.inferences}
\eanoraggedright Exceptive inference: There are actual (known)  or possible exceptions or counterexamples to the generalization (Semantic Property 2 in (\ref{ex:semprop2})).
\ex Realization inference: There are verifying instances in the actual world that count as evidence for the truth of the generalization (Semantic Property 3 in (\ref{ex:semprop3})).
\z
\z
\noindent
The co-occurrence of these two uncancellable inferences  may not be a matter of mere accident. The generic \textit{-va-}, as I argue,
semantically functions as a modal generic quantifier sui generis whose use involves reasoning with exceptions and counterexamples; it
involves the speaker's (implicit) background understanding of what counts as salient reasons for/against the correctness of applying the generalization, in the light of exceptions or counterexamples which the speaker thinks cannot be safely ignored or tolerated, out of abundance
of caution, to ensure plausible deniability and the like. This kind of reasoning
with exceptions and counterexamples is motivated by the speaker's more or less direct experience of the facts, particulars in the actual world, of how various states of affairs are realized in the \textit{actual} world.  If so, then it should not be surprising that \textit{-va-} also carries the uncancellable inference that there are verifying instances in the actual world that count as evidence for the truth of generic sentences  (realization inference).

Due to these two uncancellable inferences, the Slavic
generic morpheme \textit{-va-} delimits only a subset of generic sentences that the null generic
operator \cnst{Gen} underpins; \cnst{Gen} is introduced into the
logical representation of generics that not only sanction exceptions, but also of generics that (necessarily) hold with no exceptions like the universal laws of nature,
and also purely hypothetical generics that are true even if there have been no instances that count as
evidence for their truth, and may never be.

As in English, the null generic operator \cnst{Gen}
may stand for a representation of the generic (and habitual)
meaning
in the logical structure of
Slavic generic (and habitual) sentences that are headed by verbs
that are \textit{not} marked for genericity, i.e.,
non-generic verbs (see \ref{fig:filip:1}):
these are either
primary imperfective verbs (formally unmarked for imperfectivity) or secondary imperfectives (which
are marked for imperfectivity with
the imperfectivizing suffix), and also perfective verbs, in particular in West Slavic languages where they are commonly used in generic (and habitual) sentences.

% original text: "Slavic languages also express generic statements by means of  verbs that are formally unmarked for genericity,  and these introduce \cnst{Gen} into the logical representation. They contain non-generic \textit{imperfective} verbs, i.e.,  all imperfectives \textit{without} the generic morpheme \textit{-va-}, either primary or secondary (whereby the latter are marked for imperfectivity with  the imperfectivizing suffix). In West Slavic languages, perfective verbs are also commonly used for the expression  of generic statements that introduce \cnst{Gen} into the logical representation.


\section{Proposal} \label{proposal}

\subsection{Reasoning with exceptions}\label{proposal5.1}

One key difference between generic sentences formally marked with \textit{-va-} and those that lack this marker has to do with the \textsc{strength} of the generalization they express which is calibrated in terms of the speaker's (public) commitment to the actual or possible exceptions or counterexamples that the generalization \textit{may, must} or \textit{must not} have.\footnote{Assuming that making a statement publicly commits the speaker to its content: A ``speaker is committed to a statement when he makes it himself, or agrees to it as made by someone else, or if he makes or agrees to other statements from which it clearly follows” (\citealt[136]{hamblin1971mathematical,hamblin1987}).}
(These three salient cases were discussed in \sectref{delimitingVA}.)

%What is here understood by \textsc{weak} and \textsc{strong} does not depend on informative strength, but rather on the strength of the generalization that the speaker calibrates based on their knowledge of
Intuitively, the strongest generalizations have no exceptions or counterexamples whatsoever in every possible world/in all the known circumstances (e.g.,
universal laws of nature
\textit{The Earth revolves around the Sun} or definitional statements like \textit{A triangle has three sides}).
Lower on this scale are
non-universal generalizations (including laws and rules
that fall under special sciences such as
psychology, biology, and economics, e.g., \textit{As the price of a good or service increases, the quantity demanded decreases, and vice versa}.)
which in principle allow for exceptions, but may be compatible with no exceptions, followed by generalizations that have (known) exceptions that we tolerate (e.g. \textit{Dogs bark}).
%or counterexamples that falsify a generic sentence, unless it contains some overt restrictor (e.g., \#\textit{Books are paperback} versus \textit{Typically/Normally books are paperback}).

For all these types of generalizations, in Czech there are
freely available non-generic (im)perfective verb forms (i.e., forms without
the generic \textit{-va-}), unless, they have exceptions that cannot be tolerated
(positive counterinstances), in which case, just as in English, some overt expression  is needed to qualify the generalization and ensure its truth.
By using a non-generic verb in a non-universal generic sentence, which in principle allows for exceptions, the speaker may remain publicly
non-committal as to whether there are exceptions to the expressed generalization, but also, in an appropriate context and the nature of the generalization,
allow for pragmatic strengthening to an exceptionless generalization.


In contrast, the use of a specifically generic verb with the generic marker \textit{-va-} commits the speaker to there being actual exceptions or counterexamples, or possible exceptions to the generalization
that the speaker cannot categorically exclude (due to the lack of sufficient information) in every possible world (see Semantic Property 2 in (\ref{ex:semprop2})). This means that the
generic morpheme \textit{-va-} can only be used to mark
generic sentences
that express non-universal generalizations, and
which exclude pragmatic strengthening to an exceptionless generalization.
It may be used as a hedge to safeguard a generalization against refutation by states of affairs  of which the speaker is ignorant, and so to ensure its plausible deniability.
Such generalizations are akin
to generalizations with exceptives like \textit{but}, prepositions \textit{except (for)},
or with a ceteris paribus clause.

The difference in the strength of the generalization expressed by generic sentences marked with \textit{-va-}, on the one hand, and of the generalization expressed
by formally unmarked generics, on the other hand, can be represented, as I propose,
in terms of
a \textit{rank order} (\citealt{Lehrer1974}; \citealt{hirschberg1985}; \citealt{Geurts_2010}; \citealt{horn1989,horn2000,horn2010,horn2017}, i.a.).
Rank orders, which have the general format of
$\langle \langle \text{strong} \ | \  \text{weak} \rangle \rangle$,
are related to, but distinct from, scales for entailment-based scalar implicature. Both rank-order based and entailment based scalar implicatures are a kind of quantity implicature. However, in
rank order based implicatures, stronger elements do not entail weaker ones, but rather we can infer their negations:

\eanoraggedright
\textsc{scale order} versus \textsc{rank order} \citep{horn2017}

%\smallskip
%\textbf{Scale}: In a scale of the form $\langle Y, X \rangle$,
%\ldots {} $Y$ \ldots \, unilaterally entails \ldots {} $X$ \ldots\\
%Examples: \\
%$\langle scalding, hot, warm \rangle$, `if it is hot, it is warm'\\
%$\langle certain, likely, possible \rangle$,
%`if it is certain, it is likely'

\smallskip
\textsc{Scale}: In a scale of the form $\langle Y, X \rangle$,
\ldots {} $Y$ \ldots \, unilaterally entails \ldots {} $X$ \ldots\\
Example: $\langle all, most, many, some \rangle$\\
\smallskip

\vspace{-0.3cm}
\textsc{Rank order}: In a rank order of the form
$\langle\langle Y$ | $X \rangle\rangle$, \ldots {} $Y$ \ldots \, unilaterally entails  \ldots {} $\neg$ $X$ \, \ldots\\
Examples:\\
$\langle\langle$\textit{married} | \textit{engaged}$\rangle\rangle$:
`if a person is married, it is false that that person is engaged (to be married)'\\
$\langle \langle$\textit{dead} | \textit{sick}$\rangle \rangle$: `if a person is dead,  it is false that that person is sick'
\z

\noindent
Unlike values within a scale, ranked elements are mutually incompatible. For instance, being married
does not entail being engaged.
Applied to generic sentences with and without the generic morpheme
\textit{-va-} in Slavic languages,
we get the following rank order:
\ea \label{rank.order}
\textsc{Rank order} (preliminary version):
$\langle\langle[\cnst{Gen}]P$ | $[\cnst{va}]P\,\rangle\rangle$\\
where, \\
$[\cnst{Gen}]P$: compatible with no exceptions\\
$[\cnst{va}]P$: incompatible with no exceptions \\
\smallskip
In a rank order $\langle \langle[\cnst{Gen}]P$  |  $[\cnst{va}]P\,\rangle \rangle$, $[\cnst{Gen}]P$ unilaterally entails $\neg([\cnst{va}]P)$, all else being equal.
\z

\noindent
The rank order in (\ref{rank.order})
is based on two main lexical alternatives
that Slavic languages have available for the expression of generic statements:
namely, (i) verbs that are formally marked for genericity with the generic morpheme \textit{-va-}, which introduces the generic operator \cnst{va} into the logical representation and here represented with \cnst{va} as part
of the weak element of the rank order,  and (ii) non-generic verbs without the generic morpheme \textit{-va-}, which have contextually determined generic uses, and occur in generic sentences that introduce the null generic operator \cnst{gen} into the logical representation, and here represented with \cnst{gen} as part of the strong element of the rank order.
The two ranked alternatives, strong and weak,
are equally informative, and do not stand in an entailment relation to each other:
namely, the strong element does not a fortiori imply the truth of the weak one, but rather from
the strong element we can infer the negation of the weak element.
Let me illustrate this with the following contrast:
%\ea{\label{busdriver}
%\gll Petra u nás řídí\textsuperscript{IPFV} autobus. Jako malá chtěla být malířkou, ale maminka jí to nakonec rozmluvila.  \\
%Petra at us drive.\textsc{3sg.prs} bus. As little want.\textsc{3sg.pst} be.\textsc{inf} painter but mom her it ultimately dissuade.\textsc{3sg.prs}	\\}
%\glt `Petra drives a bus here [i.e., is a bus driver]. When she was little she wanted to be a painter, but her mom ultimately talked her out of it.'
\ea
\ea{\label{busdriver}
\gll Petra u nás řídí\textsuperscript{IPFV} autobus.  \\
Petra at us drive.\textsc{3sg.prs} bus \\}
\glt `Petra drives a bus here [i.e., is a bus driver].'
\ex{\label{drive.va}
\gll Petra u nás řídívá autobus, když je nedostatek řidičů.  \\
Petra at us drive.\textsc{gen.3sg.pst} bus when be.\textsc{3sg.pst} shortage driver.\textsc{masc.gen.pl}	\\}
\glt `Petra tends to drive a bus here [i.e., drives a bus now and then, with some regularity], when there's a shortage of bus drivers.'
\z
\z
Unless the context indicates otherwise,
from (\ref{busdriver}) we may naturally infer that Petra is employed as a bus driver, being a bus driver is her profession.
As such it is her characterizing property that holds of her stably over a large chunk of her life-time, i.e., at all its subintervals and moments of time without interruptions, including times when she happens to be off work, e.g., on vacation or sick.  In
contrast, (\ref{drive.va}) is naturally understood as meaning that
Petra drives a bus from time to time, often, occasionally,
just in case the appropriate occasion arises, and the like.
Hence,
(\ref{busdriver}) does not entail (\ref{drive.va}),
but rather from (\ref{busdriver}) we can infer the negation of (\ref{drive.va}).
And vice versa, (\ref{drive.va}) does not entail (\ref{busdriver}).
This can be also confirmed
by the observation that in (\ref{cancel.test}), B's cancellation with `yes' is odd or at least dispreferred:
\ea\label{cancel.test}
\begin{xlist}
\exi{A:}[]{\gll Řídívá Petra u Vás autobus?    \\
drive.\textsc{gen.3sg.prs} Petra at you bus \\
\glt `Does Petra \{occasionally / from time to time / often\dots\} drive a bus in your city?'}
\exi{B:}[]{\gll \minsp{\{} Ne / \minsp{\#} Ano\}, Petra u nás \minsp{(} dokonce) řídí\textsuperscript{IPFV} autobus.\\
{} no {} {} yes  Petra at us {} even drive.\textsc{3sg.prs} bus	\\
\glt `No, (well) Petra (in fact) is a bus driver here/in our city.' \label{cancel.va.test}\\
{`\#Yes, Petra (in fact) is a bus driver here/in our city.'}}
\end{xlist}
\z


\noindent Rank orders, similarly to scale orders, give
rise to quantity implicatures,
which are inferences triggered by an utterance and the alternative utterance(s) a speaker could have made instead. The rank order (\ref{rank.order}) may trigger
one of the two types of inference, depending on context: certainty or ignorance regarding the
speaker's commitment to exceptions.
The choice of the non-generic alternative, i.e., the strong element in the rank order in (\ref{rank.order}),
may in an appropriate context convey a certainty inference (by pragmatic strengthening) that there are \textit{no exceptions} to the generalization
in all the known conditions, in all the speaker's doxastic alternatives. This constitutes the strongest, upper bound of the rank order.

The use of the generic morpheme \textit{-va-}, which corresponds
to the weak alternative in the rank order (\ref{rank.order}), signals that there are exceptions or counterexamples
to the generalization that are of the kind that
the speaker thinks cannot be safely ignored,
given the speaker's (implicit) understanding
of the correct application of the generalization,
weighing some salient reasons for or against it.
The use of\linebreak \textit{-va-} triggers one of the two inferences, depending on context.
The speaker wishes to convey the inference that they know
that there are exceptions or counterexamples (certainty inference), \textit{or} the speaker
does not have enough information
to be in the position to exclude them (ignorance inference).

Let us examine these two inferences in more detail.
By using the generic \textit{-va-}verb (the weak alternative) when the verb without the generic morpheme
(the strong alternative) could have also been used (salva veritate), the speaker
may convey the ignorance inference that they are not in the position to use the strong alternative, because it allows
pragmatic strengthening to an exceptionless
generalization, unless it is explicitly indicated otherwise,
and the hearer might be misled.
This case has been illustrated with the minimal pair in \REF{classes.toplevel} above, repeated below.

\ea\label{classesX-toplevel}
\ea {\label{classesX}
\gll Výuka začíná\textsuperscript{IPFV} od 8.00 a končí v poledne.\\
teaching start.3\textsc{sg.prs} from 8:00 and ends in noon\\}
\glt `Classes start at 8 a.m. and end at noon.'
\ex {\label{classes.genX}
\gll  Výuka začínává  od 8.00 a končí v poledne. \\
teaching starts.\textsc{gen.3sg.prs} from 8:00 and ends in noon \\}
\glt `Classes tend to start at 8 a.m. and end at noon.' [as a rule, usually]
\z
\z

\noindent Let us suppose there is a subset of \REF{classesX}'s speaker's doxastic alternatives where the school opens at 8 a.m. and a subset where it does not; in other words, the speaker of \REF{classesX} does not know which of these alternatives holds, and considers both as live possibilities
(ignorance inference). Uttering (\ref{classes.genX})
with the generic morpheme, rather than (\ref{classesX}) without it, readily allows for plausible deniability in the face of refutation by information of which the speaker of \REF{classesX} is ignorant,
and specifically should it turn out that the school does not always open at 8 a.m.
In contrast, by uttering (\ref{classesX}) with
the non-generic verb,
rather than (\ref{classes.genX}) with the generic one,  entails the negation of the latter, and,
provided the context
does not indicate otherwise, naturally leads to a pragmatic strengthening of (\ref{classesX}) to an exceptionless generalization. I.e.,
(\ref{classesX}) is compatible with the certainty inference that
in all the speaker's doxastic alternatives associated with \REF{classesX} the school opens at 8 a.m.

The exceptive inference of the generic verb marked with \textit{-va-} (see Semantic Property 2 in (\ref{ex:semprop2})) can also be exploited to
convey the certainty inference that there are counterexamples to the generalization that would have made
the corresponding sentence without \textit{-va-} false, all else being equal.
The best examples are generalizations with counterexamples
that count as positive counterinstances
(in the sense of \citealt{leslie2008}), such as hardbacks
that falsify \textit{Books are paperbacks}.  These are inadmissible
non-confirming cases which \textit{cannot} be (easily) ignored when judging
the truth of the generalization.
Above, this was illustrated by \REF{books.paperback-toplevel} and the relevant contrast is repeated below. By using the generic morpheme in (\ref{books.paperback.true.vaX}), as opposed to the alternative without \textit{-va-} in \REF{books.paperback.falseCzechX}, the speaker (publicly) commits to the existence of
counterexamples to the generalization, which, in turn, prohibits pragmatic strengthening to an exceptionless generalization,
and guarantees the truth of (\ref{books.paperback.true.vaX}).
% by the contrast between the true generic sentence (\ref{books.paperback.true.va}) formally marked with the generic morpheme \textit{-va-} and the false generic sentence (\ref{books.paperback.falseCzech}) without \textit{-va-}.

\ea\label{books.paperbackX.toplevel}
\ea {\label{books.paperback.falseCzechX}	\gll Knihy	jsou\textsuperscript{IPFV} 	bro\v{z}ovan\'{e}.\\
books	be.3\textsc{pl.prs}	paperback \\}  \jambox*{\textsc{false}}
\glt `Books are paperback.'
\ex {\label{books.paperback.true.vaX}
\gll  Knihy	b\'{y}vaj\'{i}	bro\v{z}ovan\'{e}.	\\
books	be.\textsc{gen}.3\textsc{pl.prs}	paperback \\} \jambox*{\textsc{true}}
\glt `Books tend to be paperback.'
\z
\z


%also because it entail the negation of (\ref{books.paperback.true.va}), which is  formally marked with the generic morpheme, and it allows for pragmatic strengthening to a false exceptionless statement.

\noindent In representing such speaker-oriented certainty and ignorance inferences,
we may rely on the covert doxastic operator $\Box_{\cnst{sp}(c)}$
(inspired by \citealt{Meyer2013}),
which is understood as universally quantifying over the speaker’s
doxastic alternatives \citep{Deganoetal2025}:

\ea \label{doxastic}
A covert doxastic operator: `$\Box_{\cnst{sp}(c)}(P)$' where, \\
`$\cnst{sp}(c)$': a doxastic source, the speaker in context $c$\\
 `$\Box_{\cnst{sp}(c)}$': a covert doxastic operator\\
 `$\Box_{\cnst{sp}(c)}(P)$': in all the speaker's doxastic alternatives, the speaker knows that $P$\\
\z

\noindent
Specifically for the representation of the ignorance inference, we may employ the
classic definition of ignorance of \citet{Hintikka1962-HINKAB}, characterized as `not knowing whether'. In our case, and assuming that
$P$ corresponds to the strongest value in the rank order of the form
$\langle \langle \text{Strong} \ | \  \text{Weak} \rangle \rangle$
and compatible with ``all exceptions categorically excluded'',
we have: (i)~the speaker does not know whether $P$ is true, (ii) the speaker does not
know whether $\neg P$ is true,
and (iii) the speaker also
considers both $P$ and $\neg P$
as live possibilities compatible with the speaker's knowledge.\footnote{Traditionally, and based on Grice (e.g., \citealt{Grice1975,Grice1989}, i.a.), ignorance inferences are derived
by means of maxims of conversation (quantity and quality),
based on general principles of conversation and cooperative behavior of
interlocutors. Ignorance inferences are taken to consist of uncertainty and possibility
inferences, where the possibility inference is derived from uncertainty.
Recently, however, \citet{Deganoetal2025} have argued
for a possibility inference without the uncertainty inference. Due
the limits on this chapter, I cannot here address the debates concerning
ignorance inferences in more detail.}


With these ingredients in place, adding the two speaker-oriented
inferences to
our original rank order (\ref{rank.order}), its final formulation is given below.



\ea \label{rank.order.inferences}
\textsc{Rank order} (final version):  $\langle\langle [\cnst{Gen}]P$ | $[\cnst{va}]P\,\rangle\rangle$ \\
where, \\
$[\cnst{Gen}]P$: compatible with no exceptions  \\
$[\cnst{va}]P$: incompatible with no exceptions \\\medskip
\begin{xlist}
\exi{(i)}{Asserting that $[\cnst{Gen}]P$ holds, given that the speaker could have instead asserted that $[\cnst{va}]P$ holds, may trigger \\\smallskip
\textsc{certainty inference}: `$\Box_{\cnst{sp}(c)}(P)$'\\
unless the context indicates otherwise.\medskip}
\exi{(ii)}{Asserting that $[\cnst{va}]P$ holds, given that the speaker could have instead asserted that $[\cnst{gen}]P$ holds, triggers
one of the following inferences depending on context:
\begin{xlist}
\exi{(a)}{\textsc{certainty inference}: $\Box_{\cnst{sp}(c)} \neg (P)$}
\exi{(b)}{\textsc{ignorance inference}:
$\neg\Box_{\cnst{sp}(c)}(P)\wedge \neg\Box_{\cnst{sp}(c)} \neg(P)  \leftrightarrow  \Diamond_{\cnst{sp}(c)}(P)\wedge\Diamond_{\cnst{sp}(c)}\neg(P)$\smallskip\\
\hfill(following \citealt{Hintikka1962-HINKAB})\label{ex:hintign}}
\end{xlist}}
\end{xlist}
\z

\noindent In effect, generic sentences formally marked with the
generic morpheme \textit{-va-} are endowed with two modal (intensional) components: one that is tied to their generic, law-like, or gnomic force, which it shares with the null generic operator $\cnst{Gen}$, and the other which involves reasoning with exceptions and counterexamples,  which
may be characterized in terms of the speaker-oriented doxastic (and epistemic) modality, in interaction with general pragmatic principles of interpretation.\footnote{The speaker-oriented modality that concerns the speaker's commitment to the strength of the generalization, which is  based on the speaker's information state about exceptions or counterexamples to it can be seen as related to observations in the generic literature that
generic sentences track not only (relatively)
objective facts about the world, but also a variety of subjective factors,
including the speaker's information state, beliefs about the facts of the world, and a variety of attitudes that
the speaker may take towards them.
As has been observed above,
some generics track what we
consider to be particularly salient,
striking, dangerous or impressive facts (e.g., \textit{Sharks attack swimmers}), which, for instance, \citet{leslie2008} integrates into her cognitive approach to generic sentences, or they track
psychologically salient partitions of the kind (\citealt{cohen2004}), or reflect judgments based on stereotypes (e.g., \textit{Asian people are good at mathematics}).}
%  The W(eaker) element  in the rank order (\ref{rank.order})/(\ref{rank.order.inferences}) reflects the Semantic Property 2 in (\ref{ex:semprop2}), which means that it corresponds to generic statements that have actual exceptions or possible ones, while the  S(tronger) element is compatible with exceptions. The difference between \textit{-va-}marked verbs (``\ldots{}W \ldots'')  and the corresponding simpler unmarked verbs (``\ldots{}S \ldots'')with which they are substitutable  salva veritate in a given context, when the generalization in question is amenable to exceptions, lies in that the generalization corresponding to (``\ldots{} W \ldots'') can never be pragmatically strengthened to an exceptionlessgeneric statement, while that corresponding to  (``\ldots{} S \ldots'') can.\footnote{For an overview of the role of different types of alternatives and mechanisms in pragmatic strengthening see \citet{bade2021model}.}
%Facts about the world comprise not only `brute physical facts' (in the sense of \citealt{searle1995}) like those captured by universal laws of nature, as in (\ref{earthva}), but also `social facts' (\citealt{searle1995}) that are made possible by constitutive rules, as in (\ref{waltz}), and regulatory rules, as in (\ref{board.va}).


% \subsection{Supporting Evidence: The $\cnst{Va}$ operator with overt restrictors}

\subsection{Supporting Evidence: The \textsc{va} operator with overt restrictors}

The properties of the generic marker \textit{-va-} are relevant to discussions of generic sentences in the context of ceteris paribus (\textsc{cp}) generalizations, given that, as I propose,
its use is driven by reasoning
with exceptions, counterexamples, or various restricting
conditions which cannot be ignored or tolerated.
Philosophers draw on generic sentences in order to
advance our understanding of the nature and status of \textsc{cp} generalizations in special sciences (see e.g. \citealt{nickel2016} and references therein). A line is drawn between exceptions to the \textsc{cp} generalization that are compatible with its truth and counterexamples which are not.

The distribution of \textit{-va-} over different types of generic sentences taps into this distinction, as has been illustrated by a number of examples above. %Generally, verbs with the generic marker \textit{-va-} are not only formally but also distributionally marked in opposition to the corresponding non-generic verbs without this marker.
To what has already been observed, we may add that in contrast to non-generic verbs, there are contexts in which specifically generic verbs marked with \textit{-va-} may be odd or dispreferred with no overt restrictor, but
perfectly acceptable if they are licensed by an appropriate overt restrictor.
This can be shown in generic sentences with kind reference that describe what are taken to be characterizing properties of a kind to which there are tolerable exceptions compatible with their truth, but which would refute the corresponding universal generalization.
Consider the following contrast:

\eanoraggedright \label{exceptions.dog.sound}
\textsc{exceptions}: non-barking dogs (negative counterinstances, see \citealt{leslie2008})  \\
\textit{Context:} `What sound do dogs make?'
\ea[]{\label{psi.IPF.repeat}
\gll Psi štěkají\textsuperscript{IPFV}. \\
dogs bark.3\textsc{pl.prs} \\ \jambox*{[ =(\ref{psi.IPF})]}
\glt  `Dogs bark.'}
%Approximately: Generally, for situations \textit{C(s)} that are appropriate for dogs to bark, dogs bark in \textit{s}.
\ex[(\#)]{\label{psiczechodd}
\gll Psi štěkávají.\\
dogs bark.\textsc{gen.3\textsc{pl.prs}} \\
\glt `Dogs tend to bark [when{\dots} / if{\dots} / other things being equal / unless prevented].'}
% I.e., \ldots when/if\ldots / \ldots other things being equal / \ldots unless prevented \ldots}
%Approximately: Generally, for situations \textit{C(s)} that are appropriate for dogs to bark, dogs bark in \textit{s},  \textit{other things being equal/other things being right}.
\z
\z

\noindent In this context, the preferred and the most natural answer is (\ref{psi.IPF.repeat}) with the verb without \textit{-va-}.
%and which corresponds to the English \textit{Dogs bark}.
(\ref{psi.IPF.repeat}) and its closest English correspondent in this context \textit{Dogs bark} instantiates a prototypical generic sentence (see \sectref{CharGenericity}): namely, they have no overt restrictor or quantificational element, their bare plural subject term denotes a kind whose characterizing property
is denoted by a verb form that is the least marked form in their respective aspect and tense system, i.e., they exhibit the minimal marking tendency of genericity with respect to tense and aspect \citep{dahl1995}.\footnote{In languages with tense and/or aspect systems, prototypical generic sentences
exhibit the ``minimal marking tendency'' of genericity
with respect to tense and aspect, as \citet{dahl1995} puts it: namely,
their main verb is either devoid of any tense and aspect marking
(e.g., \textit{bark} in \textit{Dogs bark})
or corresponds to the least marked form in the verb system.
In Slavic languages it is the primary (im)perfective verb; it is
a verb form with no derivational morphemes, overt aspect marking,
but possibly only marked with tense inflectional morphology. The Czech sentence (\ref{psi.IPF.repeat}) on its generic interpretation
is headed by the imperfective simplex \textit{štěkají} `bark.\textsc{3pl.prs}', which is marked for tense, number and person, but has no marker of imperfective aspect.}

Both (\ref{psi.IPF.repeat}) and \textit{Dogs bark} are
expected and felicitous answers to the question
at hand, because it calls for the basic information about one particular characterizing property of dogs that can be \textit{truly} predicated of dogs as a whole kind, based on our shared kind-related expectations about what it means to be a (realistically) possible dog, by default, typically, or generally, in any possible world; if something is a dog, it will likely bark. Negative counterinstances (in the sense of \citealt{leslie2008}), i.e. dogs failing to bark
(e.g., breeds like greyhounds and individual dogs that do not bark for a variety of reasons like an injury), and all kinds of disturbing or facilitating factors which may be potentially infinite, are largely irrelevant in this context.

Contrary to this, by using \textit{-va-} in (\ref{psiczechodd}) the speaker shifts focus to exceptions, counterexamples, (possible)
restricting conditions on the application of the generalization, implying that (i) members of the kind \textsc{dog} are not homogeneous with respect to the characterizing property of barking, or (ii) do not bark in all the appropriate situations in which not exceptional barking dogs might be expected to bark.
This straightforwardly follows from Semantic Property 2 in (\ref{ex:semprop2}).
By using the marked alternative (\ref{psiczechodd}) with \textit{-va-}, when (\ref{psi.IPF.repeat}) without it is preferred as an answer to `What sound do dogs make?',
the speaker might be (mis)understood as conveying the ignorance inference that they are uncertain whether barking is a \textit{characterizing} property of dogs. But this would be odd, because it would be inconsistent with our general knowledge that barking is a  \textit{characterizing} property of dogs, and the speaker is expected to know this, as well as that
there exceptional dogs that fail to bark (negative counterinstances), which we tend to ignore.

Or, and
most likely, the speaker, by using the generic \textit{-va-}, may  convey the certainty inference that they know that there are other exceptions (apart from the tolerable negative counterinstances which we set aside) that are noteworthy, relevant, striking or possibly not known to the hearer in a given context, which cannot be ignored. And so the speaker intends to draw attention to such exceptional individuals, perhaps to show off his knowledge of dogs. But if so, it is odd for the speaker not to explicitly specify what these exceptions are and why they are worthy of being highlighted. This is also the reason, I submit, why (\ref{psiczechodd}) without an overt restrictor, and out of the blue, is odd or dispreferred.

Once an overt appropriate restrictor
is added to the unrestricted (\ref{psiczechodd}), we get a restricted generic sentence in which the generic marker \textit{-va-} is fully felicitous.
With an overt restrictor, the generic verb marked with
\textit{-va-} and its unmarked counterpart without \textit{-va-} are  substitutable salva veritate:
\ea
\ea {\label{psi.greyhound}
\gll Psi \minsp{\{} štěkají\textsuperscript{IPFV} / štěkávají\}, s výjimkou některých druhů psů jako chrti.\\
dogs  {} bark.3\textsc{pl.prs} {} bark.\textsc{gen.3pl.prs}  with exception some breeds dogs like greyhounds\\}
\glt `Dogs tend to bark, except for some dog breeds like greyhounds.'
\ex {\label{psi.notknow}
\gll Psi \minsp{\{} štěkají\textsuperscript{IPFV} / štěkávají\}  na ty, které neznají.  \\
dogs {} bark.3\textsc{pl.prs} {} bark.\textsc{gen.3\textsc{pl.prs}} at those whom \textsc{neg}.know\\}
\glt `Dogs bark at those whom they don't know.'
\z
\z
\noindent In (\ref{psi.greyhound}) the overt restrictor
specifies a subkind of dogs, greyhounds, that are explicitly
exempt from the generalization. In (\ref{psi.notknow}), the restrictor constrains a subset of situations in
which stages of prototypical, unexceptional dogs generally bark.

These two types of overt qualifications
are to be expected, if we assume that the
generic morpheme \textit{-va-} functions as a kind of
generic quantifier, and generic sentences with
kind reference involve a ``double generalization''
(in terms of \citealt{carlson2011}). For instance, \textit{Dogs bark} is true by virtue of (i) \textit{individual dogs} that realize the kind \textsc{dog}  having the disposition of barking (unless they belong to a non-barking kind like greyhounds), and also by virtue of (ii) \textit{particular situations} of barking by a stage of an individual dog.

On the quantificational theory of generic sentences (see \sectref{CharGenericity}), the relevant quantifier in their logical form is either
the covert $\cnst{Gen}$ operator or an overt generic Q-adverb
(e.g., \textit{usually, typically}).  And also, as I argue,
the generic $\cnst{va}$ operator, introduced by the generic marker \textit{-va-}.
All three -- $\cnst{Gen}$, a Q-Adverb, and $\cnst{va}$ -- have either a covert restrictor or an overt restrictor.\footnote{When there is no overt restrictor, as in (\ref{psi.IPF.repeat}) \textit{Dogs bark}, the context is standardly assumed to restrict the quantification to contextually appropriate situations. Roughly, (\ref{psi.IPF.repeat}) may then be paraphrased as ‘Generally, for situations $C(s)$ that are appropriate for dogs to bark, dogs bark in $s$’.}
Overt restrictors of generic quantifiers constitute a large and a heterogeneous class which
includes time/place adverbials (e.g., \textit{in/every winter}), conditional \textit{if/when}-clauses, restrictive modifiers, speaker-oriented epistemics concerning the degree of commitment to the generalization (e.g., \textit{in my opinion}), exceptive constructions (\textit{except/but/besides/apart from X}), and a variety of ceteris paribus clauses that specify (causally) disturbing or facilitating factors.

As we have also seen, the generic marker \textit{-va-} has perfectly acceptable and felicitous uses in sentences with no overt restrictor.
If so, the speaker uses it to either
ensure the truth of a generic sentence
in the face of counterexamples that would falsify the corresponding sentence without it, all else being equal (see \sectref{mustbeused}), or
to ensure plausible
deniability by preempting pragmatic strengthening to  universal/exceptionless generalization (see Section \ref{maybeused}).
%Some of the most compelling examples for the first case are generics in which the presence of \textit{-va-} ensures their truth in the face of counterexamples that falsify the corresponding generics without \textit{-va-}, and which were here identified as counterexamples that count as positive counterinstances to the generalization. They were discussed in \sectref{mustbeused}, in particular in connection with the contrast in \REF{books.paperback-toplevel}, repeated in \REF{books.paperbackX.toplevel}.
%The second case in  which \textit{-va-} does not require licensing by means of an overt restrictor concerns its use to ensure plausible deniability. This arises when the corresponding stronger alternative without \textit{-va-} would also be possible (all else being equal), but the speaker cannot commit to it because it is compatible with the corresponding universal or exceptionless generalization, and hence permits pragmatic strengthening to a generalization with no exceptions, unless the context indicates otherwise.  Examples illustrating this case are in (\ref{classes.toplevel}), repeated in \REF{classesX-toplevel}.Here, \textit{-va-} functions as a kind of hedging device to safeguard the generalization against refutation by states of affairs that the speaker knows they are ignorant of.
A case in point illustrating this use of \textit{-va-} is (\ref{classes.genX}). Here, \textit{-va-} may be used to convey one of the two inferences depending on context.  \textit{Either} the speaker conveys the certainty inference that they know that there are counterexamples to the generalization that would refute the corresponding exceptionless generalization, \textit{or} the ignorance inference that they lack enough information for making a stronger claim corresponding to an exceptionless  generalization, e.g., they have incomplete beliefs about various possible disturbing factors that may refute that stronger claim.

When the generic marker \textit{-va-} has perfectly acceptable and felicitous uses in sentences with no overt restrictor, it might
be thought of as assuming a function
that is comparable to that of a \textsc{cp}-condition, besides introducing a generic quantifier. This property clearly sets \textit{-va-} apart  from
both the covert $\cnst{Gen}$ operator, and also from overt Q-adverbs.
The term \textit{ceteris paribus} covers different kinds of phenomena, and has different interpretations. Here, \textit{-va-} may be thought of as roughly amounting to an exclusive and indefinite
(or indeterminate) \textsc{cp}-condition, e.g., ``other things being equal'', ``unless prevented'', ``in the absence of disturbing factors'', and the like. The exclusion-clause, as \citet{Cartwright1983} suggests,
may always be reformulated as a clause which requires that certain truth conditions for the exclusive \textsc{cp} law hold, namely those which exclude disturbing factors: ``the literal translation is `other things being equal’, yet it would be more apt to read `ceteris paribus’ as `other things being right’”
\citep[45]{Cartwright1983}. Generally, the number of possible disturbing or enabling factors
is potentially infinite, and therefore, the restricting condition in application of the generalization cannot always be strictly and fully specified by means of an overt restrictor.







\section{Conclusion and further implications} \label{independence.tenseaspect}

The Czech marker \textit{-va-}, which, as all agree, is separate from
the imperfectivizing suffix, is, as I argue,
best treated as a marker of genericity (\textit{pace} e.g., \citealt{dahl1995}). It is not a
marker of iterativity, frequentativity, multiplicativity, in short a mere plurality of situations,
or habituality, and as such a marker that is subsumed under tense and/or imperfective aspect, contrary to what is commonly assumed.
One of the prima facie reasons why it is wrong to classify and label it as an iterative marker is the simple fact that it is incompatible with iterative adverbials like `three times', used to count particular episodic situations that are not a part of a larger pattern.

It has been shown
that this marker has properties that prohibit its subsumption under tense and that do not neatly fit under imperfectivity, either in Slavic languages
or as it is generally understood cross-linguistically (see e.g., \citealt{Comrie1976}).
Neither can its meaning be reduced to `most' or `usually' (contrary to \citet{dahl1995}, i.a.), because it imposes no requirement on a specific quantity of situations or instances counting as evidence for the truth of the generalization.
By the same token, its meaning cannot be reduced to any other single quantifier or expression of quantity
no matter how modalized or probabilistic it might be.

Among the novel key insights of this study is that
the marker \textit{-va-} introduces a generic quantifier into the logical representation of sentences, which, similarly to the null generic operator \cnst{gen},
has a modal (intensional) import that is tied to its predictive, law-like, force, and hence requires access to other possible worlds than just the actual one. Second, it exhibits variable binding properties that are akin to \cnst{gen}.
Third, generic sentences formally marked \textit{-va-} differ from generic sentences without this marker
in so far as their truth conditions and felicity depend on the speaker's reasoning with exceptions. Specifically, they reflect the speaker's commitment
to the strength of the generalization calibrated in terms of exceptions that the speaker believes it may have or has, possibly also including counterexamples that cannot be ignored.
This additional modal component  may be characterized in terms of epistemic and doxastic modality, and its integration into the meaning of sentences also exploits general pragmatic principles of interpretation.

An exciting domain of study for future research is
the exploration of the cognate generic markers \textit{-va-}, their
common citation form across Slavic languages (see e.g., \citealt{Genis2021}, i.a.), in other Slavic languages, namely, in Slovak, where it is also
productive, in Polish, where it is less productive, but
common enough, and also in other Slavic languages where its occurrences
are treated as a part of a lexicalized combination with a verb base,
as in Russian, for instance. Hypothesizing that the properties
that are here identified for the Czech marker \textit{-va-} may plausibly
also characterize its cognates in other Slavic languages would
be at least a concrete point of departure for such (comparative) studies.

Turning to cross-linguistic consequences beyond Czech, and Slavic languages,
the properties of the Czech marker \textit{-va-} preclude its assimilation
to markers
that merely encode a plurality of situations (at a given world/time), i.e., markers that are variously labeled as iterative, frequentative, pluractional, habitual, and the like, in traditional descriptive grammars or contemporary linguistic studies.
The empirical and theoretical arguments presented here in support of the claim that the marker
\textit{-va-} in Czech, and also in other Slavic languages where it is attested in comparable examples, is a generic marker might serve as a blueprint for their exploration.
As observed above, \citet{dahl1995} after all takes
the Czech marker \textit{-va-} as a paradigm example of a whole class of markers
in a number of typologically distinct languages that he labels ``habitual tense-aspect'' markers,
and which, according to him, are not generic markers (the list of languages, is given \citet [421,~fn.~8]{dahl1995} and above).
To the extent that his claims concerning the Czech marker \textit{-va-} are invalid, as I showed, this raises at least some doubts whether
other members of Dahl's class that the Czech marker \textit{-va-} is supposed to represent are all indeed ``habitual tense-aspect'' markers, as Dahl claims.
In short, we may pose the following question for future research:

\eanoraggedright
What is the status of so-called ``habitual tense-aspect'' markers (in terms of \citealt{dahl1995}, i.a.)
with respect to the categories of the tense and aspect systems?
Do they have (at least some of the) formal and semantic properties that have been here identified
for the Czech marker \textit{-va-} which make it a plausible candidate for a generic marker?
\z

\noindent Markers that are treated as habitual markers in Tlingit \citep{cable2022}, Modern Hebrew \citep{bonehdoron2010}, African-American Vernacular English \citep{green2000aspectual}, just to name a few, are reported to introduce the realization inference just like
the Czech generic morpheme \textit{-va-} does. It is an open empirical question whether they also carry the uncancellable inference that the generalization requires (at least the possibility of)
exceptions or counterexamples, given that these two meaning components
might be seen as related in the case of the Czech  generic marker \textit{-va-}. What is also striking is that
such previous studies focus on sentences that describe habits in the strict sense, i.e., regularities of action of ordinary animate individuals, mostly human agents, such as \textit{Fred smokes}. It is unclear whether such so-called ``habitual'' markers can also be used to mark generic sentences that express regularities that
have nothing to do with habits proper, just
like the Czech marker \textit{-va-} does.

Perhaps future research will show that at least some
markers that are referred to as habitual, and also as iterative, frequentative, multiplicative, and the like, are best
classified as markers of genericity. If so, it
would further bolster the claim that markers of genericity are not as rare as is thought
(\textit{pace} \citealt{dahl1995}),
and invalidate the widespread traditional belief
that they do not exist in natural languages (\textit{pace}
\citealt{leslie2008}, \citealt{LeslieGelman2012}, \citealt{Johnston2012},  \citealt{Rhodes2018}, \citealt{cohen2022}, \citealt{Schiller2023}, i.a.)
This would put the semantic episodic/generic distinction on a firmer empirical ground as a semantic distinction with direct grammatical reflexes, in addition to its indirect reflexes in the tense and aspect systems of natural languages, and in a variety of semantic and
pragmatic phenomena.
Consequently, this would add further arguments for the independence  of genericity from
the categories of the TMA system in natural languages, and specifically from tense and imperfective aspect.
So far such arguments have mainly been of semantic/pragmatic nature, or based on reasoning strategies involved in generalizations which are developed in philosophy, cognitive psychology and computer science (see e.g., \citealt{carlson1977unified,carlson1982generic,krifka1995genericity,pelletier1997generics}).

Assuming that the Czech marker \textit{-va-},
and plausibly also the Slavic marker \textit{-va-}, is a generic marker that introduces a generic operator into the logical representation whose
properties overlap with, but are not identical to, those of the null generic operator \cnst{gen}, brings us to the fundamental question in the domain of genericity posed at the outset (see also Carlson 1977, 1995): Can we provide a unified semantic analysis for all generic sentences?
In Czech, generic sentences fall into two main types: one formally marked with \textit{-va-}, and the other formally unmarked. In the latter case,
and similarly as in English, the generic (and habitual) reading is a function of the combined meanings of the subject and predicate interacting with context,
and it can be represented by means of a tripartite logical structure with the \cnst{gen} operator. Recall that in all Slavic languages,
generic sentences that are not formally marked
with the generic \textit{-va-} are headed by
primary or secondary imperfectives (marked for imperfectivity with the imperfectivizing suffix),
and in West Slavic languages also commonly by perfective verbs.

These two main formal means for the expression of characterizing genericity, one formally marked with the morpheme \textit{-va-} for genericity and the other formally unmarked in this respect, are not neatly aligned with the inductive and the rules-and-regulations model of genericity \citep{carlson1995truth}, respectively.
The generic morpheme \textit{-va-} is naturally used in generic sentences that provide the best examples in support of the inductive model of genericity, namely generics that express descriptive (weak) generalizations, including habits in the strict sense (i.e., regularities of action ascribed to ordinary animate individuals, as in \REF{va-tenses}). However, it is important to emphasize that the Czech generic \textit{-va-} is also used in generic sentences that best fit the rules-and-regulations model of genericity, including generics with a normative reading.
% previous version:
% The use of the generic morpheme \textit{-va-}, as we have seen, while it is naturally used in generic sentences that provide the best examples in support of the inductive model of genericity, namely generics that express descriptive (weak) generalizations, including habits in the strict sense (i.e., regularities of action ascribed to ordinary animate individuals, as in \REF{va-tenses}), it is also used in generic sentences that best fit the rules-and-regulations model of genericity, including generics with a normative reading.
In short, its use cuts across the inductive versus rules-and-regulations distinction, and also across the related descriptive versus normative distinction.

% At the same time, these two different linguistic
% means of expressing genericity, one with the generic \textit{-va-} and the other without it, correspond to two different subtypes of generics, each requiring different semantic and pragmatic commitments, and separate models for their interpretation, albeit not neatly aligned with the divide
% between the inductive and rules-and-regulations models.
% Similar conclusions were independently reached
% for other languages,
% based on the analysis of different linguistic forms
% that they have available or obligatory for
% the encoding of generic statements (e.g., \citealt{greenberg2007exceptions,pelletier-equal2009,boneh2008habituality,cable2022}, i.a.).
% Tentatively, we may suggest that this
% raises doubts about whether
% a single model for the interpretation of all generic sentences is viable,
% or at best it leaves the question about the possibility of characterizing such a single model still open.

\newpage
Tentatively, it may be concluded that this leaves open the question about the feasibility of a unified semantic analysis for all generic sentences in Czech, and perhaps also in other languages that have a  dedicated marker of genericity that is comparable to that in Czech.

\section*{Abbreviations}

\begin{tabularx}{.5\textwidth}{@{}lQ}
\textsc{1}&first person\\
\textsc{2}&second person\\
\textsc{3}&third person\\
\textsc{acc}&accusative\\
\textsc{aux}&auxiliary\\
\textsc{cp}&ceteris paribus\\
\textsc{dat}&dative\\
\textsc{gen}&generic\\
\textsc{inf}&infinitive\\
\textsc{ipfv}&imperfective\\
\textsc{nci}&negative concord item\\
\end{tabularx}
\begin{tabularx}{.5\textwidth}{lQ@{}}
\textsc{neg}&negation\\
\textsc{nom}&nominative\\
\textsc{pfv}&perfective\\
\textsc{pl}&plural\\
\textsc{pref}&prefix\\
\textsc{prs}&present tense\\
\textsc{pst}&past tense\\
\textsc{refl}&reflexive\\
\textsc{sg}&singular\\
\textsc{va}&morpheme \textit{-va-}\\
&\\ % this dummy row achieves correct vertical alignment of both tables
\end{tabularx}

\section*{Acknowledgements}
I would like to thank Greg Carlson and Manfred Krifka for their suggestions, and inspiration, as well as the audiences at the \textit{Formal Diachronic Semantics 6} conference (University of
Cologne, Germany), \textit{The 6th International Conference on Grammar and Text} (School of Social Sciences and Humanities of NOVA University Lisbon, Portugal), \textit{The 20th Szklarska Poręba Workshop on the Roots of
Pragmasemantics} (Szklarska Poręba, Poland), \textit{The workshop on Formal Semantics and
Its Interfaces} (Zhejiang University, Hangzhou, China),
and \textit{The 19th Conference on Formal Approaches to Slavic
Linguistics} (University of Maryland, USA) for all the valuable criticism and suggestions. All opinions, omissions, and errors remain my own.

%\section*{Contributions}
%John Doe contributed to conceptualization, methodology, and validation.
%Jane Doe contributed to writing of the original draft, review, and editing.

\section*{Literary sources of data}

\noindent Hana Bělohradská: \textit{Poslední večeř}e. Praha: Český spisovatel. 1992.\\
\noindent Karel Čapek: \textit{Obyčejný život.} Olomouc: Otakar II. 1934.\\
\noindent Karel Čapek: \textit{Hovory s T.G. Masarykem.} Praha: Fr. Borový: Čin, 1936.\\
%https://catalog.library.tamu.edu/Record/in00001486166
Karel Kryl: \textit{7 básniček na zrcadlo}. Mnichov: Soukromá publikace, 1974.

\sloppy
\printbibliography[heading=subbibliography,notkeyword=this]

\end{document}
